\documentclass{article}
\usepackage[utf8]{inputenc}
\usepackage{amsmath, amssymb, graphicx, xcolor, fancyvrb}

% --- UBAH JUDUL, NAMA, DAN NIM ANDA DI SINI ---
\title{TugasAplikom6Adi25030630001}
\author{Adi Kurniawan \\ NIM: 25030630001}
\date{\today}

% --- PENERJEMAH OTOMATIS KODE EMT (JANGAN DIUBAH) ---
\newenvironment{eulernotebook}{\par}{\par}
\newcommand{\eulerheading}[1]{\section{#1}}
\newcommand{\eulersubheading}[1]{\subsection{#1}}
\newenvironment{eulercomment}{\par\medskip}{\par\medskip}
\newenvironment{eulerformula}{}{}
\DefineVerbatimEnvironment{eulerprompt}{Verbatim}{formatcom=\color{red}}
\DefineVerbatimEnvironment{euleroutput}{Verbatim}{formatcom=\color{blue}}
\DefineVerbatimEnvironment{eulerttcomment}{Verbatim}{formatcom=\color{teal}}
\newcommand{\eulerimg}[2]{\begin{center}\includegraphics[width=0.8\textwidth]{#2}\end{center}}

\begin{document}
\maketitle

% ====================================================================
% TEMPELKAN KODE DARI \begin{eulernotebook} SAMPAI \end{eulernotebook}
% DI BAWAH BARIS INI:
\begin{eulernotebook}
\eulerheading{Kalkulus dengan EMT}
\begin{eulercomment}
Materi Kalkulus mencakup di antaranya:

- Fungsi (fungsi aljabar, trigonometri, eksponensial, logaritma, komposisi fungsi)\\
- Limit Fungsi,\\
- Turunan Fungsi,\\
- Integral Tak Tentu,\\
- Integral Tentu dan Aplikasinya,\\
- Barisan dan Deret (kekonvergenan barisan dan deret).

EMT (bersama Maxima) dapat digunakan untuk melakukan semua perhitungan di dalam kalkulus, baik secara numerik maupun analitik
(eksak).

\end{eulercomment}
\eulersubheading{Mendefinisikan Fungsi}
\begin{eulercomment}
Terdapat beberapa cara mendefinisikan fungsi pada EMT, yakni:

- Menggunakan format nama\_fungsi := rumus fungsi (untuk fungsi numerik),\\
- Menggunakan format nama\_fungsi \&= rumus fungsi (untuk fungsi simbolik, namun dapat dihitung secara numerik),\\
- Menggunakan format nama\_fungsi \&\&= rumus fungsi (untuk fungsi simbolik murni, tidak dapat dihitung langsung),\\
- Fungsi sebagai program EMT.

Setiap format harus diawali dengan perintah function (bukan sebagai ekspresi).

Berikut adalah adalah beberapa contoh cara mendefinisikan fungsi:

\end{eulercomment}
\begin{eulerformula}
\[
f(x)=2x^2+e^{\sin(x)}.
\]
\end{eulerformula}
\begin{eulerprompt}
>function f(x) := 2*x^2+exp(sin(x)) // fungsi numerik
>f(0), f(1), f(pi)
\end{eulerprompt}
\begin{euleroutput}
  1
  4.31977682472
  20.7392088022
\end{euleroutput}
\begin{eulerprompt}
>f(a) // tidak dapat dihitung nilainya
\end{eulerprompt}
\begin{euleroutput}
  Variable or function a not found.
  Error in:
  f(a) // tidak dapat dihitung nilainya ...
     ^
\end{euleroutput}
\begin{eulercomment}
Silakan Anda plot kurva fungsi di atas!
\end{eulercomment}
\begin{eulerprompt}
> aspect(1.5); plot2d("f(x)", -pi, pi):
\end{eulerprompt}
\eulerimg{17}{Image Use/TugasAplikom6Adi25030630001-001.png}
\begin{eulercomment}
Berikutnya kita definisikan fungsi:

\end{eulercomment}
\begin{eulerformula}
\[
g(x)=\frac{\sqrt{x^2-3x}}{x+1}.
\]
\end{eulerformula}
\begin{eulerprompt}
>function g(x) := sqrt(x^2-3*x)/(x+1)
>g(3)
\end{eulerprompt}
\begin{euleroutput}
  0
\end{euleroutput}
\begin{eulerprompt}
>g(0)
\end{eulerprompt}
\begin{euleroutput}
  0
\end{euleroutput}
\begin{eulerprompt}
>g(1) // kompleks, tidak dapat dihitung oleh fungsi numerik
\end{eulerprompt}
\begin{euleroutput}
  Floating point error!
  Error in sqrt
  Try "trace errors" to inspect local variables after errors.
  g:
      useglobal; return sqrt(x^2-3*x)/(x+1) 
  Error in:
  g(1) // kompleks, tidak dapat dihitung oleh fungsi numerik ...
      ^
\end{euleroutput}
\begin{eulercomment}
Silakan Anda plot kurva fungsi di atas!
\end{eulercomment}
\begin{eulerprompt}
>aspect(1.5); plot2d("g(x)", 3, 10):
\end{eulerprompt}
\eulerimg{17}{Image Use/TugasAplikom6Adi25030630001-003.png}
\begin{eulerprompt}
>f(g(5)) // komposisi fungsi
\end{eulerprompt}
\begin{euleroutput}
  2.20920171961
\end{euleroutput}
\begin{eulerprompt}
>g(f(5))
\end{eulerprompt}
\begin{euleroutput}
  0.950898070639
\end{euleroutput}
\begin{eulerprompt}
>function h(x) := f(g(x)) // definisi komposisi fungsi 
>h(5) // sama dengan f(g(5))
\end{eulerprompt}
\begin{euleroutput}
  2.20920171961
\end{euleroutput}
\begin{eulercomment}
Silakan Anda plot kurva fungsi komposisi fungsi f dan g:

\end{eulercomment}
\begin{eulerformula}
\[
h(x)=f(g(x))
\]
\end{eulerformula}
\begin{eulercomment}
dan

\end{eulercomment}
\begin{eulerformula}
\[
u(x)=g(f(x))
\]
\end{eulerformula}
\begin{eulercomment}
bersama-sama kurva fungsi f dan g dalam satu bidang koordinat.
\end{eulercomment}
\begin{eulerprompt}
>function f(x) := 2*x^2+exp(sin(x))
>function g(x) := sqrt(x^2-3*x)/(x+1)
>function h(x) := f(g(x))
>function u(x) := g(f(x))
>f(g(5))
\end{eulerprompt}
\begin{euleroutput}
  2.20920171961
\end{euleroutput}
\begin{eulerprompt}
>g(f(5))
\end{eulerprompt}
\begin{euleroutput}
  0.950898070639
\end{euleroutput}
\begin{eulerprompt}
>h(5)
\end{eulerprompt}
\begin{euleroutput}
  2.20920171961
\end{euleroutput}
\begin{eulerprompt}
>aspect(1.5); plot2d(["f(x)","g(x)","h(x)","u(x)"], 3, 10, color=[blue,red,green,orange]):
\end{eulerprompt}
\eulerimg{17}{Image Use/TugasAplikom6Adi25030630001-004.png}
\begin{eulerprompt}
> 
>f(0:10) // nilai-nilai f(0), f(1), f(2), ..., f(10)
\end{eulerprompt}
\begin{euleroutput}
  [1,  4.31978,  10.4826,  19.1516,  32.4692,  50.3833,  72.7562,
  99.929,  130.69,  163.51,  200.58]
\end{euleroutput}
\begin{eulerprompt}
>fmap(0:10) // sama dengan f(0:10), berlaku untuk semua fungsi
\end{eulerprompt}
\begin{euleroutput}
  [1,  4.31978,  10.4826,  19.1516,  32.4692,  50.3833,  72.7562,
  99.929,  130.69,  163.51,  200.58]
\end{euleroutput}
\begin{eulerprompt}
>gmap(200:210)
\end{eulerprompt}
\begin{euleroutput}
  [0.987534,  0.987596,  0.987657,  0.987718,  0.987778,  0.987837,
  0.987896,  0.987954,  0.988012,  0.988069,  0.988126]
\end{euleroutput}
\begin{eulercomment}
Misalkan kita akan mendefinisikan fungsi

\end{eulercomment}
\begin{eulerformula}
\[
f(x) = \begin{cases} x^3 & x>0 \\ x^2 & x\le 0. \end{cases}
\]
\end{eulerformula}
\begin{eulercomment}
Fungsi tersebut tidak dapat didefinisikan sebagai fungsi numerik
secara "inline" menggunakan format :=, melainkan didefinisikan sebagai
program. Perhatikan, kata "map" digunakan agar fungsi dapat menerima
vektor sebagai input, dan hasilnya berupa vektor. Jika tanpa kata
"map" fungsinya hanya dapat menerima input satu nilai.
\end{eulercomment}
\begin{eulerprompt}
>function map f(x) ...
\end{eulerprompt}
\begin{eulerudf}
    if x>0 then return x^3
    else return x^2
    endif;
  endfunction
\end{eulerudf}
\begin{eulerprompt}
>f(1)
\end{eulerprompt}
\begin{euleroutput}
  1
\end{euleroutput}
\begin{eulerprompt}
>f(-2)
\end{eulerprompt}
\begin{euleroutput}
  4
\end{euleroutput}
\begin{eulerprompt}
>f(-5:5)
\end{eulerprompt}
\begin{euleroutput}
  [25,  16,  9,  4,  1,  0,  1,  8,  27,  64,  125]
\end{euleroutput}
\begin{eulerprompt}
>aspect(1.5); plot2d("f(x)",-5,5):
\end{eulerprompt}
\eulerimg{17}{Image Use/TugasAplikom6Adi25030630001-006.png}
\begin{eulerprompt}
>function f(x) &= 2*E^x // fungsi simbolik
\end{eulerprompt}
\begin{euleroutput}
  
                                      x
                                   2 E
  
\end{euleroutput}
\begin{eulerprompt}
>$f(a) // nilai fungsi secara simbolik
\end{eulerprompt}
\begin{eulerformula}
\[
2\,e^{a}
\]
\end{eulerformula}
\begin{eulerprompt}
>f(E) // nilai fungsi berupa bilangan desimal
\end{eulerprompt}
\begin{euleroutput}
  30.308524483
\end{euleroutput}
\begin{eulerprompt}
>$f(E), $float(%)
\end{eulerprompt}
\begin{eulerformula}
\[
2\,e^{e}
\]
\end{eulerformula}
\begin{eulerformula}
\[
30.30852448295852
\]
\end{eulerformula}
\begin{eulerprompt}
>function g(x) &= 3*x+1
\end{eulerprompt}
\begin{euleroutput}
  
                                 3 x + 1
  
\end{euleroutput}
\begin{eulerprompt}
>function h(x) &= f(g(x)) // komposisi fungsi
\end{eulerprompt}
\begin{euleroutput}
  
                                   3 x + 1
                                2 E
  
\end{euleroutput}
\begin{eulerprompt}
>plot2d("h(x)",-1,1):
\end{eulerprompt}
\eulerimg{17}{Image Use/TugasAplikom6Adi25030630001-010.png}
\eulerheading{Latihan}
\begin{eulercomment}
Bukalah buku Kalkulus. Cari dan pilih beberapa (paling sedikit 5
fungsi berbeda tipe/bentuk/jenis) fungsi dari buku tersebut, kemudian
definisikan fungsi-fungsi tersebut dan komposisinya di EMT pada
baris-baris perintah berikut (jika perlu tambahkan lagi). Untuk setiap
fungsi, hitung beberapa nilainya, baik untuk satu nilai maupun vektor.
Gambar grafik fungsi-fungsi tersebut dan komposisi-komposisi 2 fungsi.

Juga, carilah fungsi beberapa (dua) variabel. Lakukan hal sama seperti
di atas.


\end{eulercomment}
\begin{eulerttcomment}
 MENDEFINISIKAN 5 FUNGSI BERBEDA TIPE 
\end{eulerttcomment}
\begin{eulerprompt}
>function f1(x) := x^3 - 4*x            // 1. Fungsi Polinomial
>function f2(x) := (x^2 - 1)/(x^2 + 1)  // 2. Fungsi Rasional
>function f3(x) := sin(2*x) + cos(x)    // 3. Fungsi Trigonometri
>function f4(x) := exp(-x^2)            // 4. Fungsi Eksponensial
>function f5(x) := sqrt(x^2 + 9)        // 5. Fungsi Akar/Irrasional                
\end{eulerprompt}
\begin{eulercomment}
HITUNG BEBERAPA NILAI (SATU NILAI DAN VEKTOR) 
\end{eulercomment}
\begin{eulerprompt}
>f1(2), f2(0), f3(pi/4), f4(1), f5(4)
\end{eulerprompt}
\begin{euleroutput}
  0
  -1
  1.70710678119
  0.367879441171
  5
\end{euleroutput}
\begin{eulerprompt}
>f1(-2:2)
\end{eulerprompt}
\begin{euleroutput}
  [0,  3,  0,  -3,  0]
\end{euleroutput}
\begin{eulerprompt}
>f2(0:5)
\end{eulerprompt}
\begin{euleroutput}
  [-1,  0,  0.6,  0.8,  0.882353,  0.923077]
\end{euleroutput}
\begin{eulerprompt}
>f3(linspace(0,pi,5))
\end{eulerprompt}
\begin{euleroutput}
  [1,  1.76007,  0.896802,  -0.896802,  -1.76007,  -1]
\end{euleroutput}
\begin{eulerprompt}
>f4(-2:2)
\end{eulerprompt}
\begin{euleroutput}
  [0.0183156,  0.367879,  1,  0.367879,  0.0183156]
\end{euleroutput}
\begin{eulerprompt}
>f5(0:4)
\end{eulerprompt}
\begin{euleroutput}
  [3,  3.16228,  3.60555,  4.24264,  5]
\end{euleroutput}
\begin{eulercomment}
KOMPOSISI DUA FUNGSI

Mendefinisikan komposisi h1(x) = f1(f3(x)) dan h2(x) = f4(f5(x))
\end{eulercomment}
\begin{eulerprompt}
>function h1(x) := f1(f3(x))
>function h2(x) := f4(f5(x))
\end{eulerprompt}
\begin{eulercomment}
Menghitung nilai komposisi pada satu titik dan vektor
\end{eulercomment}
\begin{eulerprompt}
>h1(pi/6), h2(2)
\end{eulerprompt}
\begin{euleroutput}
  -1.73205080757
  2.26032940698e-06
\end{euleroutput}
\begin{eulerprompt}
>h1(linspace(-pi,pi,5))
\end{eulerprompt}
\begin{euleroutput}
  [3,  -1.09341,  0.565292,  -1.58783,  2.86595,  3]
\end{euleroutput}
\begin{eulercomment}
MENGGAMBAR GRAFIK FUNGSI DAN KOMPOSISI FUNGSI
\end{eulercomment}
\begin{eulerprompt}
>aspect(1.5); plot2d(["f1(x)","f3(x)","h1(x)"], -pi, pi, color=[blue,red,green]):
\end{eulerprompt}
\eulerimg{17}{Image Use/TugasAplikom6Adi25030630001-011.png}
\begin{eulerprompt}
>aspect(1.5); plot2d(["f4(x)","f5(x)","h2(x)"], -3, 3, color=[blue,red,green]):
\end{eulerprompt}
\eulerimg{17}{Image Use/TugasAplikom6Adi25030630001-012.png}
\begin{eulercomment}
FUNGSI DUA VARIABEL 
\end{eulercomment}
\begin{eulerprompt}
>function g2v(x,y) := x^2 + y^2 - 4
>g2v(2,3)
\end{eulerprompt}
\begin{euleroutput}
  9
\end{euleroutput}
\begin{eulerprompt}
>g2v([1,2],[3,4])
\end{eulerprompt}
\begin{euleroutput}
  [6,  16]
\end{euleroutput}
\begin{eulerprompt}
>plot3d("g2v(x,y)", r=3):
\end{eulerprompt}
\eulerimg{17}{Image Use/TugasAplikom6Adi25030630001-013.png}
\begin{eulerprompt}
>     
\end{eulerprompt}
\begin{eulercomment}
\begin{eulercomment}
\eulerheading{Menghitung Limit}
\begin{eulercomment}
Perhitungan limit pada EMT dapat dilakukan dengan menggunakan fungsi Maxima, yakni "limit".
Fungsi "limit" dapat digunakan untuk menghitung limit fungsi dalam bentuk ekspresi maupun fungsi
yang sudah didefinisikan sebelumnya. Nilai limit dapat dihitung pada sebarang nilai atau pada tak
hingga (-inf, minf, dan inf). Limit kiri dan limit kanan juga dapat dihitung, dengan cara memberi
opsi "plus" atau "minus". Hasil limit dapat berupa nilai, "und" (tak definisi), "ind" (tak tentu
namun terbatas), "infinity" (kompleks tak hingga).

Perhatikan beberapa contoh berikut. Perhatikan cara menampilkan perhitungan secara lengkap, tidak
hanya menampilkan hasilnya saja.
\end{eulercomment}
\begin{eulerprompt}
>$showev('limit(sqrt(x^2-3*x)/(x+1),x,inf))
\end{eulerprompt}
\begin{eulerformula}
\[
\lim_{x\rightarrow \infty }{\frac{\sqrt{x^2-3\,x}}{x+1}}=1
\]
\end{eulerformula}
\begin{eulerprompt}
>$limit((x^3-13*x^2+51*x-63)/(x^3-4*x^2-3*x+18),x,3)
\end{eulerprompt}
\begin{eulerformula}
\[
-\frac{4}{5}
\]
\end{eulerformula}
\begin{eulercomment}
maxima: 'limit((x\textasciicircum{}3-13*x\textasciicircum{}2+51*x-63)/(x\textasciicircum{}3-4*x\textasciicircum{}2-3*x+18),x,3)=limit((x\textasciicircum{}3-13*x\textasciicircum{}2+51*x-63)/(x\textasciicircum{}3-4*x\textasciicircum{}2-3*x+18),x,3)

Fungsi tersebut diskontinu di titik x=3. Berikut adalah grafik fungsinya.
\end{eulercomment}
\begin{eulerprompt}
>aspect(1.5); plot2d("(x^3-13*x^2+51*x-63)/(x^3-4*x^2-3*x+18)",0,4); plot2d(3,-4/5,>points,style="ow",>add):
\end{eulerprompt}
\eulerimg{17}{Image Use/TugasAplikom6Adi25030630001-016.png}
\begin{eulerprompt}
>$limit(2*x*sin(x)/(1-cos(x)),x,0)
\end{eulerprompt}
\begin{eulerformula}
\[
4
\]
\end{eulerformula}
\begin{eulercomment}
maxima: 'limit(2*x*sin(x)/(1-cos(x)),x,0)=limit(2*x*sin(x)/(1-cos(x)),x,0)

Fungsi tersebut diskontinu di titik x=0. Berikut adalah grafik fungsinya.
\end{eulercomment}
\begin{eulerprompt}
>plot2d("2*x*sin(x)/(1-cos(x))",-pi,pi); plot2d(0,4,>points,style="ow",>add):
\end{eulerprompt}
\eulerimg{17}{Image Use/TugasAplikom6Adi25030630001-018.png}
\begin{eulerprompt}
>$limit(cot(7*h)/cot(5*h),h,0)
\end{eulerprompt}
\begin{eulerformula}
\[
\frac{5}{7}
\]
\end{eulerformula}
\begin{eulercomment}
maxima: showev('limit(cot(7*h)/cot(5*h),h,0))

Fungsi tersebut juga diskontinu (karena tidak terdefinisi) di x=0. Berikut adalah grafiknya.
\end{eulercomment}
\begin{eulerprompt}
>plot2d("cot(7*x)/cot(5*x)",-0.001,0.001); plot2d(0,5/7,>points,style="ow",>add):
\end{eulerprompt}
\eulerimg{17}{Image Use/TugasAplikom6Adi25030630001-020.png}
\begin{eulerprompt}
>$showev('limit(((x/8)^(1/3)-1)/(x-8),x,8))
\end{eulerprompt}
\begin{eulerformula}
\[
\lim_{x\rightarrow 8}{\frac{\frac{x^{\frac{1}{3}}}{2}-1}{x-8}}=  \frac{1}{24}
\]
\end{eulerformula}
\begin{eulercomment}
Tunjukkan limit tersebut dengan grafik, seperti contoh-contoh sebelumnya.
\end{eulercomment}
\begin{eulerprompt}
>plot2d("((x/8)^(1/3)-1)/(x-8)", 0, 16); plot2d(8, 1/24, >points, style="ow", >add):
\end{eulerprompt}
\eulerimg{17}{Image Use/TugasAplikom6Adi25030630001-022.png}
\begin{eulerprompt}
>$showev('limit(1/(2*x-1),x,0))
\end{eulerprompt}
\begin{eulerformula}
\[
\lim_{x\rightarrow 0}{\frac{1}{2\,x-1}}=-1
\]
\end{eulerformula}
\begin{eulercomment}
Tunjukkan limit tersebut dengan grafik, seperti contoh-contoh sebelumnya.
\end{eulercomment}
\begin{eulerprompt}
>plot2d("1/(2*x-1)", -1, 1); plot2d(0, -1, >points, style="ow", >add):
\end{eulerprompt}
\eulerimg{17}{Image Use/TugasAplikom6Adi25030630001-024.png}
\begin{eulerprompt}
>$showev('limit((x^2-3*x-10)/(x-5),x,5))
\end{eulerprompt}
\begin{eulerformula}
\[
\lim_{x\rightarrow 5}{\frac{x^2-3\,x-10}{x-5}}=7
\]
\end{eulerformula}
\begin{eulercomment}
Tunjukkan limit tersebut dengan grafik, seperti contoh-contoh sebelumnya.
\end{eulercomment}
\begin{eulerprompt}
>plot2d("(x^2-3*x-10)/(x-5)", 0, 10); plot2d(5, 7, >points, style="ow", >add):
\end{eulerprompt}
\eulerimg{17}{Image Use/TugasAplikom6Adi25030630001-026.png}
\begin{eulerprompt}
>$showev('limit(sqrt(x^2+x)-x,x,inf))
\end{eulerprompt}
\begin{eulerformula}
\[
\lim_{x\rightarrow \infty }{\sqrt{x^2+x}-x}=\frac{1}{2}
\]
\end{eulerformula}
\begin{eulercomment}
Tunjukkan limit tersebut dengan grafik, seperti contoh-contoh sebelumnya.
\end{eulercomment}
\begin{eulerprompt}
>plot2d("sqrt(x^2+x)-x", 0, 50):
\end{eulerprompt}
\eulerimg{17}{Image Use/TugasAplikom6Adi25030630001-028.png}
\begin{eulerprompt}
>$showev('limit(abs(x-1)/(x-1),x,1,minus))
\end{eulerprompt}
\begin{eulerformula}
\[
\lim_{x\uparrow 1}{\frac{\left| x-1\right| }{x-1}}=-1
\]
\end{eulerformula}
\begin{eulercomment}
Hitung limit di atas untuk x menuju 1 dari kanan.\\
Tunjukkan limit tersebut dengan grafik, seperti contoh-contoh sebelumnya.
\end{eulercomment}
\begin{eulerprompt}
>$showev('limit(abs(x-1)/(x-1),x,1,plus))
\end{eulerprompt}
\begin{eulerformula}
\[
\lim_{x\downarrow 1}{\frac{\left| x-1\right| }{x-1}}=1
\]
\end{eulerformula}
\begin{eulerprompt}
>plot2d("abs(x-1)/(x-1)", -1, 3); plot2d(1, -1, >points, style="ow", >add); plot2d(1, 1, >points, style="ow", >add):
\end{eulerprompt}
\eulerimg{17}{Image Use/TugasAplikom6Adi25030630001-031.png}
\begin{eulerprompt}
>$showev('limit(sin(x)/x,x,0))
\end{eulerprompt}
\begin{eulerformula}
\[
\lim_{x\rightarrow 0}{\frac{\sin x}{x}}=1
\]
\end{eulerformula}
\begin{eulerprompt}
>plot2d("sin(x)/x",-pi,pi); plot2d(0,1,>points,style="ow",>add):
\end{eulerprompt}
\eulerimg{17}{Image Use/TugasAplikom6Adi25030630001-033.png}
\begin{eulerprompt}
>$showev('limit(sin(x^3)/x,x,0))
\end{eulerprompt}
\begin{eulerformula}
\[
\lim_{x\rightarrow 0}{\frac{\sin x^3}{x}}=0
\]
\end{eulerformula}
\begin{eulercomment}
Tunjukkan limit tersebut dengan grafik, seperti contoh-contoh sebelumnya.
\end{eulercomment}
\begin{eulerprompt}
>plot2d("sin(x^3)/x", -pi, pi); plot2d(0, 0, >points, style="ow", >add):
\end{eulerprompt}
\eulerimg{17}{Image Use/TugasAplikom6Adi25030630001-035.png}
\begin{eulerprompt}
>$showev('limit(log(x), x, minf))
\end{eulerprompt}
\begin{eulerformula}
\[
\lim_{x\rightarrow  -\infty }{\log x}={\it infinity}
\]
\end{eulerformula}
\begin{eulerprompt}
>$showev('limit((-2)^x,x, inf))
\end{eulerprompt}
\begin{eulerformula}
\[
\lim_{x\rightarrow \infty }{\left(-2\right)^{x}}={\it infinity}
\]
\end{eulerformula}
\begin{eulerprompt}
>$showev('limit(t-sqrt(2-t),t,2,minus))
\end{eulerprompt}
\begin{eulerformula}
\[
\lim_{t\uparrow 2}{t-\sqrt{2-t}}=2
\]
\end{eulerformula}
\begin{eulerprompt}
>$showev('limit(t-sqrt(2-t),t,2,plus))
\end{eulerprompt}
\begin{eulerformula}
\[
\lim_{t\downarrow 2}{t-\sqrt{2-t}}=2
\]
\end{eulerformula}
\begin{eulerprompt}
>$showev('limit(t-sqrt(2-t),t,5,plus)) // Perhatikan hasilnya
\end{eulerprompt}
\begin{eulerformula}
\[
\lim_{t\downarrow 5}{t-\sqrt{2-t}}=5-\sqrt{3}\,i
\]
\end{eulerformula}
\begin{eulerprompt}
>plot2d("x-sqrt(2-x)",0,2):
\end{eulerprompt}
\eulerimg{17}{Image Use/TugasAplikom6Adi25030630001-041.png}
\begin{eulerprompt}
>$showev('limit((x^2-9)/(2*x^2-5*x-3),x,3))
\end{eulerprompt}
\begin{eulerformula}
\[
\lim_{x\rightarrow 3}{\frac{x^2-9}{2\,x^2-5\,x-3}}=\frac{6}{7}
\]
\end{eulerformula}
\begin{eulercomment}
Tunjukkan limit tersebut dengan grafik, seperti contoh-contoh sebelumnya.
\end{eulercomment}
\begin{eulerprompt}
>plot2d("(x^2-9)/(2*x^2-5*x-3)", 0, 6); plot2d(3, 6/7, >points, style="ow", >add):
\end{eulerprompt}
\eulerimg{17}{Image Use/TugasAplikom6Adi25030630001-043.png}
\begin{eulerprompt}
>$showev('limit((1-cos(x))/x,x,0))
\end{eulerprompt}
\begin{eulerformula}
\[
\lim_{x\rightarrow 0}{\frac{1-\cos x}{x}}=0
\]
\end{eulerformula}
\begin{eulercomment}
Tunjukkan limit tersebut dengan grafik, seperti contoh-contoh sebelumnya.
\end{eulercomment}
\begin{eulerprompt}
> plot2d("(1-cos(x))/x", -pi, pi); plot2d(0, 0, >points, style="ow", >add):
\end{eulerprompt}
\eulerimg{17}{Image Use/TugasAplikom6Adi25030630001-045.png}
\begin{eulerprompt}
> $showev('limit((x^2+abs(x))/(x^2-abs(x)),x,0))
\end{eulerprompt}
\begin{eulerformula}
\[
\lim_{x\rightarrow 0}{\frac{\left| x\right| +x^2}{x^2-\left| x  \right| }}=-1
\]
\end{eulerformula}
\begin{eulercomment}
Tunjukkan limit tersebut dengan grafik, seperti contoh-contoh sebelumnya.
\end{eulercomment}
\begin{eulerprompt}
>plot2d("(x^2+abs(x))/(x^2-abs(x))", -1, 1); plot2d(0, -1, >points, style="ow", >add):
\end{eulerprompt}
\eulerimg{17}{Image Use/TugasAplikom6Adi25030630001-047.png}
\begin{eulerprompt}
>$showev('limit((1+1/x)^x,x,inf))
\end{eulerprompt}
\begin{eulerformula}
\[
\lim_{x\rightarrow \infty }{\left(\frac{1}{x}+1\right)^{x}}=e
\]
\end{eulerformula}
\begin{eulerprompt}
>plot2d("(1+1/x)^x",0,1000):
\end{eulerprompt}
\eulerimg{17}{Image Use/TugasAplikom6Adi25030630001-049.png}
\begin{eulerprompt}
>$showev('limit((1+k/x)^x,x,inf))
\end{eulerprompt}
\begin{eulerformula}
\[
\lim_{x\rightarrow \infty }{\left(\frac{k}{x}+1\right)^{x}}=e^{k}
\]
\end{eulerformula}
\begin{eulerprompt}
>$showev('limit((1+x)^(1/x),x,0))
\end{eulerprompt}
\begin{eulerformula}
\[
\lim_{x\rightarrow 0}{\left(x+1\right)^{\frac{1}{x}}}=e
\]
\end{eulerformula}
\begin{eulercomment}
Tunjukkan limit tersebut dengan grafik, seperti contoh-contoh sebelumnya.
\end{eulercomment}
\begin{eulerprompt}
>plot2d("(1+x)^(1/x)", -0.9, 2); plot2d(0, E, >points, style="ow", >add):
\end{eulerprompt}
\eulerimg{17}{Image Use/TugasAplikom6Adi25030630001-052.png}
\begin{eulerprompt}
>$showev('limit((x/(x+k))^x,x,inf))
\end{eulerprompt}
\begin{eulerformula}
\[
\lim_{x\rightarrow \infty }{\left(\frac{x}{x+k}\right)^{x}}=e^ {- k   }
\]
\end{eulerformula}
\begin{eulerprompt}
>$showev('limit((E^x-E^2)/(x-2),x,2))
\end{eulerprompt}
\begin{eulerformula}
\[
\lim_{x\rightarrow 2}{\frac{e^{x}-e^2}{x-2}}=e^2
\]
\end{eulerformula}
\begin{eulercomment}
Tunjukkan limit tersebut dengan grafik, seperti contoh-contoh sebelumnya.
\end{eulercomment}
\begin{eulerprompt}
>plot2d("(exp(x)-exp(2))/(x-2)", 0, 4); plot2d(2, exp(2), >points, style="ow", >add):
\end{eulerprompt}
\eulerimg{17}{Image Use/TugasAplikom6Adi25030630001-055.png}
\begin{eulerprompt}
>$showev('limit(sin(1/x),x,0))
\end{eulerprompt}
\begin{eulerformula}
\[
\lim_{x\rightarrow 0}{\sin \left(\frac{1}{x}\right)}={\it ind}
\]
\end{eulerformula}
\begin{eulerprompt}
>$showev('limit(sin(1/x),x,inf))
\end{eulerprompt}
\begin{eulerformula}
\[
\lim_{x\rightarrow \infty }{\sin \left(\frac{1}{x}\right)}=0
\]
\end{eulerformula}
\begin{eulerprompt}
>plot2d("sin(1/x)",-0.001,0.001):
\end{eulerprompt}
\eulerimg{17}{Image Use/TugasAplikom6Adi25030630001-058.png}
\eulerheading{Latihan}
\begin{eulercomment}
Bukalah buku Kalkulus. Cari dan pilih beberapa (paling sedikit 5 fungsi berbeda
tipe/bentuk/jenis) fungsi dari buku tersebut, kemudian definisikan di EMT pada
baris-baris perintah berikut (jika perlu tambahkan lagi). Untuk setiap fungsi, hitung
nilai limit fungsi tersebut di beberapa nilai dan di tak hingga. Gambar grafik fungsi
tersebut untuk mengkonfirmasi nilai-nilai limit tersebut.
\end{eulercomment}
\begin{eulercomment}
MENDEFINISIKAN FUNGSI, HUNG LIMIT, DAN PLOT GRAFIK

1. FUNGSI TRIGONOMETRI
\end{eulercomment}
\begin{eulerprompt}
>$showev('limit(sin(5*x)/(3*x), x, 0))
\end{eulerprompt}
\begin{eulerformula}
\[
\frac{\lim_{x\rightarrow 0}{\frac{\sin \left(5\,x\right)}{x}}}{3}=  \frac{5}{3}
\]
\end{eulerformula}
\begin{eulerprompt}
>$showev('limit(sin(5*x)/(3*x), x, inf))
\end{eulerprompt}
\begin{eulerformula}
\[
\frac{\lim_{x\rightarrow \infty }{\frac{\sin \left(5\,x\right)}{x}}  }{3}=0
\]
\end{eulerformula}
\begin{eulerprompt}
>aspect(1.5); plot2d("sin(5*x)/(3*x)", -pi, pi); plot2d(0, 5/3, >points, style="ow", >add):
\end{eulerprompt}
\eulerimg{17}{Image Use/TugasAplikom6Adi25030630001-061.png}
\begin{eulercomment}
2. FUNGSI RASIONAL
\end{eulercomment}
\begin{eulerprompt}
> $showev('limit((x^2-4)/(x-2), x, 2))
\end{eulerprompt}
\begin{eulerformula}
\[
\lim_{x\rightarrow 2}{\frac{x^2-4}{x-2}}=4
\]
\end{eulerformula}
\begin{eulerprompt}
>$showev('limit((x^2-4)/(x-2), x, inf))
\end{eulerprompt}
\begin{eulerformula}
\[
\lim_{x\rightarrow \infty }{\frac{x^2-4}{x-2}}=\infty 
\]
\end{eulerformula}
\begin{eulerprompt}
>aspect(1.5); plot2d("(x^2-4)/(x-2)", -2, 6); plot2d(2, 4, >points, style="ow", >add):
\end{eulerprompt}
\eulerimg{17}{Image Use/TugasAplikom6Adi25030630001-064.png}
\begin{eulercomment}
3. FUNGSI EKSPONENSIAL
\end{eulercomment}
\begin{eulerprompt}
>$showev('limit((exp(2*x)-1)/x, x, 0))
\end{eulerprompt}
\begin{eulerformula}
\[
\lim_{x\rightarrow 0}{\frac{e^{2\,x}-1}{x}}=2
\]
\end{eulerformula}
\begin{eulerprompt}
>$showev('limit((exp(2*x)-1)/x, x, inf))
\end{eulerprompt}
\begin{eulerformula}
\[
\lim_{x\rightarrow \infty }{\frac{e^{2\,x}-1}{x}}=\infty 
\]
\end{eulerformula}
\begin{eulerprompt}
>aspect(1.5); plot2d("(exp(2*x)-1)/x", -2, 2); plot2d(0, 2, >points, style="ow", >add):
\end{eulerprompt}
\eulerimg{17}{Image Use/TugasAplikom6Adi25030630001-067.png}
\begin{eulercomment}
4. FUNGSI AKAR / IRRASIONAL
\end{eulercomment}
\begin{eulerprompt}
>$showev('limit((sqrt(x+4)-2)/x, x, 0))
\end{eulerprompt}
\begin{eulerformula}
\[
\lim_{x\rightarrow 0}{\frac{\sqrt{x+4}-2}{x}}=\frac{1}{4}
\]
\end{eulerformula}
\begin{eulerprompt}
>$showev('limit((sqrt(x+4)-2)/x, x, inf))
\end{eulerprompt}
\begin{eulerformula}
\[
\lim_{x\rightarrow \infty }{\frac{\sqrt{x+4}-2}{x}}=0
\]
\end{eulerformula}
\begin{eulerprompt}
>aspect(1.5); plot2d("(sqrt(x+4)-2)/x", -3, 5); plot2d(0, 1/4, >points, style="ow", >add):
\end{eulerprompt}
\eulerimg{17}{Image Use/TugasAplikom6Adi25030630001-070.png}
\begin{eulercomment}
5. FUNGSI LOGARITMA
\end{eulercomment}
\begin{eulerprompt}
>$showev('limit(x*log(1+1/x), x, inf))
\end{eulerprompt}
\begin{eulerformula}
\[
\lim_{x\rightarrow \infty }{\log \left(\frac{1}{x}+1\right)\,x}=1
\]
\end{eulerformula}
\begin{eulerprompt}
>$showev('limit(x*log(1+1/x), x, 0, plus))
\end{eulerprompt}
\begin{eulerformula}
\[
\lim_{x\downarrow 0}{\log \left(\frac{1}{x}+1\right)\,x}=0
\]
\end{eulerformula}
\begin{eulerprompt}
>aspect(1.5); plot2d("x*log(1+1/x)", 0.01, 50):
\end{eulerprompt}
\eulerimg{17}{Image Use/TugasAplikom6Adi25030630001-073.png}
\begin{eulercomment}
\begin{eulercomment}
\eulerheading{Turunan Fungsi}
\begin{eulercomment}
Definisi turunan:

\end{eulercomment}
\begin{eulerformula}
\[
f'(x) = \lim_{h\to 0} \frac{f(x+h)-f(x)}{h}
\]
\end{eulerformula}
\begin{eulercomment}
Berikut adalah contoh-contoh menentukan turunan fungsi dengan
menggunakan definisi turunan (limit).
\end{eulercomment}
\begin{eulerprompt}
>$showev('limit(((x+h)^2-x^2)/h,h,0)) // turunan x^2
\end{eulerprompt}
\begin{eulerformula}
\[
\lim_{h\rightarrow 0}{\frac{\left(x+h\right)^2-x^2}{h}}=2\,x
\]
\end{eulerformula}
\begin{eulerprompt}
>p &= expand((x+h)^2-x^2)|simplify; $p //pembilang dijabarkan dan disederhanakan
\end{eulerprompt}
\begin{eulerformula}
\[
2\,h\,x+h^2
\]
\end{eulerformula}
\begin{eulerprompt}
>q &=ratsimp(p/h); $q // ekspresi yang akan dihitung limitnya disederhanakan
\end{eulerprompt}
\begin{eulerformula}
\[
2\,x+h
\]
\end{eulerformula}
\begin{eulerprompt}
>$limit(q,h,0) // nilai limit sebagai turunan
\end{eulerprompt}
\begin{eulerformula}
\[
2\,x
\]
\end{eulerformula}
\begin{eulerprompt}
>$showev('limit(((x+h)^n-x^n)/h,h,0)) // turunan x^n
\end{eulerprompt}
\begin{eulerformula}
\[
\lim_{h\rightarrow 0}{\frac{\left(x+h\right)^{n}-x^{n}}{h}}=n\,x^{n  -1}
\]
\end{eulerformula}
\begin{eulercomment}
Mengapa hasilnya seperti itu? Tuliskan atau tunjukkan bahwa hasil
limit tersebut benar, sehingga benar turunan fungsinya benar.  Tulis
penjelasan Anda di komentar ini.

Sebagai petunjuk, ekspansikan (x+h)\textasciicircum{}n dengan menggunakan teorema
binomial.

Penjelasan Pembuktian Turunan f(x) = x\textasciicircum{}n Menggunakan Teorema
BinomialMenggunakan definisi turunan:

Berdasarkan Teorema Binomial, ekspansi dari (x+h)\textasciicircum{}n adalah:\\
\end{eulercomment}
\begin{eulerttcomment}
 (x+h)^n = x^n + n*x^(n-1)*h + (n*(n-1)/2)*x^(n-2)*h^2 + ... + h^n
\end{eulerttcomment}
\begin{eulercomment}

Pembilang ekspresi limit menjadi:\\
\end{eulercomment}
\begin{eulerttcomment}
 (x+h)^n - x^n = n*x^(n-1)*h + (n*(n-1)/2)*x^(n-2)*h^2 + ... + h^n
               = h * [n*x^(n-1) + (n*(n-1)/2)*x^(n-2)*h + ... + h^
\end{eulerttcomment}
\begin{eulercomment}
(n-1)]

Saat dibagi dengan h dan dihitung limit h -\textgreater{} 0:\\
\end{eulercomment}
\begin{eulerttcomment}
 limit_(h->0) [(x+h)^n - x^n]/h = limit_(h->0) [n*x^(n-1) + (n*(n-
\end{eulerttcomment}
\begin{eulercomment}
1)/2)*x\textasciicircum{}(n-2)*h + ... + h\textasciicircum{}(n-1)]

Karena semua suku setelah suku pertama memiliki faktor h, maka ketika
h -\textgreater{} 0\\
\end{eulercomment}
\begin{eulerttcomment}
 nilai suku-suku tersebut menjadi 0, sehingga tersisa n*x^(n-1).
 Oleh karena itu, terbukti bahwa turunan dari f(x) = x^n adalah f'(x)
\end{eulerttcomment}
\begin{eulercomment}
= n*x\textasciicircum{}(n-1).
\end{eulercomment}
\begin{eulerprompt}
> 
>$showev('limit((sin(x+h)-sin(x))/h,h,0)) // turunan sin(x)
\end{eulerprompt}
\begin{eulerformula}
\[
\lim_{h\rightarrow 0}{\frac{\sin \left(x+h\right)-\sin x}{h}}=\cos   x
\]
\end{eulerformula}
\begin{eulercomment}
Mengapa hasilnya seperti itu? Tuliskan atau tunjukkan bahwa hasil
limit tersebut\\
benar, sehingga benar turunan fungsinya benar.  Tulis penjelasan Anda
di komentar ini.

Sebagai petunjuk, ekspansikan sin(x+h) dengan menggunakan rumus jumlah
dua sudut.

Penjelasan:\\
\end{eulercomment}
\begin{eulerttcomment}
 Berdasarkan rumus jumlah dua sudut trigonometri:
 sin(x+h) = sin(x)*cos(h) + cos(x)*sin(h)
\end{eulerttcomment}
\begin{eulercomment}

\end{eulercomment}
\begin{eulerttcomment}
 Pembilang ekspresi limit menjadi:
 sin(x+h) - sin(x) = sin(x)*cos(h) + cos(x)*sin(h) - sin(x)
                   = sin(x)*(cos(h) - 1) + cos(x)*sin(h)
\end{eulerttcomment}
\begin{eulercomment}

\end{eulercomment}
\begin{eulerttcomment}
 Masukkan ke dalam definisi limit turunan:
 limit_(h->0) [sin(x+h) - sin(x)] / h
   = limit_(h->0) [ sin(x)*(cos(h) - 1)/h + cos(x)*(sin(h)/h) ]
   = sin(x) * limit_(h->0) [(cos(h) - 1)/h] + cos(x) * limit_(h->0)
\end{eulerttcomment}
\begin{eulercomment}
[sin(h)/h]

\end{eulercomment}
\begin{eulerttcomment}
 Menggunakan nilai limit dasar trigonometri:
 1. limit_(h->0) sin(h)/h = 1
 2. limit_(h->0) (cos(h) - 1)/h = 0
\end{eulerttcomment}
\begin{eulercomment}

\end{eulercomment}
\begin{eulerttcomment}
 Maka diperoleh:
   = sin(x) * (0) + cos(x) * (1) = cos(x)
\end{eulerttcomment}
\begin{eulercomment}

\end{eulercomment}
\begin{eulerttcomment}
 Terbukti bahwa turunan dari f(x) = sin(x) adalah f'(x) = cos(x).
\end{eulerttcomment}
\begin{eulerprompt}
>$showev('limit((log(x+h)-log(x))/h,h,0)) // turunan log(x)
\end{eulerprompt}
\begin{eulerformula}
\[
\lim_{h\rightarrow 0}{\frac{\log \left(x+h\right)-\log x}{h}}=  \frac{1}{x}
\]
\end{eulerformula}
\begin{eulercomment}
Mengapa hasilnya seperti itu? Tuliskan atau tunjukkan bahwa hasil
limit tersebut\\
benar, sehingga benar turunan fungsinya benar.  Tulis penjelasan Anda
di komentar ini.

Sebagai petunjuk, gunakan sifat-sifat logaritma dan hasil limit pada
bagian sebelumnya di atas.

Penejelasan Mengapa Hasil Limit Tersebut d/dx[log(x)] = 1/x:

1. Berdasarkan definisi limit turunan:\\
\end{eulercomment}
\begin{eulerttcomment}
   lim_\{h->0\} (log(x+h) - log(x)) / h
\end{eulerttcomment}
\begin{eulercomment}

2. Menggunakan sifat pengurangan logaritma: log(a) - log(b) = log(a/b)\\
\end{eulercomment}
\begin{eulerttcomment}
   = lim_\{h->0\} (1/h) * log((x+h)/x)
   = lim_\{h->0\} (1/h) * log(1 + h/x)
\end{eulerttcomment}
\begin{eulercomment}

3. Menggunakan sifat pangkat logaritma: c * log(a) = log(a\textasciicircum{}c)\\
\end{eulercomment}
\begin{eulerttcomment}
   = lim_\{h->0\} log(1 + h/x)^(1/h)
\end{eulerttcomment}
\begin{eulercomment}

4. Pemisalan u = h/x  ==\textgreater{}  h = x*u. Saat h -\textgreater{} 0, maka u -\textgreater{} 0.\\
\end{eulercomment}
\begin{eulerttcomment}
   = lim_\{u->0\} log(1 + u)^(1/(x*u))
   = lim_\{u->0\} log([ (1 + u)^(1/u) ]^(1/x))
\end{eulerttcomment}
\begin{eulercomment}

5. Karena fungsi logaritma bersifat kontinu, limit masuk ke dalam:\\
\end{eulercomment}
\begin{eulerttcomment}
   = log( lim_\{u->0\} (1 + u)^(1/u) )^(1/x)
\end{eulerttcomment}
\begin{eulercomment}

6. Berdasarkan definisi bilangan Euler: lim\_\{u-\textgreater{}0\} (1 + u)\textasciicircum{}(1/u) = e\\
\end{eulercomment}
\begin{eulerttcomment}
   = log(e^(1/x))
   = (1/x) * log(e)
\end{eulerttcomment}
\begin{eulercomment}

7. Karena log(e) = 1 (logaritma natural):\\
\end{eulercomment}
\begin{eulerttcomment}
   = 1/x
\end{eulerttcomment}
\begin{eulercomment}

Kesimpulan: Terbukti bahwa turunan dari log(x) adalah 1/x.
\end{eulercomment}
\begin{eulerprompt}
>$showev('limit((1/(x+h)-1/x)/h,h,0)) // turunan 1/x
\end{eulerprompt}
\begin{eulerformula}
\[
\lim_{h\rightarrow 0}{\frac{\frac{1}{x+h}-\frac{1}{x}}{h}}=-\frac{1  }{x^2}
\]
\end{eulerformula}
\begin{eulerprompt}
>$showev('limit((E^(x+h)-E^x)/h,h,0)) // turunan f(x)=e^x
\end{eulerprompt}
\begin{euleroutput}
  Answering "Is x an integer?" with "integer"
  Answering "Is x an integer?" with "integer"
  Answering "Is x an integer?" with "integer"
  Answering "Is x an integer?" with "integer"
  Answering "Is x an integer?" with "integer"
  Maxima is asking
  Acceptable answers are: yes, y, Y, no, n, N, unknown, uk
  Is x an integer?
  
  Use assume!
  Error in:
  $showev('limit((E^(x+h)-E^x)/h,h,0)) // turunan f(x)=e^x ...
                                       ^
\end{euleroutput}
\begin{eulercomment}
Maxima bermasalah dengan limit:

\end{eulercomment}
\begin{eulerformula}
\[
\lim_{h\to 0}\frac{e^{x+h}-e^x}{h}.
\]
\end{eulerformula}
\begin{eulercomment}
Oleh karena itu diperlukan trik khusus agar hasilnya benar.

1. Secara analitis, bentuk limit turunan e\textasciicircum{}x adalah:\\
\end{eulercomment}
\begin{eulerttcomment}
  lim_\{h->0\} (e^(x+h) - e^x) / h
\end{eulerttcomment}
\begin{eulercomment}

2. Menggunakan sifat eksponen: e\textasciicircum{}(x+h) = e\textasciicircum{}x * e\textasciicircum{}h, maka:\\
\end{eulercomment}
\begin{eulerttcomment}
  lim_\{h->0\} (e^x * e^h - e^x) / h
  = lim_\{h->0\} e^x * (e^h - 1) / h
\end{eulerttcomment}
\begin{eulercomment}

3. Karena e\textasciicircum{}x konstan terhadap h, keluarkan e\textasciicircum{}x dari limit:\\
\end{eulercomment}
\begin{eulerttcomment}
  = e^x * lim_\{h->0\} (e^h - 1) / h
\end{eulerttcomment}
\begin{eulercomment}

4. Karena nilai lim\_\{h-\textgreater{}0\} (e\textasciicircum{}h - 1)/h = 1, maka hasil akhirnya:\\
\end{eulercomment}
\begin{eulerttcomment}
  = e^x * 1 = e^x
\end{eulerttcomment}
\begin{eulercomment}

Di EMT/Maxima, trik ini ditulis dengan memisahkan \%e\textasciicircum{}x di luar
perintah limit\\
agar Maxima tidak mengalami loop pertanyaan variabel:

\textdollar{}showev(\%e\textasciicircum{}x * 'limit((\%e\textasciicircum{}h - 1)/h, h, 0))
\end{eulercomment}
\begin{eulerprompt}
> $showev(%e^x * 'limit((%e^h - 1)/h, h, 0))
\end{eulerprompt}
\begin{eulerformula}
\[
\left(\lim_{h\rightarrow 0}{\frac{e^{h}-1}{h}}\right)\,e^{x}=e^{x}
\]
\end{eulerformula}
\begin{eulerprompt}
>$showev('limit((E^h-1)/h,h,0))
\end{eulerprompt}
\begin{eulerformula}
\[
\lim_{h\rightarrow 0}{\frac{e^{h}-1}{h}}=1
\]
\end{eulerformula}
\begin{eulerprompt}
>$showev('factor(E^(x+h)-E^x))
\end{eulerprompt}
\begin{eulerformula}
\[
{\it factor}\left(e^{x+h}-e^{x}\right)=\left(e^{h}-1\right)\,e^{x}
\]
\end{eulerformula}
\begin{eulerprompt}
>$showev('limit(factor((E^(x+h)-E^x)/h),h,0)) // turunan f(x)=e^x
\end{eulerprompt}
\begin{eulerformula}
\[
\left(\lim_{h\rightarrow 0}{\frac{e^{h}-1}{h}}\right)\,e^{x}=e^{x}
\]
\end{eulerformula}
\begin{eulerprompt}
>function f(x) &= x^x
\end{eulerprompt}
\begin{euleroutput}
  
                                     x
                                    x
  
\end{euleroutput}
\begin{eulerprompt}
>$showev('limit(f(x),x,0))
\end{eulerprompt}
\begin{eulerformula}
\[
\lim_{x\rightarrow 0}{x^{x}}=1
\]
\end{eulerformula}
\begin{eulercomment}
Silakan Anda gambar kurva

\end{eulercomment}
\begin{eulerformula}
\[
y=x^x.
\]
\end{eulerformula}
\begin{eulerprompt}
> plot2d("f(x)",-2,5,-2,2); plot2d(0,1,>points,style=".",>add):
\end{eulerprompt}
\eulerimg{17}{Image Use/TugasAplikom6Adi25030630001-090.png}
\begin{eulerprompt}
>$showev('limit((f(x+h)-f(x))/h,h,0)) // turunan f(x)=x^x
\end{eulerprompt}
\begin{eulerformula}
\[
\lim_{h\rightarrow 0}{\frac{\left(x+h\right)^{x+h}-x^{x}}{h}}=  {\it infinity}
\]
\end{eulerformula}
\begin{eulercomment}
Di sini Maxima juga bermasalah terkait limit:

\end{eulercomment}
\begin{eulerformula}
\[
\lim_{h\to 0} \frac{(x+h)^{x+h}-x^x}{h}.
\]
\end{eulerformula}
\begin{eulercomment}
Dalam hal ini diperlukan asumsi nilai x.
\end{eulercomment}
\begin{eulerprompt}
> &assume(x>0); $showev('limit((f(x+h)-f(x))/h,h,0)) // turunan f(x)=x^x
\end{eulerprompt}
\begin{eulerformula}
\[
\lim_{h\rightarrow 0}{\frac{\left(x+h\right)^{x+h}-x^{x}}{h}}=x^{x}  \,\left(\log x+1\right)
\]
\end{eulerformula}
\begin{eulercomment}
Mengapa hasilnya seperti itu? Tuliskan atau tunjukkan bahwa hasil
limit tersebut benar, sehingga benar turunan fungsinya benar. Tulis
penjelasan Anda di komentar ini.

Penjelasan Analitis Hasil Limit Turunan f(x) = x\textasciicircum{}x:

1. Definisi Fungsi dan Logaritma Natural:\\
\end{eulercomment}
\begin{eulerttcomment}
   Misalkan y = f(x) = x^x.
   Untuk menurunkan fungsi eksponensial dengan basis dan pangkat
\end{eulerttcomment}
\begin{eulercomment}
berupa variabel,\\
\end{eulercomment}
\begin{eulerttcomment}
   kita gunakan turunan logaritmik dengan mengenakan logaritma natural
\end{eulerttcomment}
\begin{eulercomment}
(ln) di kedua sisi:\\
\end{eulercomment}
\begin{eulerttcomment}
   ln(y) = ln(x^x) = x * ln(x)
\end{eulerttcomment}
\begin{eulercomment}

2. Turunan Implisit terhadap x:\\
\end{eulercomment}
\begin{eulerttcomment}
   Turunkan kedua sisi terhadap x menggunakan aturan rantai dan aturan
\end{eulerttcomment}
\begin{eulercomment}
perkalian:\\
\end{eulercomment}
\begin{eulerttcomment}
   (1/y) * (dy/dx) = (d/dx [x]) * ln(x) + x * (d/dx [ln(x)])
   (1/y) * (dy/dx) = 1 * ln(x) + x * (1/x)
   (1/y) * (dy/dx) = ln(x) + 1
\end{eulerttcomment}
\begin{eulercomment}

3. Substitusi Kembali Nilai y = x\textasciicircum{}x:\\
\end{eulercomment}
\begin{eulerttcomment}
   dy/dx = y * (ln(x) + 1)
   dy/dx = x^x * (ln(x) + 1)
\end{eulerttcomment}
\begin{eulercomment}

4. Hubungan dengan Limit Turunan:\\
\end{eulercomment}
\begin{eulerttcomment}
   Berdasarkan definisi limit turunan f'(x):
   lim_\{h->0\} ((x+h)^(x+h) - x^x) / h = f'(x)
\end{eulerttcomment}
\begin{eulercomment}

\end{eulercomment}
\begin{eulerttcomment}
   Karena f'(x) = x^x * (ln(x) + 1), maka nilai limit tersebut
\end{eulerttcomment}
\begin{eulercomment}
terbukti sama dengan:\\
\end{eulercomment}
\begin{eulerttcomment}
   x^x * (log(x) + 1)  [Di Maxima/EMT, log(x) menyatakan logaritma
\end{eulerttcomment}
\begin{eulercomment}
natural / ln(x)]
\end{eulercomment}
\begin{eulerprompt}
>&forget(x>0) // jangan lupa, lupakan asumsi untuk kembali ke semula
\end{eulerprompt}
\begin{euleroutput}
  
                                 [x > 0]
  
\end{euleroutput}
\begin{eulerprompt}
>&forget(x<0)
\end{eulerprompt}
\begin{euleroutput}
  
                                 [x < 0]
  
\end{euleroutput}
\begin{eulerprompt}
>&facts()
\end{eulerprompt}
\begin{euleroutput}
  
                                    []
  
\end{euleroutput}
\begin{eulerprompt}
>$showev('limit((asin(x+h)-asin(x))/h,h,0)) // turunan arcsin(x)
\end{eulerprompt}
\begin{eulerformula}
\[
\lim_{h\rightarrow 0}{\frac{\arcsin \left(x+h\right)-\arcsin x}{h}}=  \frac{1}{\sqrt{1-x^2}}
\]
\end{eulerformula}
\begin{eulercomment}
Mengapa hasilnya seperti itu? Tuliskan atau tunjukkan bahwa hasil
limit tersebut benar, sehingga benar turunan fungsinya benar. Tulis
penjelasan Anda di komentar ini.

Penjelasan Analitis Pembuktian Turunan f(x) = arcsin(x):

1. Pemisalan dan Fungsi Invers:\\
\end{eulercomment}
\begin{eulerttcomment}
   Misalkan y = arcsin(x), maka berdasarkan definisi fungsi invers:
   x = sin(y)  dengan domain -1 < x < 1 dan range -pi/2 < y < pi/2.
\end{eulerttcomment}
\begin{eulercomment}

2. Turunan Implisit terhadap x:\\
\end{eulercomment}
\begin{eulerttcomment}
   Turunkan kedua sisi x = sin(y) terhadap x menggunakan aturan
\end{eulerttcomment}
\begin{eulercomment}
rantai:\\
\end{eulercomment}
\begin{eulerttcomment}
   d/dx [x] = d/dx [sin(y)]
   1 = cos(y) * (dy/dx)
   dy/dx = 1 / cos(y)
\end{eulerttcomment}
\begin{eulercomment}

3. Mengubah cos(y) ke dalam Bentuk x:\\
\end{eulercomment}
\begin{eulerttcomment}
   Gunakan identitas trigonometri dasar: sin^2(y) + cos^2(y) = 1
   cos^2(y) = 1 - sin^2(y)
   cos(y) = sqrt(1 - sin^2(y))
\end{eulerttcomment}
\begin{eulercomment}

\end{eulercomment}
\begin{eulerttcomment}
   Karena x = sin(y), maka:
   cos(y) = sqrt(1 - x^2)
\end{eulerttcomment}
\begin{eulercomment}

4. Substitusi Kembali:\\
\end{eulercomment}
\begin{eulerttcomment}
   dy/dx = 1 / sqrt(1 - x^2)
\end{eulerttcomment}
\begin{eulercomment}

5. Hubungan dengan Limit Turunan:\\
\end{eulercomment}
\begin{eulerttcomment}
   Berdasarkan definisi limit turunan f'(x):
   lim_\{h->0\} (arcsin(x+h) - arcsin(x)) / h = f'(x)
\end{eulerttcomment}
\begin{eulercomment}

\end{eulercomment}
\begin{eulerttcomment}
   Terbukti bahwa nilai limit tersebut adalah 1 / sqrt(1 - x^2).
\end{eulerttcomment}
\begin{eulerprompt}
>$showev('limit((tan(x+h)-tan(x))/h,h,0)) // turunan tan(x)
\end{eulerprompt}
\begin{eulerformula}
\[
\lim_{h\rightarrow 0}{\frac{\tan \left(x+h\right)-\tan x}{h}}=  \frac{1}{\cos ^2x}
\]
\end{eulerformula}
\begin{eulercomment}
Mengapa hasilnya seperti itu? Tuliskan atau tunjukkan bahwa hasil
limit tersebut benar, sehingga benar turunan fungsinya benar. Tulis
penjelasan Anda di komentar ini.

Penjelasan Analitis Pembuktian Turunan f(x) = tan(x):

1. Definisi Identitas Trigonometri:\\
\end{eulercomment}
\begin{eulerttcomment}
   f(x) = tan(x) = sin(x) / cos(x)
\end{eulerttcomment}
\begin{eulercomment}

2. Aturan Pembagian Turunan (Quotient Rule):\\
\end{eulercomment}
\begin{eulerttcomment}
   Jika f(x) = u(x) / v(x), maka f'(x) = (u' * v - u * v') / (v^2)
   Dengan:
   - u(x) = sin(x)  => u'(x) = cos(x)
   - v(x) = cos(x)  => v'(x) = -sin(x)
\end{eulerttcomment}
\begin{eulercomment}

3. Substitusi ke Rumus Turunan:\\
\end{eulercomment}
\begin{eulerttcomment}
   f'(x) = (cos(x) * cos(x) - sin(x) * (-sin(x))) / cos^2(x)
   f'(x) = (cos^2(x) + sin^2(x)) / cos^2(x)
\end{eulerttcomment}
\begin{eulercomment}

4. Identitas Pythagoras Trigonometri:\\
\end{eulercomment}
\begin{eulerttcomment}
   Karena sin^2(x) + cos^2(x) = 1, maka:
   f'(x) = 1 / cos^2(x)  (atau sec^2(x))
\end{eulerttcomment}
\begin{eulercomment}

5. Hubungan dengan Limit Turunan:\\
\end{eulercomment}
\begin{eulerttcomment}
   Berdasarkan definisi limit turunan f'(x):
   lim_\{h->0\} (tan(x+h) - tan(x)) / h = f'(x)
\end{eulerttcomment}
\begin{eulercomment}

\end{eulercomment}
\begin{eulerttcomment}
   Terbukti bahwa nilai limit tersebut adalah 1 / cos^2(x).
\end{eulerttcomment}
\begin{eulerprompt}
>function f(x) &= sinh(x) // definisikan f(x)=sinh(x)
\end{eulerprompt}
\begin{euleroutput}
  
                                 sinh(x)
  
\end{euleroutput}
\begin{eulerprompt}
>function df(x) &= limit((f(x+h)-f(x))/h,h,0); $df(x) // df(x) = f'(x)
\end{eulerprompt}
\begin{eulerformula}
\[
\frac{e^ {- x }\,\left(e^{2\,x}+1\right)}{2}
\]
\end{eulerformula}
\begin{eulercomment}
Hasilnya adalah cosh(x), karena

\end{eulercomment}
\begin{eulerformula}
\[
\frac{e^x+e^{-x}}{2}=\cosh(x).
\]
\end{eulerformula}
\begin{eulerprompt}
>plot2d(["f(x)","df(x)"],-pi,pi,color=[blue,red]):
\end{eulerprompt}
\eulerimg{17}{Image Use/TugasAplikom6Adi25030630001-098.png}
\begin{eulerprompt}
>function f(x) &= sin(3*x^5+7)^2
\end{eulerprompt}
\begin{euleroutput}
  
                                 2    5
                              sin (3 x  + 7)
  
\end{euleroutput}
\begin{eulerprompt}
>diff(f,3), diffc(f,3)
\end{eulerprompt}
\begin{euleroutput}
  1198.32948904
  1198.72863721
\end{euleroutput}
\begin{eulercomment}
Apakah perbedaan diff dan diffc?

Perbedaan Fungsi diff dan diffc di Euler Math Toolbox (EMT):

1. Fungsi diff(f, x):\\
\end{eulercomment}
\begin{eulerttcomment}
   - Menggunakan metode beda maju (Forward Difference) untuk
\end{eulerttcomment}
\begin{eulercomment}
menghitung numerik turunan.\\
\end{eulercomment}
\begin{eulerttcomment}
   - Rumus: f'(x) ˜ (f(x + h) - f(x)) / h
   - Metode ini hanya memperhitungkan titik di sebelah kanan (x + h).
\end{eulerttcomment}
\begin{eulercomment}

2. Fungsi diffc(f, x):\\
\end{eulercomment}
\begin{eulerttcomment}
   - Menggunakan metode beda pusat (Central Difference) untuk
\end{eulerttcomment}
\begin{eulercomment}
menghitung numerik turunan.\\
\end{eulercomment}
\begin{eulerttcomment}
   - Rumus: f'(x) ˜ (f(x + h) - f(x - h)) / (2 * h)
   - Metode ini memperhitungkan titik di sebelah kiri (x - h) dan
\end{eulerttcomment}
\begin{eulercomment}
kanan (x + h).

3. Kesimpulan:\\
\end{eulercomment}
\begin{eulerttcomment}
   - diffc jauh lebih akurat dengan tingkat galat O(h^2) dibandingkan
\end{eulerttcomment}
\begin{eulercomment}
diff yang memiliki galat O(h).\\
\end{eulercomment}
\begin{eulerttcomment}
   - Oleh karena itu, hasil numerik diffc(f, 3) lebih mendekati nilai
\end{eulerttcomment}
\begin{eulercomment}
turunan analitis yang sebenarnya dibanding diff(f, 3).
\end{eulercomment}
\begin{eulerprompt}
>$showev('diff(f(x),x))
\end{eulerprompt}
\begin{eulerformula}
\[
\frac{d}{d\,x}\,\sin ^2\left(3\,x^5+7\right)=30\,x^4\,\cos \left(3  \,x^5+7\right)\,\sin \left(3\,x^5+7\right)
\]
\end{eulerformula}
\begin{eulerprompt}
>$% with x=3
\end{eulerprompt}
\begin{eulerformula}
\[
{\it \%at}\left(\frac{d}{d\,x}\,\sin ^2\left(3\,x^5+7\right) , x=3  \right)=2430\,\cos 736\,\sin 736
\]
\end{eulerformula}
\begin{eulerprompt}
>$float(%)
\end{eulerprompt}
\begin{eulerformula}
\[
{\it \%at}\left(\frac{d^{1.0}}{d\,x^{1.0}}\,\sin ^2\left(3.0\,x^5+  7.0\right) , x=3.0\right)=1198.728637211748
\]
\end{eulerformula}
\begin{eulerprompt}
>plot2d(f,0,3.1):
\end{eulerprompt}
\eulerimg{17}{Image Use/TugasAplikom6Adi25030630001-102.png}
\begin{eulerprompt}
>function f(x) &=5*cos(2*x)-2*x*sin(2*x) // mendifinisikan fungsi f
\end{eulerprompt}
\begin{euleroutput}
  
                        5 cos(2 x) - 2 x sin(2 x)
  
\end{euleroutput}
\begin{eulerprompt}
>function df(x) &=diff(f(x),x) // fd(x) = f'(x)
\end{eulerprompt}
\begin{euleroutput}
  
                       - 12 sin(2 x) - 4 x cos(2 x)
  
\end{euleroutput}
\begin{eulerprompt}
>$'f(1)=f(1), $float(f(1)), $'f(2)=f(2), $float(f(2)) // nilai f(1) dan f(2)
\end{eulerprompt}
\begin{eulerformula}
\[
-0.2410081230863468
\]
\end{eulerformula}
\eulerimg{0}{Image Use/TugasAplikom6Adi25030630001-104-large.png}
\eulerimg{0}{Image Use/TugasAplikom6Adi25030630001-105-large.png}
\eulerimg{0}{Image Use/TugasAplikom6Adi25030630001-106-large.png}
\begin{eulerprompt}
>xp=solve("df(x)",1,2,0) // solusi f'(x)=0 pada interval [1, 2]
\end{eulerprompt}
\begin{euleroutput}
  1.35822987384
\end{euleroutput}
\begin{eulerprompt}
>df(xp), f(xp) // cek bahwa f'(xp)=0 dan nilai ekstrim di titik tersebut
\end{eulerprompt}
\begin{euleroutput}
  0
  -5.67530133759
\end{euleroutput}
\begin{eulerprompt}
>plot2d(["f(x)","df(x)"],0,2*pi,color=[blue,red]): //grafik fungsi dan turunannya
\end{eulerprompt}
\eulerimg{17}{Image Use/TugasAplikom6Adi25030630001-107.png}
\begin{eulercomment}
Perhatikan titik-titik "puncak" grafik y=f(x) dan nilai turunan pada
saat grafik fungsinya mencapai titik "puncak" tersebut.

Penjelasan Analitis Hubungan Titik Puncak f(x) dan Turunannya df(x):

1. Hubungan Titik Puncak (Ekstremum) dengan Turunan Pertama:\\
\end{eulercomment}
\begin{eulerttcomment}
   - Ketika grafik fungsi y = f(x) (kurva biru) mencapai titik
\end{eulerttcomment}
\begin{eulercomment}
"puncak"\\
\end{eulercomment}
\begin{eulerttcomment}
     (baik maksimum lokal maupun minimum lokal), garis singgungnya
\end{eulerttcomment}
\begin{eulercomment}
mendatar.\\
\end{eulercomment}
\begin{eulerttcomment}
   - Hal ini berarti kemiringan (gradien) garis singgung pada titik
\end{eulerttcomment}
\begin{eulercomment}
puncak bernilai nol.

2. Pengamatan pada Grafik:\\
\end{eulercomment}
\begin{eulerttcomment}
   - Pada saat kurva f(x) (biru) berada di puncak maksimum (sekitar x
\end{eulerttcomment}
\begin{eulercomment}
˜ 2.8 dan x ˜ 5.8)\\
\end{eulercomment}
\begin{eulerttcomment}
     maupun puncak minimum (sekitar x ˜ 1.2 dan x ˜ 4.3), grafik
\end{eulerttcomment}
\begin{eulercomment}
turunannya df(x) (merah)\\
\end{eulercomment}
\begin{eulerttcomment}
     tepat memotong sumbu-X (bernilai nol, df(x) = 0).
\end{eulerttcomment}
\begin{eulercomment}

3. Kesimpulan Matematika:\\
\end{eulercomment}
\begin{eulerttcomment}
   - Syarat stasioner untuk menentukan titik puncak suatu fungsi f(x)
\end{eulerttcomment}
\begin{eulercomment}
adalah:\\
\end{eulercomment}
\begin{eulerttcomment}
     f'(x) = 0  atau  df(x) = 0
   - Grafik menunjukkan secara visual bahwa titik stasioner (puncak)
\end{eulerttcomment}
\begin{eulercomment}
f(x) selalu\\
\end{eulercomment}
\begin{eulerttcomment}
     berkorespondensi dengan pembuat nol (titik potong sumbu-X) dari
\end{eulerttcomment}
\begin{eulercomment}
fungsi turunannya df(x).
\end{eulercomment}
\begin{eulerprompt}
>  
\end{eulerprompt}
\eulerheading{Latihan}
\begin{eulercomment}
Bukalah buku Kalkulus. Cari dan pilih beberapa (paling sedikit 5
fungsi berbeda tipe/bentuk/jenis) fungsi dari buku tersebut, kemudian
definisikan di EMT pada baris-baris perintah berikut (jika perlu
tambahkan lagi). Untuk setiap fungsi, tentukan turunannya dengan
menggunakan definisi turunan (limit), menggunakan perintah diff, dan
secara manual (langkah demi langkah yang dihitung dengan Maxima)
seperti contoh-contoh di atas. Gambar grafik fungsi asli dan fungsi
turunannya pada sumbu koordinat yang sama.


\end{eulercomment}
\begin{eulerttcomment}
 1. FUNGSI POLINOMIAL: f(x) = 2*x^3 - 3*x^2 + 5
\end{eulerttcomment}
\begin{eulerprompt}
>function f1(x) := 2*x^3 - 3*x^2 + 5
\end{eulerprompt}
\begin{eulercomment}
a. Menggunakan definisi limit
\end{eulercomment}
\begin{eulerprompt}
>$showev('limit((2*(x+h)^3 - 3*(x+h)^2 + 5 - (2*x^3 - 3*x^2 + 5))/h, h, 0))
\end{eulerprompt}
\begin{eulerformula}
\[
\lim_{h\rightarrow 0}{\frac{2\,\left(x+h\right)^3-3\,\left(x+h  \right)^2-2\,x^3+3\,x^2}{h}}=6\,x^2-6\,x
\]
\end{eulerformula}
\begin{eulercomment}
b. Menggunakan perintah diff
\end{eulercomment}
\begin{eulerprompt}
>$showev('diff(2*x^3 - 3*x^2 + 5, x))
\end{eulerprompt}
\begin{eulerformula}
\[
\frac{d}{d\,x}\,\left(2\,x^3-3\,x^2+5\right)=6\,x^2-6\,x
\]
\end{eulerformula}
\begin{eulercomment}
c. Langkah manual
\end{eulercomment}
\begin{eulerprompt}
>&define(f1(x), 2*x^3 - 3*x^2 + 5);
>$showev('diff(f1(x), x))
\end{eulerprompt}
\begin{eulerformula}
\[
\frac{d}{d\,x}\,\left(2\,x^3-3\,x^2+5\right)=6\,x^2-6\,x
\]
\end{eulerformula}
\begin{eulercomment}
d. Plot grafik fungsi asli dan turunannya
\end{eulercomment}
\begin{eulerprompt}
>plot2d(["2*x^3 - 3*x^2 + 5", "6*x^2 - 6*x"], -2, 2, color=[blue, red], grid=1, title="1. Polinomial: f(x) dan f'(x)");
\end{eulerprompt}
\begin{eulercomment}
2. FUNGSI RASIONAL: f(x) = (x + 1)/(x - 2)
\end{eulercomment}
\begin{eulerprompt}
>function f2(x) := (x + 1)/(x - 2)
\end{eulerprompt}
\begin{eulercomment}
a. Menggunakan definisi limit
\end{eulercomment}
\begin{eulerprompt}
>$showev('limit((((x+h)+1)/((x+h)-2) - (x+1)/(x-2))/h, h, 0))
\end{eulerprompt}
\begin{eulerformula}
\[
\lim_{h\rightarrow 0}{\frac{\frac{x+h+1}{x+h-2}-\frac{x+1}{x-2}}{h}  }=-\frac{3}{x^2-4\,x+4}
\]
\end{eulerformula}
\begin{eulercomment}
b. Menggunakan perintah diff
\end{eulercomment}
\begin{eulerprompt}
>$showev('diff((x + 1)/(x - 2), x))
\end{eulerprompt}
\begin{eulerformula}
\[
\frac{d}{d\,x}\,\left(\frac{x+1}{x-2}\right)=\frac{1}{x-2}-\frac{x+  1}{\left(x-2\right)^2}
\]
\end{eulerformula}
\begin{eulercomment}
c. Langkah manual
\end{eulercomment}
\begin{eulerprompt}
>&u: x + 1; &v: x - 2;
>$showev(('diff(u,x)*v - u*'diff(v,x))/v^2)
\end{eulerprompt}
\begin{eulerformula}
\[
\frac{\left(x-2\right)\,\left(\frac{d}{d\,x}\,\left(x+1\right)  \right)-\left(x+1\right)\,\left(\frac{d}{d\,x}\,\left(x-2\right)  \right)}{\left(x-2\right)^2}=-\frac{3}{\left(x-2\right)^2}
\]
\end{eulerformula}
\begin{eulercomment}
d. Plot grafik fungsi asli dan turunannya 
\end{eulercomment}
\begin{eulerprompt}
> plot2d(["(x+1)/(x-2)", "-3/(x-2)^2"], -3, 5, c=-10, d=10, color=[blue, red], grid=1, title="2. Rasional");
\end{eulerprompt}
\begin{eulercomment}
3. FUNGSI TRIGONOMETRI: f(x) = sin(2*x)
\end{eulercomment}
\begin{eulerprompt}
> function f3(x) := sin(2*x)
\end{eulerprompt}
\begin{eulercomment}
a. Menggunakan definisi limit
\end{eulercomment}
\begin{eulerprompt}
>$showev('limit((sin(2*(x+h)) - sin(2*x))/h, h, 0))
\end{eulerprompt}
\begin{eulerformula}
\[
\lim_{h\rightarrow 0}{\frac{\sin \left(2\,\left(x+h\right)\right)-  \sin \left(2\,x\right)}{h}}=2\,\cos \left(2\,x\right)
\]
\end{eulerformula}
\begin{eulercomment}
b. Menggunakan perintah diff
\end{eulercomment}
\begin{eulerprompt}
>$showev('diff(sin(2*x), x))
\end{eulerprompt}
\begin{eulerformula}
\[
\frac{d}{d\,x}\,\sin \left(2\,x\right)=2\,\cos \left(2\,x\right)
\]
\end{eulerformula}
\begin{eulercomment}
c. Langkah manual
\end{eulercomment}
\begin{eulerprompt}
>remvalue u; &depends(u, x);
>$showev('diff(sin(u), u) * 'diff(2*x, x))
\end{eulerprompt}
\begin{eulerformula}
\[
\frac{d}{d\,u}\,\sin u\,\left(\frac{d}{d\,x}\,\left(2\,x\right)  \right)=2\,\cos u
\]
\end{eulerformula}
\begin{eulerprompt}
>$showev(subst(u=2*x, 2*cos(u)))
\end{eulerprompt}
\begin{eulerformula}
\[
2\,\cos \left(2\,x\right)=2\,\cos \left(2\,x\right)
\]
\end{eulerformula}
\begin{eulercomment}
d. Plot grafik fungsi asli dan turunannya
\end{eulercomment}
\begin{eulerprompt}
>plot2d(["sin(2*x)", "2*cos(2*x)"], 0, 2*pi, color=[blue, red], grid=1, title="3. Trigonometri: f(x) dan f'(x)");
\end{eulerprompt}
\begin{eulercomment}
4. FUNGSI EKSPONENSIAL: f(x) = \%e\textasciicircum{}(3*x)
\end{eulercomment}
\begin{eulerprompt}
>function f4(x) := exp(3*x)
\end{eulerprompt}
\begin{eulercomment}
a. Menggunakan definisi limit
\end{eulercomment}
\begin{eulerprompt}
>$&'limit((%e^(3*(x+h)) - %e^(3*x))/h, h, 0) = 3*%e^(3*x)
\end{eulerprompt}
\begin{eulerformula}
\[
\lim_{h\rightarrow 0}{\frac{e^{3\,\left(x+h\right)}-e^{3\,x}}{h}}=3  \,e^{3\,x}
\]
\end{eulerformula}
\begin{eulerprompt}
>  
\end{eulerprompt}
\begin{eulercomment}
b. Perintah diff
\end{eulercomment}
\begin{eulerprompt}
>$showev('diff(%e^(3*x), x))
\end{eulerprompt}
\begin{eulerformula}
\[
\frac{d}{d\,x}\,e^{3\,x}=3\,e^{3\,x}
\]
\end{eulerformula}
\begin{eulercomment}
c. Langkah manual
\end{eulercomment}
\begin{eulerprompt}
>$showev('diff(%e^u, u) * 'diff(3*x, x))
\end{eulerprompt}
\begin{eulerformula}
\[
\frac{d}{d\,u}\,e^{u}\,\left(\frac{d}{d\,x}\,\left(3\,x\right)  \right)=3\,e^{u}
\]
\end{eulerformula}
\begin{eulercomment}
d. Plot grafik
\end{eulercomment}
\begin{eulerprompt}
>plot2d(["exp(3*x)", "3*exp(3*x)"], -1, 1, color=[blue, red], grid=1, title="4. Eksponensial");
\end{eulerprompt}
\begin{eulercomment}
5. FUNGSI LOGARITMA NATURAL: f(x) = log(x\textasciicircum{}2 + 1)
\end{eulercomment}
\begin{eulerprompt}
>&assume(x>0);
>function f5(x) := log(x^2 + 1)
\end{eulerprompt}
\begin{eulercomment}
a. Menggunakan definisi limit
\end{eulercomment}
\begin{eulerprompt}
>$showev('limit((log((x+h)^2 + 1) - log(x^2 + 1))/h, h, 0))
\end{eulerprompt}
\begin{eulerformula}
\[
\lim_{h\rightarrow 0}{\frac{\log \left(\left(x+h\right)^2+1\right)-  \log \left(x^2+1\right)}{h}}=\frac{2\,x}{x^2+1}
\]
\end{eulerformula}
\begin{eulercomment}
b. Menggunakan perintah diff
\end{eulercomment}
\begin{eulerprompt}
>$showev('diff(log(x^2 + 1), x))
\end{eulerprompt}
\begin{eulerformula}
\[
\frac{d}{d\,x}\,\log \left(x^2+1\right)=\frac{2\,x}{x^2+1}
\]
\end{eulerformula}
\begin{eulercomment}
c. Langkah manual
\end{eulercomment}
\begin{eulerprompt}
>remvalue u;
>$showev('diff(log(u), u) * 'diff(x^2 + 1, x))
\end{eulerprompt}
\begin{eulerformula}
\[
\frac{d}{d\,u}\,\log u\,\left(\frac{d}{d\,x}\,\left(x^2+1\right)  \right)=\frac{2\,x}{u}
\]
\end{eulerformula}
\begin{eulerprompt}
>$showev(subst(u=x^2+1, (1/u)*(2*x)))
\end{eulerprompt}
\begin{eulerformula}
\[
\frac{2\,x}{x^2+1}=\frac{2\,x}{x^2+1}
\]
\end{eulerformula}
\begin{eulercomment}
d. Plot grafik
\end{eulercomment}
\begin{eulerprompt}
>plot2d(["log(x^2 + 1)", "(2*x)/(x^2 + 1)"], -3, 3, color=[blue, red], grid=1, title="5. Logaritma: f(x) dan f'(x)");
\end{eulerprompt}
\eulerheading{Integral}
\begin{eulercomment}
EMT dapat digunakan untuk menghitung integral, baik integral tak tentu
maupun integral tentu. Untuk integral tak tentu (simbolik) sudah tentu
EMT menggunakan Maxima, sedangkan untuk perhitungan integral tentu EMT
sudah menyediakan beberapa fungsi yang mengimplementasikan algoritma
kuadratur (perhitungan integral tentu menggunakan metode numerik).

Pada notebook ini akan ditunjukkan perhitungan integral tentu dengan
menggunakan Teorema Dasar Kalkulus:

\end{eulercomment}
\begin{eulerformula}
\[
\int_a^b f(x)\ dx = F(b)-F(a), \quad \text{ dengan  } F'(x) = f(x).
\]
\end{eulerformula}
\begin{eulercomment}
Fungsi untuk menentukan integral adalah integrate. Fungsi ini dapat
digunakan untuk menentukan, baik integral tentu maupun tak tentu (jika
fungsinya memiliki antiderivatif). Untuk perhitungan integral tentu
fungsi integrate menggunakan metode numerik (kecuali fungsinya tidak
integrabel, kita tidak akan menggunakan metode ini).
\end{eulercomment}
\begin{eulerprompt}
>$showev('integrate(x^n,x))
\end{eulerprompt}
\begin{euleroutput}
  Answering "Is n equal to -1?" with "no"
\end{euleroutput}
\begin{eulerformula}
\[
\int {x^{n}}{\;dx}=\frac{x^{n+1}}{n+1}
\]
\end{eulerformula}
\begin{eulerprompt}
>$showev('integrate(1/(1+x),x))
\end{eulerprompt}
\begin{eulerformula}
\[
\int {\frac{1}{x+1}}{\;dx}=\log \left(x+1\right)
\]
\end{eulerformula}
\begin{eulerprompt}
>$showev('integrate(1/(1+x^2),x))
\end{eulerprompt}
\begin{eulerformula}
\[
\int {\frac{1}{x^2+1}}{\;dx}=\arctan x
\]
\end{eulerformula}
\begin{eulerprompt}
>$showev('integrate(1/sqrt(1-x^2),x))
\end{eulerprompt}
\begin{eulerformula}
\[
\int {\frac{1}{\sqrt{1-x^2}}}{\;dx}=\arcsin x
\]
\end{eulerformula}
\begin{eulerprompt}
>$showev('integrate(sin(x),x,0,pi))
\end{eulerprompt}
\begin{eulerformula}
\[
\int_{0}^{\pi}{\sin x\;dx}=2
\]
\end{eulerformula}
\begin{eulerprompt}
>plot2d("sin(x)",0,2*pi):
\end{eulerprompt}
\eulerimg{17}{Image Use/TugasAplikom6Adi25030630001-131.png}
\begin{eulerprompt}
>$showev('integrate(sin(x),x,a,b))
\end{eulerprompt}
\begin{eulerformula}
\[
\int_{a}^{b}{\sin x\;dx}=\cos a-\cos b
\]
\end{eulerformula}
\begin{eulerprompt}
>$showev('integrate(x^n,x,a,b))
\end{eulerprompt}
\begin{euleroutput}
  Answering "Is n positive, negative or zero?" with "positive"
\end{euleroutput}
\begin{eulerformula}
\[
\int_{a}^{b}{x^{n}\;dx}=\frac{b^{n+1}}{n+1}-\frac{a^{n+1}}{n+1}
\]
\end{eulerformula}
\begin{eulerprompt}
>$showev('integrate(x^2*sqrt(2*x+1),x))
\end{eulerprompt}
\begin{eulerformula}
\[
\int {x^2\,\sqrt{2\,x+1}}{\;dx}=\frac{\left(2\,x+1\right)^{\frac{7  }{2}}}{28}-\frac{\left(2\,x+1\right)^{\frac{5}{2}}}{10}+\frac{\left(  2\,x+1\right)^{\frac{3}{2}}}{12}
\]
\end{eulerformula}
\begin{eulerprompt}
>$showev('integrate(x^2*sqrt(2*x+1),x,0,2))
\end{eulerprompt}
\begin{eulerformula}
\[
\int_{0}^{2}{x^2\,\sqrt{2\,x+1}\;dx}=\frac{2\,5^{\frac{5}{2}}}{21}-  \frac{2}{105}
\]
\end{eulerformula}
\begin{eulerprompt}
>$ratsimp(%)
\end{eulerprompt}
\begin{eulerformula}
\[
\int_{0}^{2}{x^2\,\sqrt{2\,x+1}\;dx}=\frac{2\,5^{\frac{7}{2}}-2}{  105}
\]
\end{eulerformula}
\begin{eulerprompt}
>$showev('integrate((sin(sqrt(x)+a)*E^sqrt(x))/sqrt(x),x,0,pi^2))
\end{eulerprompt}
\begin{eulerformula}
\[
\int_{0}^{\pi^2}{\frac{\sin \left(\sqrt{x}+a\right)\,e^{\sqrt{x}}}{  \sqrt{x}}\;dx}=\left(-e^{\pi}-1\right)\,\sin a+\left(e^{\pi}+1  \right)\,\cos a
\]
\end{eulerformula}
\begin{eulerprompt}
>$factor(%)
\end{eulerprompt}
\begin{eulerformula}
\[
\int_{0}^{\pi^2}{\frac{\sin \left(\sqrt{x}+a\right)\,e^{\sqrt{x}}}{  \sqrt{x}}\;dx}=\left(-e^{\pi}-1\right)\,\left(\sin a-\cos a\right)
\]
\end{eulerformula}
\begin{eulerprompt}
>function map f(x) &= E^(-x^2)
\end{eulerprompt}
\begin{euleroutput}
  
                                      2
                                   - x
                                  E
  
\end{euleroutput}
\begin{eulerprompt}
>$showev('integrate(f(x),x))
\end{eulerprompt}
\begin{eulerformula}
\[
\int {e^ {- x^2 }}{\;dx}=\frac{\sqrt{\pi}\,\mathrm{erf}\left(x  \right)}{2}
\]
\end{eulerformula}
\begin{eulercomment}
Fungsi f tidak memiliki antiturunan, integralnya masih memuat integral
lain.

\end{eulercomment}
\begin{eulerformula}
\[
erf(x) = \int \frac{e^{-x^2}}{\sqrt{\pi}} \ dx.
\]
\end{eulerformula}
\begin{eulercomment}
Kita tidak dapat menggunakan teorema Dasar kalkulus untuk menghitung
integral tentu fungsi tersebut jika semua batasnya berhingga. Dalam
hal ini dapat digunakan metode numerik (rumus kuadratur).

Misalkan kita akan menghitung:

\end{eulercomment}
\begin{eulerformula}
\[
\int_{0}^{\pi}{x^{x}\;dx}
\]
\end{eulerformula}
\begin{eulercomment}
\end{eulercomment}
\begin{eulerprompt}
>x=0:0.1:pi-0.1; plot2d(x,f(x+0.1),>bar); plot2d("f(x)",0,pi,>add):
\end{eulerprompt}
\eulerimg{17}{Image Use/TugasAplikom6Adi25030630001-142.png}
\begin{eulercomment}
Integral tentu

\end{eulercomment}
\begin{eulerformula}
\[
\int_{0}^{\pi}{e^ {- x^2 }\;dx}
\]
\end{eulerformula}
\begin{eulercomment}
dapat dihampiri dengan jumlah luas persegi-persegi panjang di bawah
kurva y=f(x) tersebut. Langkah-langkahnya adalah sebagai berikut.
\end{eulercomment}
\begin{eulerprompt}
>t &= makelist(a,a,0,pi-0.1,0.1); // t sebagai list untuk menyimpan nilai-nilai x
>fx &= makelist(f(t[i]+0.1),i,1,length(t)); // simpan nilai-nilai f(x)
>// jangan menggunakan x sebagai list, kecuali Anda pakar Maxima!
\end{eulerprompt}
\begin{eulercomment}
Hasilnya adalah:

\end{eulercomment}
\begin{eulerformula}
\[
\int_{0}^{\pi}{x^{x}\;dx}=19.19514483551471
\]
\end{eulerformula}
\begin{eulercomment}
Jumlah tersebut diperoleh dari hasil kali lebar sub-subinterval (=0.1)
dan jumlah nilai-nilai f(x) untuk x = 0.1, 0.2, 0.3, ..., 3.2.
\end{eulercomment}
\begin{eulerprompt}
>0.1*sum(f(x+0.1)) // cek langsung dengan perhitungan numerik EMT
\end{eulerprompt}
\begin{euleroutput}
  0.836219610253
\end{euleroutput}
\begin{eulercomment}
Untuk mendapatkan nilai integral tentu yang mendekati nilai sebenarnya, lebar
sub-intervalnya dapat diperkecil lagi, sehingga daerah di bawah kurva tertutup
semuanya, misalnya dapat digunakan lebar subinterval 0.001. (Silakan dicoba!)

Meskipun Maxima tidak dapat menghitung integral tentu fungsi tersebut untuk
batas-batas yang berhingga, namun integral tersebut dapat dihitung secara eksak jika
batas-batasnya tak hingga. Ini adalah salah satu keajaiban di dalam matematika, yang
terbatas tidak dapat dihitung secara eksak, namun yang tak hingga malah dapat
dihitung secara eksak.
\end{eulercomment}
\begin{eulerprompt}
>$showev('integrate(f(x),x,0,inf))
\end{eulerprompt}
\begin{eulerformula}
\[
\int_{0}^{\infty }{e^ {- x^2 }\;dx}=\frac{\sqrt{\pi}}{2}
\]
\end{eulerformula}
\begin{eulercomment}
Tunjukkan kebenaran hasil di atas!

PEMBUKTIAN ANALITIS INTEGRAL GAUSS (KOORDINAT POLAR)\\
\end{eulercomment}
\begin{eulerttcomment}
 Misalkan I = int(e^(-x^2), x, 0, inf)
\end{eulerttcomment}
\begin{eulercomment}

Kuadratkan integral menjadi integral lipat dua di kuadran pertama:

I\textasciicircum{}2 = int(int(e\textasciicircum{}(-(x\textasciicircum{}2+y\textasciicircum{}2)), x, 0, inf), y, 0, inf)


\end{eulercomment}
\begin{eulerttcomment}
 Transformasi ke koordinat polar (x = r cos(theta), y = r sin(theta),
\end{eulerttcomment}
\begin{eulercomment}
dx dy = r dr dtheta):

Batas r: 0 hingga inf, Batas theta: 0 hingga pi/2


\end{eulercomment}
\begin{eulerttcomment}
 I^2 = int(dtheta, theta, 0, pi/2) * int(r * e^(-r^2), r, 0, inf)
\end{eulerttcomment}
\begin{eulercomment}

\end{eulercomment}
\begin{eulerttcomment}
    = (pi / 2) * [ -1/2 * e^(-r^2) ]_(0)^(inf)
\end{eulerttcomment}
\begin{eulercomment}

\end{eulercomment}
\begin{eulerttcomment}
    = (pi / 2) * (1 / 2)
\end{eulerttcomment}
\begin{eulercomment}

\end{eulercomment}
\begin{eulerttcomment}
    = pi / 4
\end{eulerttcomment}
\begin{eulercomment}


\end{eulercomment}
\begin{eulerttcomment}
 Akar kuadrat dari kedua sisi:
\end{eulerttcomment}
\begin{eulercomment}

I = sqrt(pi) / 2  [TERBUKTI]

Berikut adalah contoh lain fungsi yang tidak memiliki antiderivatif,
sehingga integral tentunya hanya dapat dihitung dengan metode numerik.
\end{eulercomment}
\begin{eulerprompt}
>function f(x) &= x^x
\end{eulerprompt}
\begin{euleroutput}
  
                                     x
                                    x
  
\end{euleroutput}
\begin{eulerprompt}
>$showev('integrate(f(x),x,0,1))
\end{eulerprompt}
\begin{eulerformula}
\[
\int_{0}^{1}{x^{x}\;dx}=\int_{0}^{1}{x^{x}\;dx}
\]
\end{eulerformula}
\begin{eulerprompt}
>x=0:0.1:1-0.01; plot2d(x,f(x+0.01),>bar); plot2d("f(x)",0,1,>add):
\end{eulerprompt}
\eulerimg{17}{Image Use/TugasAplikom6Adi25030630001-147.png}
\begin{eulercomment}
Maxima gagal menghitung integral tentu tersebut secara langsung menggunakan perintah
integrate. Berikut kita lakukan seperti contoh sebelumnya untuk mendapat hasil atau
pendekatan nilai integral tentu tersebut.
\end{eulercomment}
\begin{eulerprompt}
>t &= makelist(a,a,0,1-0.01,0.01);
>fx &= makelist(f(t[i]+0.01),i,1,length(t));
\end{eulerprompt}
\begin{eulerformula}
\[
\int_{0}^{1}{x^{x}\;dx}=0.7834935879025506
\]
\end{eulerformula}
\begin{eulercomment}
Apakah hasil tersebut cukup baik? perhatikan gambarnya.

Hasil pendekatan numerik tersebut CUKUP BAIK sebagai hampiran awal
(berhasil\\
endapatkan nilai hingga 3 digit desimal pertama yaitu 0.783), namun
BELUM\\
ANGAT AKURAT jika dibandingkan dengan nilai eksak teoretisnya.


ANALISIS \& ALASAN MATEMATIS:

. PERBANDINGAN NILAI EKSAK VS NUMERIK\\
\end{eulercomment}
\begin{eulerttcomment}
  * Nilai Hampiran EMT:  0.7834935879025506
  * Nilai Eksak Deret:   0.783430510712...
    (Menggunakan rumus deret Sophomore's Dream: sum((-1)^(n-1) *
\end{eulerttcomment}
\begin{eulercomment}
n\textasciicircum{}(-n)))\\
\end{eulercomment}
\begin{eulerttcomment}
  * Galat Absolut:       |0.7834935879 - 0.7834305107| = 0.0000630772
\end{eulerttcomment}
\begin{eulercomment}
(0.0063\%)\\
2. ANALISIS VISUAL GRAFIK (PENYEBAB OVERESTIMASI)\\
\end{eulercomment}
\begin{eulerttcomment}
  * Grafik fungsi f(x) = x^x bersifat monoton turun pada interval x in
\end{eulerttcomment}
\begin{eulercomment}
[0, 0.3679].\\
\end{eulercomment}
\begin{eulerttcomment}
  * Penggunaan Jumlah Riemann Kanan pada kurva yang sedang turun
\end{eulerttcomment}
\begin{eulercomment}
menyebabkan\\
tinggi batang hijau berada di ATAS kurva f(x).\\
\end{eulercomment}
\begin{eulerttcomment}
  * Terlihat jelas pada ujung kiri ( interval [0, 0.1] ), area batang
\end{eulerttcomment}
\begin{eulercomment}
hijau\\
menonjol keluar melebihi kurva, sehingga luas hampiran menjadi LEBIH
BESAR\\
dari luas daerah yang sebenarnya.

REKOMENDASI PERBAIKAN DI EMT:

Untuk mendapatkan hasil yang jauh lebih akurat, kita dapat:\\
. Memperkecil lebar sub-interval (misalnya menggunakan dx = 0.0001).\\
. Menggunakan Aturan Trapesium atau Aturan Simpson daripada Riemann
Kanan.\\
. Menggunakan fungsi integrasi numerik bawaan EMT: integrate("x\textasciicircum{}x", 0,
1).
\end{eulercomment}
\begin{eulerprompt}
>function f(x) &= sin(3*x^5+7)^2
\end{eulerprompt}
\begin{euleroutput}
  
                                 2    5
                              sin (3 x  + 7)
  
\end{euleroutput}
\begin{eulerprompt}
>integrate(f,0,1)
\end{eulerprompt}
\begin{euleroutput}
  0.542581176074
\end{euleroutput}
\begin{eulerprompt}
>&showev('integrate(f(x),x,0,1))
\end{eulerprompt}
\begin{euleroutput}
  
           1                           1              pi
          /                      gamma(-) sin(14) sin(--)
          [     2    5                 5              10
          I  sin (3 x  + 7) dx = ------------------------
          ]                                  1/5
          /                              10 6
           0
         4/5                  1          4/5                  1
   - (((6    gamma_incomplete(-, 6 I) + 6    gamma_incomplete(-, - 6 I))
                              5                               5
               4/5                    1
   sin(14) + (6    I gamma_incomplete(-, 6 I)
                                      5
      4/5                    1                       pi
   - 6    I gamma_incomplete(-, - 6 I)) cos(14)) sin(--) - 60)/120
                             5                       10
  
\end{euleroutput}
\begin{eulerprompt}
>&float(%)
\end{eulerprompt}
\begin{euleroutput}
  
           1.0
          /
          [       2      5
          I    sin (3.0 x  + 7.0) dx = 
          ]
          /
           0.0
  0.09820784258795788 - 0.008333333333333333
   (0.3090169943749474 (0.1367372182078336
   (4.192962712629476 I gamma__incomplete(0.2, 6.0 I)
   - 4.192962712629476 I gamma__incomplete(0.2, - 6.0 I))
   + 0.9906073556948704 (4.192962712629476 gamma__incomplete(0.2, 6.0 I)
   + 4.192962712629476 gamma__incomplete(0.2, - 6.0 I))) - 60.0)
  
\end{euleroutput}
\begin{eulerprompt}
>$showev('integrate(x*exp(-x),x,0,1)) // Integral tentu (eksak)
\end{eulerprompt}
\begin{eulerformula}
\[
\int_{0}^{1}{x\,e^ {- x }\;dx}=1-2\,e^ {- 1 }
\]
\end{eulerformula}
\eulerheading{Aplikasi Integral Tentu}
\begin{eulerprompt}
>plot2d("x^3-x",-0.1,1.1); plot2d("-x^2",>add);  ...
>b=solve("x^3-x+x^2",0.5); x=linspace(0,b,200); xi=flipx(x); ...
>plot2d(x|xi,x^3-x|-xi^2,>filled,style="|",fillcolor=1,>add): // Plot daerah antara 2 kurva
\end{eulerprompt}
\eulerimg{17}{Image Use/TugasAplikom6Adi25030630001-150.png}
\begin{eulerprompt}
>a=solve("x^3-x+x^2",0), b=solve("x^3-x+x^2",1) // absis titik-titik potong kedua kurva
\end{eulerprompt}
\begin{euleroutput}
  0
  0.61803398875
\end{euleroutput}
\begin{eulerprompt}
>integrate("(-x^2)-(x^3-x)",a,b) // luas daerah yang diarsir
\end{eulerprompt}
\begin{euleroutput}
  0.0758191713542
\end{euleroutput}
\begin{eulercomment}
Hasil tersebut akan kita bandingkan dengan perhitungan secara analitik.
\end{eulercomment}
\begin{eulerprompt}
>a &= solve((-x^2)-(x^3-x),x); $a // menentukan absis titik potong kedua kurva secara eksak
\end{eulerprompt}
\begin{eulerformula}
\[
\left[ x=\frac{-\sqrt{5}-1}{2} , x=\frac{\sqrt{5}-1}{2} , x=0   \right] 
\]
\end{eulerformula}
\begin{eulerprompt}
>$showev('integrate(-x^2-x^3+x,x,0,(sqrt(5)-1)/2)) // Nilai integral secara eksak
\end{eulerprompt}
\begin{eulerformula}
\[
\int_{0}^{\frac{\sqrt{5}-1}{2}}{-x^3-x^2+x\;dx}=\frac{13-5^{\frac{3  }{2}}}{24}
\]
\end{eulerformula}
\begin{eulerprompt}
>$float(%)
\end{eulerprompt}
\begin{eulerformula}
\[
\int_{0.0}^{0.6180339887498949}{-1.0\,x^3-1.0\,x^2+x\;dx}=  0.07581917135421037
\]
\end{eulerformula}
\begin{eulercomment}
KESIMPULAN PERBANDINGAN

1. Metode Numerik EMT menghasilkan nilai: 0.0758191713542\\
2. Metode Analitik (Eksak Maxima) menghasilkan nilai:
0.07581917135421037\\
Kesimpulan: Kedua metode menunjukkan kesesuaian yang sangat tinggi
(identik hingga 11 digit desimal), membuktikan bahwa hampiran numerik
EMT sangat akurat.
\end{eulercomment}
\eulersubheading{Panjang Kurva}
\begin{eulercomment}
Hitunglah panjang kurva berikut ini dan luas daerah di dalam kurva
tersebut.

\end{eulercomment}
\begin{eulerformula}
\[
\gamma(t) = (r(t) \cos(t), r(t) \sin(t))
\]
\end{eulerformula}
\begin{eulercomment}
dengan

\end{eulercomment}
\begin{eulerformula}
\[
r(t) = 1 + \dfrac{\sin(3t)}{2},\quad 0\le t\le 2\pi.
\]
\end{eulerformula}
\begin{eulerprompt}
>t=linspace(0,2pi,1000); r=1+sin(3*t)/2; x=r*cos(t); y=r*sin(t); ...
>plot2d(x,y,>filled,fillcolor=red,style="/",r=1.5): // Kita gambar kurvanya terlebih dahulu
\end{eulerprompt}
\eulerimg{17}{Image Use/TugasAplikom6Adi25030630001-156.png}
\begin{eulerprompt}
>function r(t) &= 1+sin(3*t)/2; $'r(t)=r(t)
\end{eulerprompt}
\begin{eulerformula}
\[
r\left(\left[ 0 , 0.01 , 0.02 , 0.03 , 0.04 , 0.05 , 0.06 , 0.07 ,   0.08 , 0.09 , 0.1 , 0.11 , 0.12 , 0.13 , 0.14 , 0.15 , 0.16 , 0.17   , 0.18 , 0.19 , 0.2 , 0.21 , 0.2200000000000001 ,   0.2300000000000001 , 0.2400000000000001 , 0.2500000000000001 ,   0.2600000000000001 , 0.2700000000000001 , 0.2800000000000001 ,   0.2900000000000001 , 0.3000000000000001 , 0.3100000000000001 ,   0.3200000000000001 , 0.3300000000000001 , 0.3400000000000001 ,   0.3500000000000001 , 0.3600000000000002 , 0.3700000000000002 ,   0.3800000000000002 , 0.3900000000000002 , 0.4000000000000002 ,   0.4100000000000002 , 0.4200000000000002 , 0.4300000000000002 ,   0.4400000000000002 , 0.4500000000000002 , 0.4600000000000002 ,   0.4700000000000003 , 0.4800000000000003 , 0.4900000000000003 ,   0.5000000000000002 , 0.5100000000000002 , 0.5200000000000002 ,   0.5300000000000002 , 0.5400000000000003 , 0.5500000000000003 ,   0.5600000000000003 , 0.5700000000000003 , 0.5800000000000003 ,   0.5900000000000003 , 0.6000000000000003 , 0.6100000000000003 ,   0.6200000000000003 , 0.6300000000000003 , 0.6400000000000003 ,   0.6500000000000004 , 0.6600000000000004 , 0.6700000000000004 ,   0.6800000000000004 , 0.6900000000000004 , 0.7000000000000004 ,   0.7100000000000004 , 0.7200000000000004 , 0.7300000000000004 ,   0.7400000000000004 , 0.7500000000000004 , 0.7600000000000005 ,   0.7700000000000005 , 0.7800000000000005 , 0.7900000000000005 ,   0.8000000000000005 , 0.8100000000000005 , 0.8200000000000005 ,   0.8300000000000005 , 0.8400000000000005 , 0.8500000000000005 ,   0.8600000000000005 , 0.8700000000000006 , 0.8800000000000006 ,   0.8900000000000006 , 0.9000000000000006 , 0.9100000000000006 ,   0.9200000000000006 , 0.9300000000000006 , 0.9400000000000006 ,   0.9500000000000006 , 0.9600000000000006 , 0.9700000000000006 ,   0.9800000000000006 , 0.9900000000000007 \right] \right)=\left[ 1 ,   1.014997750101248 , 1.029982003239722 , 1.044939274599006 ,   1.05985610364446 , 1.0747190662368 , 1.089514786712912 ,   1.10422994992305 , 1.118851313213567 , 1.133365718344415 ,   1.14776010333067 , 1.162021514197434 , 1.176137116637545 ,   1.190094207561581 , 1.203880226529785 , 1.217482767055615 ,   1.230889587770742 , 1.244088623441454 , 1.257067995826556 ,   1.269816024366985 , 1.282321236697518 , 1.294572378971135 ,   1.306558425986717 , 1.318268591110984 , 1.329692335985737 ,   1.340819380011667 , 1.351639709600205 , 1.362143587185071 ,   1.37232155998543 , 1.382164468512753 , 1.391663454813742 ,   1.400809970441889 , 1.409595784150499 , 1.41801298930026 ,   1.426054010974682 , 1.433711612797009 , 1.440978903442474 ,   1.447849342840024 , 1.454316748057942 , 1.460375298868068 ,   1.466019542983613 , 1.471244400965849 , 1.476045170795258 ,   1.480417532103036 , 1.484357550059133 , 1.48786167891333 ,   1.49092676518618 , 1.493550050506925 , 1.495729174095843 ,   1.49746217488879 , 1.498747493302027 , 1.499583972635738 ,   1.499970860114983 , 1.499907807567145 , 1.499394871735262 ,   1.498432514226959 , 1.497021601099038 , 1.495163402078079 ,   1.492859589417777 , 1.490112236394023 , 1.486923815439098 ,   1.483297195916649 , 1.479235641539457 , 1.474742807432315 ,   1.469822736842662 , 1.464479857501934 , 1.458718977640905 ,   1.4525452816626 , 1.44596432547669 , 1.438982031499539 ,   1.431604683324436 , 1.423838920066784 , 1.415691730389341 ,   1.407170446212898 , 1.398282736118043 , 1.38903659844396 ,   1.379440354090461 , 1.369502639029735 , 1.359232396534563 ,   1.348638869129968 , 1.337731590275575 , 1.326520375786132 ,   1.315015314997945 , 1.303226761689157 , 1.29116532476204 ,   1.278841858695708 , 1.26626745377781 , 1.253453426124026 ,   1.240411307494323 , 1.227152834915152 , 1.213689940116914 ,   1.200034738796209 , 1.186199519712527 , 1.172196733629194 ,   1.158038982108526 , 1.143739006171271 , 1.129309674830555 ,   1.114763973510631 , 1.100114992360884 , 1.085375914475572 \right] 
\]
\end{eulerformula}
\begin{eulerprompt}
>function fx(t) &= r(t)*cos(t); $'fx(t)=fx(t)
\end{eulerprompt}
\begin{eulerformula}
\[
{\it fx}\left(\left[ 0 , 0.01 , 0.02 , 0.03 , 0.04 , 0.05 , 0.06 ,   0.07 , 0.08 , 0.09 , 0.1 , 0.11 , 0.12 , 0.13 , 0.14 , 0.15 , 0.16   , 0.17 , 0.18 , 0.19 , 0.2 , 0.21 , 0.2200000000000001 ,   0.2300000000000001 , 0.2400000000000001 , 0.2500000000000001 ,   0.2600000000000001 , 0.2700000000000001 , 0.2800000000000001 ,   0.2900000000000001 , 0.3000000000000001 , 0.3100000000000001 ,   0.3200000000000001 , 0.3300000000000001 , 0.3400000000000001 ,   0.3500000000000001 , 0.3600000000000002 , 0.3700000000000002 ,   0.3800000000000002 , 0.3900000000000002 , 0.4000000000000002 ,   0.4100000000000002 , 0.4200000000000002 , 0.4300000000000002 ,   0.4400000000000002 , 0.4500000000000002 , 0.4600000000000002 ,   0.4700000000000003 , 0.4800000000000003 , 0.4900000000000003 ,   0.5000000000000002 , 0.5100000000000002 , 0.5200000000000002 ,   0.5300000000000002 , 0.5400000000000003 , 0.5500000000000003 ,   0.5600000000000003 , 0.5700000000000003 , 0.5800000000000003 ,   0.5900000000000003 , 0.6000000000000003 , 0.6100000000000003 ,   0.6200000000000003 , 0.6300000000000003 , 0.6400000000000003 ,   0.6500000000000004 , 0.6600000000000004 , 0.6700000000000004 ,   0.6800000000000004 , 0.6900000000000004 , 0.7000000000000004 ,   0.7100000000000004 , 0.7200000000000004 , 0.7300000000000004 ,   0.7400000000000004 , 0.7500000000000004 , 0.7600000000000005 ,   0.7700000000000005 , 0.7800000000000005 , 0.7900000000000005 ,   0.8000000000000005 , 0.8100000000000005 , 0.8200000000000005 ,   0.8300000000000005 , 0.8400000000000005 , 0.8500000000000005 ,   0.8600000000000005 , 0.8700000000000006 , 0.8800000000000006 ,   0.8900000000000006 , 0.9000000000000006 , 0.9100000000000006 ,   0.9200000000000006 , 0.9300000000000006 , 0.9400000000000006 ,   0.9500000000000006 , 0.9600000000000006 , 0.9700000000000006 ,   0.9800000000000006 , 0.9900000000000007 \right] \right)=\left[ 1 ,   1.014947000636657 , 1.029776013705529 , 1.044469087191079 ,   1.059008331806833 , 1.073375947255439 , 1.087554248364218 ,   1.101525691055367 , 1.11527289811021 , 1.128778684687222 ,   1.142026083553954 , 1.154998369993414 , 1.16767908634602 ,   1.180052066148761 , 1.192101457833886 , 1.203811747950136 ,   1.215167783870255 , 1.226154795949382 , 1.236758419099762 ,   1.246964713748154 , 1.256760186143285 , 1.266131807981756 ,   1.275067035321848 , 1.283553826755846 , 1.29158066081265 ,   1.29913655256367 , 1.306211069406282 , 1.312794346000405 ,   1.318877098335118 , 1.324450636903608 , 1.329506878966172 ,   1.334038359882425 , 1.338038243495345 , 1.341500331551311 ,   1.344419072141793 , 1.346789567153917 , 1.348607578718725 ,   1.349869534647481 , 1.350572532848044 , 1.350714344714907 ,   1.350293417488142 , 1.349308875578123 , 1.347760520854542 ,   1.345648831899879 , 1.342974962229111 , 1.339740737479097 ,   1.335948651572729 , 1.331601861864506 , 1.326704183275865 ,   1.321260081430156 , 1.315274664798767 , 1.308753675871437 ,   1.301703481365363 , 1.294131061489226 , 1.286043998279732 ,   1.277450463029762 , 1.268359202828647 , 1.25877952623647 ,   1.248721288115691 , 1.238194873644713 , 1.227211181539273 ,   1.215781606508839 , 1.203918020976346 , 1.191632756090801 ,   1.17893858206338 , 1.165848687858719 , 1.152376660274093 ,   1.138536462440146 , 1.124342411777761 , 1.10980915744646 ,   1.094951657320579 , 1.079785154530145 , 1.064325153604093 ,   1.04858739625406 , 1.032587836837555 , 1.0163426175398 ,   0.999868043313951 , 0.9831805566197906 , 0.9662967120012925 ,   0.9492331505436565 , 0.932006574250646 , 0.9146337203831 ,   0.897131335799599 , 0.8795161513401855 , 0.8618048562939812 ,   0.8440140729913906 , 0.8261603315613344 , 0.8082600448937051 ,   0.7903294838468643 , 0.7723847527396025 , 0.754441765166499 ,   0.7365162201750889 , 0.7186235788426429 , 0.7007790412897039 ,   0.6829975241668103 , 0.6652936386500562 , 0.6476816689803099 ,   0.6301755515800127 , 0.6127888547805567 , 0.595534759192214 \right] 
\]
\end{eulerformula}
\begin{eulerprompt}
>function fy(t) &= r(t)*sin(t); $'fy(t)=fy(t)
\end{eulerprompt}
\begin{eulerformula}
\[
{\it fy}\left(\left[ 0 , 0.01 , 0.02 , 0.03 , 0.04 , 0.05 , 0.06 ,   0.07 , 0.08 , 0.09 , 0.1 , 0.11 , 0.12 , 0.13 , 0.14 , 0.15 , 0.16   , 0.17 , 0.18 , 0.19 , 0.2 , 0.21 , 0.2200000000000001 ,   0.2300000000000001 , 0.2400000000000001 , 0.2500000000000001 ,   0.2600000000000001 , 0.2700000000000001 , 0.2800000000000001 ,   0.2900000000000001 , 0.3000000000000001 , 0.3100000000000001 ,   0.3200000000000001 , 0.3300000000000001 , 0.3400000000000001 ,   0.3500000000000001 , 0.3600000000000002 , 0.3700000000000002 ,   0.3800000000000002 , 0.3900000000000002 , 0.4000000000000002 ,   0.4100000000000002 , 0.4200000000000002 , 0.4300000000000002 ,   0.4400000000000002 , 0.4500000000000002 , 0.4600000000000002 ,   0.4700000000000003 , 0.4800000000000003 , 0.4900000000000003 ,   0.5000000000000002 , 0.5100000000000002 , 0.5200000000000002 ,   0.5300000000000002 , 0.5400000000000003 , 0.5500000000000003 ,   0.5600000000000003 , 0.5700000000000003 , 0.5800000000000003 ,   0.5900000000000003 , 0.6000000000000003 , 0.6100000000000003 ,   0.6200000000000003 , 0.6300000000000003 , 0.6400000000000003 ,   0.6500000000000004 , 0.6600000000000004 , 0.6700000000000004 ,   0.6800000000000004 , 0.6900000000000004 , 0.7000000000000004 ,   0.7100000000000004 , 0.7200000000000004 , 0.7300000000000004 ,   0.7400000000000004 , 0.7500000000000004 , 0.7600000000000005 ,   0.7700000000000005 , 0.7800000000000005 , 0.7900000000000005 ,   0.8000000000000005 , 0.8100000000000005 , 0.8200000000000005 ,   0.8300000000000005 , 0.8400000000000005 , 0.8500000000000005 ,   0.8600000000000005 , 0.8700000000000006 , 0.8800000000000006 ,   0.8900000000000006 , 0.9000000000000006 , 0.9100000000000006 ,   0.9200000000000006 , 0.9300000000000006 , 0.9400000000000006 ,   0.9500000000000006 , 0.9600000000000006 , 0.9700000000000006 ,   0.9800000000000006 , 0.9900000000000007 \right] \right)=\left[ 0 ,   0.01014980833556662 , 0.02059826678292271 , 0.03134347622283015 ,   0.04238293991838228 , 0.05371356612987439 , 0.06533167172990376 ,   0.07723298681299934 , 0.08941266029246918 , 0.1018652664755576 ,   0.1145848126064173 , 0.1275647473648353 , 0.1407979703071057 ,   0.1542768422339107 , 0.1679931964685752 , 0.1819383510275811 ,   0.1961031216637831 , 0.2104778357613507 , 0.2250523470600841 ,   0.2398160511854019 , 0.2547579019589912 , 0.2698664284638497 ,   0.2851297528362152 , 0.3005356087557041 , 0.3160713606038417 ,   0.3317240232600813 , 0.3474802825033731 , 0.3633265159863522 ,   0.3792488147482899 , 0.3952330052320643 , 0.411264671769591 ,   0.4273291794993832 , 0.4434116976792021 , 0.4594972233561165 ,   0.4755706053556919 , 0.4916165685515136 , 0.5076197383757777 ,   0.5235646655312819 , 0.5394358508648145 , 0.5552177703616642 ,   0.5708949002207642 , 0.5864517419698421 , 0.6018728475798654 ,   0.6171428445380648 , 0.6322464608388652 , 0.6471685498521687 ,   0.6618941150286309 , 0.6764083344018014 , 0.6906965848473219 ,   0.704744466059751 , 0.7185378242080237 , 0.7320627752310482 ,   0.7453057277355214 , 0.7582534054586558 , 0.7708928692592016 ,   0.7832115386008901 , 0.7951972124932317 , 0.8068380898554457 ,   0.8181227892702304 , 0.8290403680950348 , 0.8395803408995157 ,   0.8497326971989371 , 0.8594879184543822 , 0.8688369943118147 ,   0.877771438053233 , 0.8862833012344233 , 0.894365187485098 ,   0.9020102654485477 , 0.9092122808393135 , 0.91596556759876 ,   0.9222650581299157 , 0.9281062925943645 , 0.9334854272555032 ,   0.9383992418539865 , 0.9428451460027243 , 0.9468211845903713 ,   0.9503260421838114 , 0.9533590464217597 , 0.9559201703932094 ,   0.9580100339960551 , 0.9596299042728891 , 0.9607816947225576 ,   0.9614679635877484 , 0.9616919111204768 , 0.9614573758289937 ,   0.9607688297112769 , 0.9596313724818526 , 0.9580507248003547 ,   0.9560332205117796 , 0.9535857979100135 , 0.950715990037748 ,   0.9474319140374602 , 0.9437422595696462 , 0.9396562763159917 ,   0.9351837605866338 , 0.9303350410521015 , 0.9251209636219332 ,   0.9195528754933222 , 0.9136426083945087 , 0.9074024610488752   \right] 
\]
\end{eulerformula}
\begin{eulerprompt}
>function ds(t) &= trigreduce(radcan(sqrt(diff(fx(t),t)^2+diff(fy(t),t)^2))); $'ds(t)=ds(t)        
\end{eulerprompt}
\begin{eulerformula}
\[
{\it ds}\left(t\right)=\sqrt{\left(\frac{d}{d\,t}\,\left(\frac{4\,  \sin t+\cos \left(2\,t\right)-\cos \left(4\,t\right)}{4}\right)(t)  \right)^2+\left(\frac{d}{d\,t}\,\left(\frac{4\,\cos t+\sin \left(2\,  t\right)+\sin \left(4\,t\right)}{4}\right)(t)\right)^2}
\]
\end{eulerformula}
\begin{eulerprompt}
>$integrate(ds(x),x,0,2*pi) //panjang (keliling) kurva
\end{eulerprompt}
\begin{eulerformula}
\[
\int_{0}^{2\,\pi}{\sqrt{\left(\frac{d}{d\,x}\,\left(\frac{4\,\sin x  +\cos \left(2\,x\right)-\cos \left(4\,x\right)}{4}\right)(x)\right)^  2+\left(\frac{d}{d\,x}\,\left(\frac{4\,\cos x+\sin \left(2\,x\right)  +\sin \left(4\,x\right)}{4}\right)(x)\right)^2}\;dx}
\]
\end{eulerformula}
\begin{eulercomment}
Maxima gagal melakukan perhitungan eksak integral tersebut.\\
==============================================================================
PENYELESAIAN LENGKAP: PANJANG BUSUR KURVA POLAR
==============================================================================
1. Bersihkan memori total dari sesi sebelumnya
\end{eulercomment}
\begin{eulerprompt}
> reset();
\end{eulerprompt}
\begin{eulercomment}
2. Definisikan fungsi r, fx, dan fy sebagai ekspresi simbolis murni (tanpa keyword function)
\end{eulercomment}
\begin{eulerprompt}
>r &= 1 + sin(3*t)/2;
>fx &= r*cos(t);
>fy &= r*sin(t);
\end{eulerprompt}
\begin{eulercomment}
3. Hitung elemen panjang busur (dse) secara simbolis menggunakan diff
\end{eulercomment}
\begin{eulerprompt}
>dse &= trigreduce(radcan(sqrt(diff(fx,t)^2 + diff(fy,t)^2)));
> $'ds(t)=dse
\end{eulerprompt}
\begin{eulerformula}
\[
{\it ds}\left(t\right)=\frac{\sqrt{9+4\,\sin \left(3\,t\right)+4\,  \cos \left(6\,t\right)}}{2}
\]
\end{eulerformula}
\begin{eulercomment}
4. Konversi ekspresi simbolis ke fungsi numerik EMT agar bisa diintegralkan
\end{eulercomment}
\begin{eulerprompt}
>function ds(t) :=sqrt(9 + 4*sin(3*t) + 4*cos(6*t)) / 2;
\end{eulerprompt}
\begin{eulercomment}
5. Hitung panjang busur (keliling kurva) secara numerik dari 0 sampai 2*pi
\end{eulercomment}
\begin{eulerprompt}
>$integrate(dse, t, 0, 2*pi)
\end{eulerprompt}
\begin{eulerformula}
\[
\frac{\int_{0}^{2\,\pi}{\sqrt{9+4\,\sin \left(3\,t\right)+4\,\cos   \left(6\,t\right)}\;dt}}{2}
\]
\end{eulerformula}
\begin{eulerprompt}
> integrate("sqrt(9 + 4*sin(3*x) + 4*cos(6*x)) / 2", 0, 2*pi)
\end{eulerprompt}
\begin{euleroutput}
  9.0749467823
\end{euleroutput}
\begin{eulerprompt}
> 
\end{eulerprompt}
\begin{eulercomment}
Berikut kita hitung integralnya secara numerik dengan perintah EMT.
\end{eulercomment}
\begin{eulerprompt}
> integrate("ds(x)",0,2*pi) 
\end{eulerprompt}
\begin{euleroutput}
  9.0749467823
\end{euleroutput}
\begin{eulercomment}
Spiral Logaritmik

\end{eulercomment}
\begin{eulerformula}
\[
x=e^{ax}\cos x,\ y=e^{ax}\sin x.
\]
\end{eulerformula}
\begin{eulerprompt}
>a=0.1; plot2d("exp(a*x)*cos(x)","exp(a*x)*sin(x)",r=2,xmin=0,xmax=2*pi):
\end{eulerprompt}
\eulerimg{27}{Image Use/TugasAplikom6Adi25030630001-165.png}
\begin{eulerprompt}
>&kill(a) // hapus expresi a
\end{eulerprompt}
\begin{euleroutput}
  
                                   done
  
\end{euleroutput}
\begin{eulerprompt}
>function fx(t) &= exp(a*t)*cos(t); $'fx(t)=fx(t)
\end{eulerprompt}
\begin{eulerformula}
\[
{\it fx}\left(\left[ 0 , 0.01 , 0.02 , 0.03 , 0.04 , 0.05 , 0.06 ,   0.07 , 0.08 , 0.09 , 0.1 , 0.11 , 0.12 , 0.13 , 0.14 , 0.15 , 0.16   , 0.17 , 0.18 , 0.19 , 0.2 , 0.21 , 0.2200000000000001 ,   0.2300000000000001 , 0.2400000000000001 , 0.2500000000000001 ,   0.2600000000000001 , 0.2700000000000001 , 0.2800000000000001 ,   0.2900000000000001 , 0.3000000000000001 , 0.3100000000000001 ,   0.3200000000000001 , 0.3300000000000001 , 0.3400000000000001 ,   0.3500000000000001 , 0.3600000000000002 , 0.3700000000000002 ,   0.3800000000000002 , 0.3900000000000002 , 0.4000000000000002 ,   0.4100000000000002 , 0.4200000000000002 , 0.4300000000000002 ,   0.4400000000000002 , 0.4500000000000002 , 0.4600000000000002 ,   0.4700000000000003 , 0.4800000000000003 , 0.4900000000000003 ,   0.5000000000000002 , 0.5100000000000002 , 0.5200000000000002 ,   0.5300000000000002 , 0.5400000000000003 , 0.5500000000000003 ,   0.5600000000000003 , 0.5700000000000003 , 0.5800000000000003 ,   0.5900000000000003 , 0.6000000000000003 , 0.6100000000000003 ,   0.6200000000000003 , 0.6300000000000003 , 0.6400000000000003 ,   0.6500000000000004 , 0.6600000000000004 , 0.6700000000000004 ,   0.6800000000000004 , 0.6900000000000004 , 0.7000000000000004 ,   0.7100000000000004 , 0.7200000000000004 , 0.7300000000000004 ,   0.7400000000000004 , 0.7500000000000004 , 0.7600000000000005 ,   0.7700000000000005 , 0.7800000000000005 , 0.7900000000000005 ,   0.8000000000000005 , 0.8100000000000005 , 0.8200000000000005 ,   0.8300000000000005 , 0.8400000000000005 , 0.8500000000000005 ,   0.8600000000000005 , 0.8700000000000006 , 0.8800000000000006 ,   0.8900000000000006 , 0.9000000000000006 , 0.9100000000000006 ,   0.9200000000000006 , 0.9300000000000006 , 0.9400000000000006 ,   0.9500000000000006 , 0.9600000000000006 , 0.9700000000000006 ,   0.9800000000000006 , 0.9900000000000007 \right] \right)=\left[ 1 ,   0.9999500004166653\,e^{0.01\,a} , 0.9998000066665778\,e^{0.02\,a} ,   0.9995500337489875\,e^{0.03\,a} , 0.9992001066609779\,e^{0.04\,a} ,   0.9987502603949663\,e^{0.05\,a} , 0.9982005399352042\,e^{0.06\,a} ,   0.9975510002532796\,e^{0.07\,a} , 0.9968017063026194\,e^{0.08\,a} ,   0.9959527330119943\,e^{0.09\,a} , 0.9950041652780258\,e^{0.1\,a} ,   0.9939560979566968\,e^{0.11\,a} , 0.9928086358538663\,e^{0.12\,a} ,   0.9915618937147881\,e^{0.13\,a} , 0.9902159962126372\,e^{0.14\,a} ,   0.9887710779360422\,e^{0.15\,a} , 0.9872272833756269\,e^{0.16\,a} ,   0.9855847669095608\,e^{0.17\,a} , 0.9838436927881214\,e^{0.18\,a} ,   0.9820042351172703\,e^{0.19\,a} , 0.9800665778412416\,e^{0.2\,a} ,   0.9780309147241483\,e^{0.21\,a} , 0.9758974493306055\,e^{  0.2200000000000001\,a} , 0.9736663950053748\,e^{0.2300000000000001\,  a} , 0.9713379748520296\,e^{0.2400000000000001\,a} ,   0.9689124217106447\,e^{0.2500000000000001\,a} , 0.9663899781345132\,  e^{0.2600000000000001\,a} , 0.9637708963658905\,e^{  0.2700000000000001\,a} , 0.9610554383107709\,e^{0.2800000000000001\,  a} , 0.9582438755126972\,e^{0.2900000000000001\,a} ,   0.955336489125606\,e^{0.3000000000000001\,a} , 0.9523335698857134\,e  ^{0.3100000000000001\,a} , 0.9492354180824408\,e^{0.3200000000000001  \,a} , 0.9460423435283869\,e^{0.3300000000000001\,a} ,   0.9427546655283462\,e^{0.3400000000000001\,a} , 0.9393727128473789\,  e^{0.3500000000000001\,a} , 0.9358968236779348\,e^{  0.3600000000000002\,a} , 0.9323273456060344\,e^{0.3700000000000002\,  a} , 0.9286646355765101\,e^{0.3800000000000002\,a} ,   0.924909059857313\,e^{0.3900000000000002\,a} , 0.921060994002885\,e  ^{0.4000000000000002\,a} , 0.917120822816605\,e^{0.4100000000000002  \,a} , 0.9130889403123081\,e^{0.4200000000000002\,a} ,   0.9089657496748851\,e^{0.4300000000000002\,a} , 0.9047516632199634\,  e^{0.4400000000000002\,a} , 0.9004471023526768\,e^{  0.4500000000000002\,a} , 0.8960524975255252\,e^{0.4600000000000002\,  a} , 0.8915682881953289\,e^{0.4700000000000003\,a} ,   0.886994922779284\,e^{0.4800000000000003\,a} , 0.8823328586101213\,e  ^{0.4900000000000003\,a} , 0.8775825618903726\,e^{0.5000000000000002  \,a} , 0.8727445076457512\,e^{0.5100000000000002\,a} ,   0.8678191796776498\,e^{0.5200000000000002\,a} , 0.8628070705147609\,  e^{0.5300000000000002\,a} , 0.857708681363824\,e^{0.5400000000000003  \,a} , 0.8525245220595056\,e^{0.5500000000000003\,a} ,   0.847255111013416\,e^{0.5600000000000003\,a} , 0.8419009751622686\,e  ^{0.5700000000000003\,a} , 0.8364626499151868\,e^{0.5800000000000003  \,a} , 0.8309406791001633\,e^{0.5900000000000003\,a} ,   0.8253356149096781\,e^{0.6000000000000003\,a} , 0.8196480178454794\,  e^{0.6100000000000003\,a} , 0.8138784566625338\,e^{  0.6200000000000003\,a} , 0.8080275083121516\,e^{0.6300000000000003\,  a} , 0.8020957578842924\,e^{0.6400000000000003\,a} ,   0.7960837985490556\,e^{0.6500000000000004\,a} , 0.7899922314973649\,  e^{0.6600000000000004\,a} , 0.783821665880849\,e^{0.6700000000000004  \,a} , 0.7775727187509277\,e^{0.6800000000000004\,a} ,   0.7712460149971063\,e^{0.6900000000000004\,a} , 0.7648421872844882\,  e^{0.7000000000000004\,a} , 0.7583618759905079\,e^{  0.7100000000000004\,a} , 0.7518057291408947\,e^{0.7200000000000004\,  a} , 0.7451744023448701\,e^{0.7300000000000004\,a} ,   0.7384685587295876\,e^{0.7400000000000004\,a} , 0.7316888688738206\,  e^{0.7500000000000004\,a} , 0.7248360107409049\,e^{  0.7600000000000005\,a} , 0.7179106696109431\,e^{0.7700000000000005\,  a} , 0.7109135380122771\,e^{0.7800000000000005\,a} ,   0.7038453156522357\,e^{0.7900000000000005\,a} , 0.696706709347165\,e  ^{0.8000000000000005\,a} , 0.6894984329517466\,e^{0.8100000000000005  \,a} , 0.6822212072876132\,e^{0.8200000000000005\,a} ,   0.6748757600712667\,e^{0.8300000000000005\,a} , 0.6674628258413078\,  e^{0.8400000000000005\,a} , 0.6599831458849817\,e^{  0.8500000000000005\,a} , 0.6524374681640515\,e^{0.8600000000000005\,  a} , 0.6448265472400008\,e^{0.8700000000000006\,a} ,   0.6371511441985798\,e^{0.8800000000000006\,a} , 0.6294120265736964\,  e^{0.8900000000000006\,a} , 0.6216099682706641\,e^{  0.9000000000000006\,a} , 0.6137457494888111\,e^{0.9100000000000006\,  a} , 0.6058201566434623\,e^{0.9200000000000006\,a} ,   0.5978339822872978\,e^{0.9300000000000006\,a} , 0.5897880250310977\,  e^{0.9400000000000006\,a} , 0.581683089463883\,e^{0.9500000000000006  \,a} , 0.5735199860724561\,e^{0.9600000000000006\,a} ,   0.5652995311603538\,e^{0.9700000000000006\,a} , 0.5570225467662168\,  e^{0.9800000000000006\,a} , 0.548689860581587\,e^{0.9900000000000007  \,a} \right] 
\]
\end{eulerformula}
\begin{eulerprompt}
>function fy(t) &= exp(a*t)*sin(t); $'fy(t)=fy(t)
\end{eulerprompt}
\begin{eulerformula}
\[
{\it fy}\left(\left[ 0 , 0.01 , 0.02 , 0.03 , 0.04 , 0.05 , 0.06 ,   0.07 , 0.08 , 0.09 , 0.1 , 0.11 , 0.12 , 0.13 , 0.14 , 0.15 , 0.16   , 0.17 , 0.18 , 0.19 , 0.2 , 0.21 , 0.2200000000000001 ,   0.2300000000000001 , 0.2400000000000001 , 0.2500000000000001 ,   0.2600000000000001 , 0.2700000000000001 , 0.2800000000000001 ,   0.2900000000000001 , 0.3000000000000001 , 0.3100000000000001 ,   0.3200000000000001 , 0.3300000000000001 , 0.3400000000000001 ,   0.3500000000000001 , 0.3600000000000002 , 0.3700000000000002 ,   0.3800000000000002 , 0.3900000000000002 , 0.4000000000000002 ,   0.4100000000000002 , 0.4200000000000002 , 0.4300000000000002 ,   0.4400000000000002 , 0.4500000000000002 , 0.4600000000000002 ,   0.4700000000000003 , 0.4800000000000003 , 0.4900000000000003 ,   0.5000000000000002 , 0.5100000000000002 , 0.5200000000000002 ,   0.5300000000000002 , 0.5400000000000003 , 0.5500000000000003 ,   0.5600000000000003 , 0.5700000000000003 , 0.5800000000000003 ,   0.5900000000000003 , 0.6000000000000003 , 0.6100000000000003 ,   0.6200000000000003 , 0.6300000000000003 , 0.6400000000000003 ,   0.6500000000000004 , 0.6600000000000004 , 0.6700000000000004 ,   0.6800000000000004 , 0.6900000000000004 , 0.7000000000000004 ,   0.7100000000000004 , 0.7200000000000004 , 0.7300000000000004 ,   0.7400000000000004 , 0.7500000000000004 , 0.7600000000000005 ,   0.7700000000000005 , 0.7800000000000005 , 0.7900000000000005 ,   0.8000000000000005 , 0.8100000000000005 , 0.8200000000000005 ,   0.8300000000000005 , 0.8400000000000005 , 0.8500000000000005 ,   0.8600000000000005 , 0.8700000000000006 , 0.8800000000000006 ,   0.8900000000000006 , 0.9000000000000006 , 0.9100000000000006 ,   0.9200000000000006 , 0.9300000000000006 , 0.9400000000000006 ,   0.9500000000000006 , 0.9600000000000006 , 0.9700000000000006 ,   0.9800000000000006 , 0.9900000000000007 \right] \right)=\left[ 0 ,   0.009999833334166664\,e^{0.01\,a} , 0.01999866669333308\,e^{0.02\,a}   , 0.02999550020249566\,e^{0.03\,a} , 0.03998933418663416\,e^{0.04\,  a} , 0.04997916927067833\,e^{0.05\,a} , 0.0599640064794446\,e^{0.06  \,a} , 0.06994284733753277\,e^{0.07\,a} , 0.0799146939691727\,e^{  0.08\,a} , 0.08987854919801104\,e^{0.09\,a} , 0.09983341664682814\,e  ^{0.1\,a} , 0.1097783008371748\,e^{0.11\,a} , 0.1197122072889193\,e  ^{0.12\,a} , 0.1296341426196948\,e^{0.13\,a} , 0.1395431146442365\,e  ^{0.14\,a} , 0.1494381324735992\,e^{0.15\,a} , 0.159318206614246\,e  ^{0.16\,a} , 0.169182349066996\,e^{0.17\,a} , 0.1790295734258242\,e  ^{0.18\,a} , 0.1888588949765006\,e^{0.19\,a} , 0.1986693307950612\,e  ^{0.2\,a} , 0.2084598998460996\,e^{0.21\,a} , 0.2182296230808694\,e  ^{0.2200000000000001\,a} , 0.2279775235351885\,e^{0.2300000000000001  \,a} , 0.2377026264271347\,e^{0.2400000000000001\,a} ,   0.247403959254523\,e^{0.2500000000000001\,a} , 0.2570805518921552\,e  ^{0.2600000000000001\,a} , 0.2667314366888312\,e^{0.2700000000000001  \,a} , 0.2763556485641138\,e^{0.2800000000000001\,a} ,   0.2859522251048356\,e^{0.2900000000000001\,a} , 0.2955202066613397\,  e^{0.3000000000000001\,a} , 0.3050586364434436\,e^{  0.3100000000000001\,a} , 0.3145665606161179\,e^{0.3200000000000001\,  a} , 0.3240430283948685\,e^{0.3300000000000001\,a} ,   0.3334870921408145\,e^{0.3400000000000001\,a} , 0.3428978074554515\,  e^{0.3500000000000001\,a} , 0.3522742332750901\,e^{  0.3600000000000002\,a} , 0.3616154319649622\,e^{0.3700000000000002\,  a} , 0.3709204694129828\,e^{0.3800000000000002\,a} ,   0.3801884151231616\,e^{0.3900000000000002\,a} , 0.3894183423086507\,  e^{0.4000000000000002\,a} , 0.3986093279844231\,e^{  0.4100000000000002\,a} , 0.4077604530595704\,e^{0.4200000000000002\,  a} , 0.416870802429211\,e^{0.4300000000000002\,a} ,   0.4259394650659998\,e^{0.4400000000000002\,a} , 0.4349655341112304\,  e^{0.4500000000000002\,a} , 0.44394810696552\,e^{0.4600000000000002  \,a} , 0.4528862853790685\,e^{0.4700000000000003\,a} ,   0.4617791755414831\,e^{0.4800000000000003\,a} , 0.4706258881711582\,  e^{0.4900000000000003\,a} , 0.4794255386042032\,e^{  0.5000000000000002\,a} , 0.4881772468829077\,e^{0.5100000000000002\,  a} , 0.4968801378437369\,e^{0.5200000000000002\,a} ,   0.5055333412048472\,e^{0.5300000000000002\,a} , 0.5141359916531133\,  e^{0.5400000000000003\,a} , 0.5226872289306594\,e^{  0.5500000000000003\,a} , 0.5311861979208836\,e^{0.5600000000000003\,  a} , 0.5396320487339695\,e^{0.5700000000000003\,a} ,   0.5480239367918738\,e^{0.5800000000000003\,a} , 0.556361022912784\,e  ^{0.5900000000000003\,a} , 0.5646424733950356\,e^{0.6000000000000003  \,a} , 0.5728674601004815\,e^{0.6100000000000003\,a} ,   0.5810351605373053\,e^{0.6200000000000003\,a} , 0.5891447579422698\,  e^{0.6300000000000003\,a} , 0.5971954413623923\,e^{  0.6400000000000003\,a} , 0.6051864057360399\,e^{0.6500000000000004\,  a} , 0.6131168519734341\,e^{0.6600000000000004\,a} ,   0.6209859870365599\,e^{0.6700000000000004\,a} , 0.6287930240184688\,  e^{0.6800000000000004\,a} , 0.6365371822219682\,e^{  0.6900000000000004\,a} , 0.6442176872376913\,e^{0.7000000000000004\,  a} , 0.651833771021537\,e^{0.7100000000000004\,a} ,   0.6593846719714734\,e^{0.7200000000000004\,a} , 0.6668696350036982\,  e^{0.7300000000000004\,a} , 0.6742879116281454\,e^{  0.7400000000000004\,a} , 0.6816387600233345\,e^{0.7500000000000004\,  a} , 0.6889214451105516\,e^{0.7600000000000005\,a} ,   0.696135238627357\,e^{0.7700000000000005\,a} , 0.7032794192004105\,e  ^{0.7800000000000005\,a} , 0.7103532724176082\,e^{0.7900000000000005  \,a} , 0.7173560908995231\,e^{0.8000000000000005\,a} ,   0.7242871743701429\,e^{0.8100000000000005\,a} , 0.7311458297268962\,  e^{0.8200000000000005\,a} , 0.7379313711099631\,e^{  0.8300000000000005\,a} , 0.7446431199708596\,e^{0.8400000000000005\,  a} , 0.751280405140293\,e^{0.8500000000000005\,a} ,   0.7578425628952773\,e^{0.8600000000000005\,a} , 0.7643289370255054\,  e^{0.8700000000000006\,a} , 0.7707388788989696\,e^{  0.8800000000000006\,a} , 0.7770717475268242\,e^{0.8900000000000006\,  a} , 0.7833269096274837\,e^{0.9000000000000006\,a} ,   0.7895037396899508\,e^{0.9100000000000006\,a} , 0.7956016200363664\,  e^{0.9200000000000006\,a} , 0.8016199408837775\,e^{  0.9300000000000006\,a} , 0.8075581004051147\,e^{0.9400000000000006\,  a} , 0.8134155047893741\,e^{0.9500000000000006\,a} ,   0.8191915683009986\,e^{0.9600000000000006\,a} , 0.8248857133384504\,  e^{0.9700000000000006\,a} , 0.8304973704919708\,e^{  0.9800000000000006\,a} , 0.8360259786005209\,e^{0.9900000000000007\,  a} \right] 
\]
\end{eulerformula}
\begin{eulerprompt}
>function df(t) &= trigreduce(radcan(sqrt(diff(fx(t),t)^2+diff(fy(t),t)^2))); $'df(t)=df(t)
\end{eulerprompt}
\begin{euleroutput}
  Maxima said:
  diff: second argument must be a variable; found errexp1
   -- an error. To debug this try: debugmode(true);
  
  Error in:
  ... e(radcan(sqrt(diff(fx(t),t)^2+diff(fy(t),t)^2))); $'df(t)=df(t ...
                                                       ^
\end{euleroutput}
\begin{eulerprompt}
>S &=integrate(df(t),t,0,2*%pi); $S // panjang kurva (spiral)
\end{eulerprompt}
\begin{euleroutput}
  Maxima said:
  defint: variable of integration cannot be a constant; found errexp1
   -- an error. To debug this try: debugmode(true);
  
  Error in:
  S &=integrate(df(t),t,0,2*%pi); $S // panjang kurva (spiral) ...
                                ^
\end{euleroutput}
\begin{eulerprompt}
>S(a=0.1) // Panjang kurva untuk a=0.1
\end{eulerprompt}
\begin{euleroutput}
  Function S not found.
  Try list ... to find functions!
  Error in:
  S(a=0.1) // Panjang kurva untuk a=0.1 ...
          ^
\end{euleroutput}
\begin{eulercomment}
Soal:

Tunjukkan bahwa keliling lingkaran dengan jari-jari r adalah K=2.pi.r.

==============================================================================
BUKTI KELILING LINGKARAN (K = 2 * pi * r) DI EMT / MAXIMA
==============================================================================
1. Bersihkan memori total
\end{eulercomment}
\begin{eulerprompt}
>reset();
\end{eulerprompt}
\begin{eulercomment}
2. Definisikan variabel simbolis t dan konstanta r
\end{eulercomment}
\begin{eulerprompt}
>t &= 't; r &= 'r;
\end{eulerprompt}
\begin{eulercomment}
3. Dalam koordinat polar, jari-jari lingkaran r(t) konstan bernilai r
\end{eulercomment}
\begin{eulerprompt}
>rtek &= r;
\end{eulerprompt}
\begin{eulercomment}
4. Hitung elemen panjang busur (ds) secara simbolis
\end{eulercomment}
\begin{eulerprompt}
>dspolar &= trigreduce(radcan(sqrt(rtek^2 + diff(rtek, t)^2)));
>$'ds = dspolar
\end{eulerprompt}
\begin{eulerformula}
\[
{\it ds}=\left| r\right| 
\]
\end{eulerformula}
\begin{eulercomment}
5. Integralkan dari 0 sampai 2*pi untuk mendapatkan keliling total
\end{eulercomment}
\begin{eulerprompt}
>Keliling &= integrate(dspolar, t, 0, 2*pi);
>$'K = Keliling
\end{eulerprompt}
\begin{eulerformula}
\[
K=2\,\pi\,\left| r\right| 
\]
\end{eulerformula}
\begin{eulerprompt}
>         
\end{eulerprompt}
\begin{eulercomment}
Berikut adalah contoh menghitung panjang parabola.
\end{eulercomment}
\begin{eulerprompt}
>plot2d("x^2",xmin=-1,xmax=1):
\end{eulerprompt}
\eulerimg{27}{Image Use/TugasAplikom6Adi25030630001-170.png}
\begin{eulerprompt}
>$showev('integrate(sqrt(1+diff(x^2,x)^2),x,-1,1))
\end{eulerprompt}
\begin{eulerformula}
\[
\int_{-1}^{1}{\sqrt{4\,x^2+1}\;dx}=\frac{{\rm asinh}\; 2+2\,\sqrt{5  }}{2}
\]
\end{eulerformula}
\begin{eulerprompt}
>$float(%)
\end{eulerprompt}
\begin{eulerformula}
\[
\int_{-1.0}^{1.0}{\sqrt{4.0\,x^2+1.0}\;dx}=2.957885715089195
\]
\end{eulerformula}
\begin{eulerprompt}
>x=-1:0.2:1; y=x^2; plot2d(x,y);  ...
>  plot2d(x,y,points=1,style="o#",add=1):
\end{eulerprompt}
\eulerimg{27}{Image Use/TugasAplikom6Adi25030630001-173.png}
\begin{eulerprompt}
> 
\end{eulerprompt}
\begin{eulercomment}
Panjang tersebut dapat dihampiri dengan menggunakan jumlah panjang ruas-ruas garis yang menghubungkan titik-titik pada parabola
tersebut.
\end{eulercomment}
\begin{eulerprompt}
>i=1:cols(x)-1; sum(sqrt((x[i+1]-x[i])^2+(y[i+1]-y[i])^2))
\end{eulerprompt}
\begin{euleroutput}
  2.95191957027
\end{euleroutput}
\begin{eulercomment}
Hasilnya mendekati panjang yang dihitung secara eksak. Untuk
mendapatkan hampiran yang cukup akurat, jarak antar titik dapat
diperkecil, misalnya 0.1, 0.05, 0.01, dan seterusnya. Cobalah Anda
ulangi perhitungannya dengan nilai-nilai tersebut.

==============================================================================
PENGHAMPIRAN PANJANG KURVA PARABOLA y = x\textasciicircum{}2 DENGAN BERBAGAI
SUB-INTERVAL (h)
==============================================================================
1. Reset memori agar bersih dari variabel sebelumnya
\end{eulercomment}
\begin{eulerprompt}
>reset();
\end{eulerprompt}
\begin{eulercomment}
2. Eksak (Sebagai pembanding): Integral eksak panjang kurva y = x\textasciicircum{}2 dari -1 sampai 1
\end{eulercomment}
\begin{eulerprompt}
>integrate("sqrt(1 + 4*x^2)", -1, 1)
\end{eulerprompt}
\begin{euleroutput}
  2.95788571509
\end{euleroutput}
\eulersubheading{}
\begin{eulercomment}
A. Ulangi dengan h = 0.1
\end{eulercomment}
\begin{eulerprompt}
>x = -1:0.1:1; y = x^2;
>i = 1:cols(x)-1; sum(sqrt((x[i+1]-x[i])^2 + (y[i+1]-y[i])^2))
\end{eulerprompt}
\begin{euleroutput}
  2.95639479497
\end{euleroutput}
\eulersubheading{}
\begin{eulercomment}
B. Ulangi dengan h = 0.05
\end{eulercomment}
\begin{eulerprompt}
>x = -1:0.05:1; y = x^2;
>i = 1:cols(x)-1; sum(sqrt((x[i+1]-x[i])^2 + (y[i+1]-y[i])^2))
\end{eulerprompt}
\begin{euleroutput}
  2.95751302406
\end{euleroutput}
\eulersubheading{}
\begin{eulercomment}
C. Ulangi dengan h = 0.01
\end{eulercomment}
\begin{eulerprompt}
>x = -1:0.01:1; y = x^2;
>i = 1:cols(x)-1; sum(sqrt((x[i+1]-x[i])^2 + (y[i+1]-y[i])^2))
\end{eulerprompt}
\begin{euleroutput}
  2.95787080795
\end{euleroutput}
\eulersubheading{}
\begin{eulercomment}
D. Ulangi dengan h = 0.001 (Menghasilkan hampiran yang sangat akurat)
\end{eulercomment}
\begin{eulerprompt}
>x = -1:0.001:1; y = x^2;
>i = 1:cols(x)-1; sum(sqrt((x[i+1]-x[i])^2 + (y[i+1]-y[i])^2))
\end{eulerprompt}
\begin{euleroutput}
  2.95788556602
\end{euleroutput}
\eulersubheading{Koordinat Kartesius}
\begin{eulercomment}
Berikut diberikan contoh perhitungan panjang kurva menggunakan
koordinat Kartesius. Kita akan hitung panjang kurva dengan persamaan
implisit:

\end{eulercomment}
\begin{eulerformula}
\[
x^3+y^3-3xy=0.
\]
\end{eulerformula}
\begin{eulerprompt}
>z &= x^3+y^3-3*x*y; $z
\end{eulerprompt}
\begin{eulerformula}
\[
y^3-3\,x\,y+x^3
\]
\end{eulerformula}
\begin{eulerprompt}
>plot2d(z,r=2,level=0,n=100):
\end{eulerprompt}
\eulerimg{27}{Image Use/TugasAplikom6Adi25030630001-176.png}
\begin{eulercomment}
Kita tertarik pada kurva di kuadran pertama.
\end{eulercomment}
\begin{eulerprompt}
>plot2d(z,a=0,b=2,c=0,d=2,level=[-10;0],n=100,contourwidth=3,style="/"):
\end{eulerprompt}
\eulerimg{27}{Image Use/TugasAplikom6Adi25030630001-177.png}
\begin{eulercomment}
Kita selesaikan persamaannya untuk x.
\end{eulercomment}
\begin{eulerprompt}
>$z with y=l*x, sol &= solve(%,x); $sol
\end{eulerprompt}
\begin{eulerformula}
\[
\left[ x=\frac{3\,l}{l^3+1} , x=0 \right] 
\]
\end{eulerformula}
\eulerimg{1}{Image Use/TugasAplikom6Adi25030630001-179-large.png}
\begin{eulercomment}
Kita gunakan solusi tersebut untuk mendefinisikan fungsi dengan Maxima.
\end{eulercomment}
\begin{eulerprompt}
>function f(l) &= rhs(sol[1]); $'f(l)=f(l)
\end{eulerprompt}
\begin{eulerformula}
\[
f\left(l\right)=\frac{3\,l}{l^3+1}
\]
\end{eulerformula}
\begin{eulercomment}
Fungsi tersebut juga dapat digunaka untuk menggambar kurvanya. Ingat, bahwa fungsi tersebut adalah nilai x dan nilai y=l*x, yakni
x=f(l) dan y=l*f(l).
\end{eulercomment}
\begin{eulerprompt}
>plot2d(&f(x),&x*f(x),xmin=-0.5,xmax=2,a=0,b=2,c=0,d=2,r=1.5):
\end{eulerprompt}
\eulerimg{27}{Image Use/TugasAplikom6Adi25030630001-181.png}
\begin{eulercomment}
Elemen panjang kurva adalah:

\end{eulercomment}
\begin{eulerformula}
\[
ds=\sqrt{f'(l)^2+(lf'(l)+f(l))^2}.
\]
\end{eulerformula}
\begin{eulerprompt}
>function ds(l) &= ratsimp(sqrt(diff(f(l),l)^2+diff(l*f(l),l)^2)); $'ds(l)=ds(l)
\end{eulerprompt}
\begin{eulerformula}
\[
{\it ds}\left(l\right)=\frac{\sqrt{9\,l^8+36\,l^6-36\,l^5-36\,l^3+  36\,l^2+9}}{\sqrt{l^{12}+4\,l^9+6\,l^6+4\,l^3+1}}
\]
\end{eulerformula}
\begin{eulerprompt}
>$integrate(ds(l),l,0,1)
\end{eulerprompt}
\begin{eulerformula}
\[
\int_{0}^{1}{\frac{\sqrt{9\,l^8+36\,l^6-36\,l^5-36\,l^3+36\,l^2+9}  }{\sqrt{l^{12}+4\,l^9+6\,l^6+4\,l^3+1}}\;dl}
\]
\end{eulerformula}
\begin{eulercomment}
Integral tersebut tidak dapat dihitung secara eksak menggunakan Maxima. Kita hitung integral etrsebut secara numerik dengan Euler.
Karena kurva simetris, kita hitung untuk nilai variabel integrasi dari 0 sampai 1, kemudian hasilnya dikalikan 2. 
\end{eulercomment}
\begin{eulerprompt}
>2*integrate("ds(x)",0,1)
\end{eulerprompt}
\begin{euleroutput}
  4.91748872168
\end{euleroutput}
\begin{eulerprompt}
>2*romberg(&ds(x),0,1)// perintah Euler lain untuk menghitung nilai hampiran integral yang sama
\end{eulerprompt}
\begin{euleroutput}
  4.91748872168
\end{euleroutput}
\begin{eulercomment}
Perhitungan di datas dapat dilakukan untuk sebarang fungsi x dan y dengan mendefinisikan fungsi EMT, misalnya kita beri nama
panjangkurva. Fungsi ini selalu memanggil Maxima untuk menurunkan fungsi yang diberikan.
\end{eulercomment}
\begin{eulerprompt}
>function panjangkurva(fx,fy,a,b) ...
\end{eulerprompt}
\begin{eulerudf}
  ds=mxm("sqrt(diff(@fx,x)^2+diff(@fy,x)^2)");
  return romberg(ds,a,b);
  endfunction
\end{eulerudf}
\begin{eulerprompt}
>panjangkurva("x","x^2",-1,1) // cek untuk menghitung panjang kurva parabola sebelumnya
\end{eulerprompt}
\begin{euleroutput}
  2.95788571509
\end{euleroutput}
\begin{eulercomment}
Bandingkan dengan nilai eksak di atas.
\end{eulercomment}
\begin{eulerprompt}
>2*panjangkurva(mxm("f(x)"),mxm("x*f(x)"),0,1) // cek contoh terakhir, bandingkan hasilnya!
\end{eulerprompt}
\begin{euleroutput}
  4.91748872168
\end{euleroutput}
\begin{eulercomment}
[TERLIHAT MEMILIKI HASIL YANG SAMA]

Kita hitung panjang spiral Archimides berikut ini dengan fungsi
tersebut.
\end{eulercomment}
\begin{eulerprompt}
>plot2d("x*cos(x)","x*sin(x)",xmin=0,xmax=2*pi,square=1):
\end{eulerprompt}
\eulerimg{27}{Image Use/TugasAplikom6Adi25030630001-185.png}
\begin{eulerprompt}
>panjangkurva("x*cos(x)","x*sin(x)",0,2*pi)
\end{eulerprompt}
\begin{euleroutput}
  21.2562941482
\end{euleroutput}
\begin{eulercomment}
Berikut kita definisikan fungsi yang sama namun dengan Maxima, untuk perhitungan eksak.
\end{eulercomment}
\begin{eulerprompt}
>&kill(ds,x,fx,fy)
\end{eulerprompt}
\begin{euleroutput}
  
                                   done
  
\end{euleroutput}
\begin{eulerprompt}
>function ds(fx,fy) &&= sqrt(diff(fx,x)^2+diff(fy,x)^2)
\end{eulerprompt}
\begin{euleroutput}
  
                             2              2
                    sqrt(diff (fy, x) + diff (fx, x))
  
\end{euleroutput}
\begin{eulerprompt}
>sol &= ds(x*cos(x),x*sin(x)); $sol // Kita gunakan untuk menghitung panjang kurva terakhir di atas
\end{eulerprompt}
\begin{eulerformula}
\[
\sqrt{\left(\cos x-x\,\sin x\right)^2+\left(\sin x+x\,\cos x\right)  ^2}
\]
\end{eulerformula}
\begin{eulerprompt}
>$sol | trigreduce | expand, $integrate(%,x,0,2*pi), %()
\end{eulerprompt}
\begin{eulerformula}
\[
\frac{{\rm asinh}\; \left(2\,\pi\right)+2\,\pi\,\sqrt{4\,\pi^2+1}}{  2}
\]
\end{eulerformula}
\eulerimg{1}{Image Use/TugasAplikom6Adi25030630001-188-large.png}
\begin{euleroutput}
  21.2562941482
\end{euleroutput}
\begin{eulercomment}
Hasilnya sama dengan perhitungan menggunakan fungsi EMT.

Berikut adalah contoh lain penggunaan fungsi Maxima tersebut.
\end{eulercomment}
\begin{eulerprompt}
>plot2d("3*x^2-1","3*x^3-1",xmin=-1/sqrt(3),xmax=1/sqrt(3),square=1):
\end{eulerprompt}
\eulerimg{27}{Image Use/TugasAplikom6Adi25030630001-189.png}
\begin{eulerprompt}
>sol &= radcan(ds(3*x^2-1,3*x^3-1)); $sol
\end{eulerprompt}
\begin{eulerformula}
\[
3\,x\,\sqrt{9\,x^2+4}
\]
\end{eulerformula}
\begin{eulerprompt}
>$showev('integrate(sol,x,0,1/sqrt(3))), $2*float(%) // panjang kurva di atas
\end{eulerprompt}
\begin{eulerformula}
\[
6.0\,\int_{0.0}^{0.5773502691896258}{x\,\sqrt{9.0\,x^2+4.0}\;dx}=  2.337835372767141
\]
\end{eulerformula}
\eulerimg{1}{Image Use/TugasAplikom6Adi25030630001-192-large.png}
\eulersubheading{Sikloid}
\begin{eulercomment}
Berikut kita akan menghitung panjang kurva lintasan (sikloid) suatu
titik pada lingkaran yang berputar ke kanan pada permukaan datar.
Misalkan jari-jari lingkaran tersebut adalah r. Posisi titik pusat
lingkaran pada saat t adalah:

\end{eulercomment}
\begin{eulerformula}
\[
(rt,r).
\]
\end{eulerformula}
\begin{eulercomment}
Misalkan posisi titik pada lingkaran tersebut mula-mula (0,0) dan
posisinya pada saat t adalah:

\end{eulercomment}
\begin{eulerformula}
\[
(r(t-\sin(t)),r(1-\cos(t))).
\]
\end{eulerformula}
\begin{eulercomment}
Berikut kita plot lintasan tersebut dan beberapa posisi lingkaran
ketika t=0, t=pi/2, t=r*pi.
\end{eulercomment}
\begin{eulerprompt}
>r=1;
>ex &= r*(t-sin(t))
\end{eulerprompt}
\begin{euleroutput}
  
                              r (t - sin(t))
  
\end{euleroutput}
\begin{eulerprompt}
>ey &= r*(1-cos(t))
\end{eulerprompt}
\begin{euleroutput}
  
                              r (1 - cos(t))
  
\end{euleroutput}
\begin{eulercomment}
Berikut kita gambar sikloid untuk r=1.
\end{eulercomment}
\begin{eulerprompt}
>  ex &= x-sin(x); ey &= 1-cos(x); aspect(1);
>plot2d(ex,ey,xmin=0,xmax=4pi,square=1); ...
>  plot2d("2+cos(x)","1+sin(x)",xmin=0,xmax=2pi,>add,color=blue); ...
>  plot2d([2,ex(2)],[1,ey(2)],color=red,>add); ...
>  plot2d(ex(2),ey(2),>points,>add,color=red); ...
>  plot2d("2pi+cos(x)","1+sin(x)",xmin=0,xmax=2pi,>add,color=blue); ...
>  plot2d([2pi,ex(2pi)],[1,ey(2pi)],color=red,>add);  ...
>  plot2d(ex(2pi),ey(2pi),>points,>add,color=red):        
\end{eulerprompt}
\begin{euleroutput}
  Syntax error in expression, or unfinished expression!
  Error in:
  ... (ex(2pi),ey(2pi),>points,>add,color=red):         ...
                                                       ^
\end{euleroutput}
\begin{eulercomment}
Berikut dihitung panjang lintasan untuk 1 putaran penuh. (Jangan salah menduga bahwa panjang lintasan 1 putaran penuh sama dengan
keliling lingkaran!)
\end{eulercomment}
\begin{eulerprompt}
>ds &= radcan(sqrt(diff(ex,x)^2+diff(ey,x)^2)); $ds=trigsimp(ds) // elemen panjang kurva sikloid
\end{eulerprompt}
\begin{eulerformula}
\[
\sqrt{\sin ^2x+\cos ^2x-2\,\cos x+1}=\sqrt{2-2\,\cos x}
\]
\end{eulerformula}
\begin{eulerprompt}
>ds &= trigsimp(ds); $ds
\end{eulerprompt}
\begin{eulerformula}
\[
\sqrt{2-2\,\cos x}
\]
\end{eulerformula}
\begin{eulerprompt}
>$showev('integrate(ds,x,0,2*pi)) // hitung panjang sikloid satu putaran penuh
\end{eulerprompt}
\begin{eulerformula}
\[
\int_{0}^{2\,\pi}{\sqrt{2-2\,\cos x}\;dx}=8
\]
\end{eulerformula}
\begin{eulerprompt}
>integrate(mxm("ds"),0,2*pi) // hitung secara numerik
\end{eulerprompt}
\begin{euleroutput}
  8
\end{euleroutput}
\begin{eulerprompt}
>romberg(mxm("ds"),0,2*pi) // cara lain hitung secara numerik
\end{eulerprompt}
\begin{euleroutput}
  8
\end{euleroutput}
\begin{eulercomment}
Perhatikan, seperti terlihat pada gambar, panjang sikloid lebih besar daripada keliling lingkarannya, yakni:

\end{eulercomment}
\begin{eulerformula}
\[
2\pi.
\]
\end{eulerformula}
\eulersubheading{Kurvatur (Kelengkungan) Kurva}
\begin{eulercomment}
image: Osculating.png

Aslinya, kelengkungan kurva diferensiabel (yakni, kurva mulus yang tidak lancip) di titik P didefinisikan melalui lingkaran
oskulasi (yaitu, lingkaran yang melalui titik P dan terbaik memperkirakan, paling banyak menyinggung kurva di sekitar P). Pusat
dan radius kelengkungan kurva di P adalah pusat dan radius lingkaran oskulasi. Kelengkungan adalah kebalikan dari radius
kelengkungan:

\end{eulercomment}
\begin{eulerformula}
\[
\kappa =\frac {1}{R}
\]
\end{eulerformula}
\begin{eulercomment}
dengan R adalah radius kelengkungan. (Setiap lingkaran memiliki kelengkungan ini pada setiap titiknya, dapat diartikan, setiap
lingkaran berputar 2pi sejauh 2piR.)\\
Definisi ini sulit dimanipulasi dan dinyatakan ke dalam rumus untuk kurva umum. Oleh karena itu digunakan definisi lain yang
ekivalen.

\end{eulercomment}
\eulersubheading{Definisi Kurvatur dengan Fungsi Parametrik Panjang Kurva}
\begin{eulercomment}
Setiap kurva diferensiabel dapat dinyatakan dengan persamaan parametrik terhadap panjang kurva s:

\end{eulercomment}
\begin{eulerformula}
\[
\gamma(s) = (x(s),\ y(s)),
\]
\end{eulerformula}
\begin{eulercomment}
dengan x dan y adalah fungsi riil yang diferensiabel, yang memenuhi:

\end{eulercomment}
\begin{eulerformula}
\[
\|\gamma'(s)\|=\sqrt{x'(s)^2+y'(s)^2}=1.
\]
\end{eulerformula}
\begin{eulercomment}
Ini berarti bahwa vektor singgung


\end{eulercomment}
\begin{eulerformula}
\[
\mathbf{T}(s)=(x'(s),\ y'(s))
\]
\end{eulerformula}
\begin{eulercomment}
memiliki norm 1 dan merupakan vektor singgung satuan.

Apabila kurvanya memiliki turunan kedua, artinya turunan kedua x dan y ada, maka T'(s) ada. Vektor ini merupakan normal kurva yang
arahnya menuju pusat kurvatur, norm-nya merupakan nilai kurvatur (kelengkungan):

\end{eulercomment}
\begin{eulerformula}
\[
 \begin{aligned}\mathbf{T}(s) &= \mathbf{\gamma}'(s),\\ \mathbf{T}^{2}(s) &=1\ \text{(konstanta)}\Rightarrow \mathbf{T}'(s)\cdot \mathbf{T}(s)=0\\ \kappa(s) &=\|\mathbf {T}'(s)\|= \|\mathbf{\gamma}''(s)\|=\sqrt{x''(s)^{2}+y''(s)^{2}}.\end{aligned}
\]
\end{eulerformula}
\begin{eulercomment}
Nilai

\end{eulercomment}
\begin{eulerformula}
\[
R(s)=\frac{1}{\kappa(s)}
\]
\end{eulerformula}
\begin{eulercomment}
disebut jari-jari (radius) kelengkungan kurva.

Bilangan riil

\end{eulercomment}
\begin{eulerformula}
\[
 k(s) = \pm\kappa(s)
\]
\end{eulerformula}
\begin{eulercomment}
disebut nilai kelengkungan bertanda.

Contoh:\\
Akan ditentukan kurvatur lingkaran

\end{eulercomment}
\begin{eulerformula}
\[
x=r\cos t,\ y= r\sin t.
\]
\end{eulerformula}
\begin{eulerprompt}
>fx &= r*cos(t); fy &=r*sin(t);
>&assume(t>0,r>0); s &=integrate(sqrt(diff(fx,t)^2+diff(fy,t)^2),t,0,t); s // elemen panjang kurva, panjang busur lingkaran (s)
\end{eulerprompt}
\begin{euleroutput}
  
                                   r t
  
\end{euleroutput}
\begin{eulerprompt}
>&kill(s); fx &= r*cos(s/r); fy &=r*sin(s/r); // definisi ulang persamaan parametrik terhadap s dengan substitusi t=s/r
>k &= trigsimp(sqrt(diff(fx,s,2)^2+diff(fy,s,2)^2)); $k // nilai kurvatur lingkaran dengan menggunakan definisi di atas
\end{eulerprompt}
\begin{eulerformula}
\[
\frac{1}{r}
\]
\end{eulerformula}
\begin{eulercomment}
Untuk representasi parametrik umum, misalkan

\end{eulercomment}
\begin{eulerformula}
\[
x = x(t),\ y= y(t)
\]
\end{eulerformula}
\begin{eulercomment}
merupakan persamaan parametrik untuk kurva bidang yang terdiferensialkan dua kali. Kurvatur untuk kurva tersebut didefinisikan
sebagai

\end{eulercomment}
\begin{eulerformula}
\[
\begin{aligned}\kappa &= \frac{d\phi}{ds}=\frac{\frac{d\phi}{dt}}{\frac{ds}{dt}}\quad (\phi \text{ adalah sudut kemiringan garis singgung dan }s \text{ adalah panjang kurva})\\ &=\frac{\frac{d\phi}{dt}}{\sqrt{(\frac{dx}{dt})^2+(\frac{dy}{dt})^2}}= \frac{\frac{d\phi}{dt}}{\sqrt{x'(t)^2+y'(t)^2}}.\end{aligned}.
\]
\end{eulerformula}
\begin{eulercomment}
Selanjutnya, pembilang pada persamaan di atas dapat dicari sebagai berikut.

\end{eulercomment}
\begin{eulerformula}
\[
\begin{aligned}\sec^2\phi\frac{d\phi}{dt} &= \frac{d}{dt}\left(\tan\phi\right)= \frac{d}{dt}\left(\frac{dy}{dx}\right)= \frac{d}{dt}\left(\frac{dy/dt}{dx/dt}\right)= \frac{d}{dt}\left(\frac{y'(t)}{x'(t)}\right)=\frac{x'(t)y''(t)-x''(t)y'(t)}{x'(t)^2}.\\ & \\ \frac{d\phi}{dt} &= \frac{1}{\sec^2\phi}\frac{x'(t)y''(t)-x''(t)y'(t)}{x'(t)^2}\\ &= \frac{1}{1+\tan^2\phi}\frac{x'(t)y''(t)-x''(t)y'(t)}{x'(t)^2}\\ &= \frac{1}{1+\left(\frac{y'(t)}{x'(t)}\right)^2}\frac{x'(t)y''(t)-x''(t)y'(t)}{x'(t)^2}\\ &= \frac{x'(t)y''(t)-x''(t)y'(t)}{x'(t)^2+y'(t)^2}.\end{aligned}
\]
\end{eulerformula}
\begin{eulercomment}
Jadi, rumus kurvatur untuk kurva parametrik

\end{eulercomment}
\begin{eulerformula}
\[
x=x(t),\ y=y(t)
\]
\end{eulerformula}
\begin{eulercomment}
adalah

\end{eulercomment}
\begin{eulerformula}
\[
\kappa(t) = \frac{x'(t)y''(t)-x''(t)y'(t)}{\left(x'(t)^2+y'(t)^2\right)^{3/2}}.
\]
\end{eulerformula}
\begin{eulercomment}
Jika kurvanya dinyatakan dengan persamaan parametrik pada koordinat
kutub

\end{eulercomment}
\begin{eulerformula}
\[
x=r(\theta)\cos\theta,\ y=r(\theta)\sin\theta,
\]
\end{eulerformula}
\begin{eulercomment}
maka rumus kurvaturnya adalah

\end{eulercomment}
\begin{eulerformula}
\[
\kappa(\theta) = \frac{r(\theta)^2+2r'(\theta)^2-r(\theta)r''(\theta)}{\left(r'(\theta)^2+r'(\theta)^2\right)^{3/2}}.
\]
\end{eulerformula}
\begin{eulercomment}
(Silakan Anda turunkan rumus tersebut!)

Contoh:\\
Lingkaran dengan pusat (0,0) dan jari-jari r dapat dinyatakan dengan
persamaan parametrik

\end{eulercomment}
\begin{eulerformula}
\[
x=r\cos t,\ y=r\sin t.
\]
\end{eulerformula}
\begin{eulercomment}
Nilai kelengkungan lingkaran tersebut adalah

\end{eulercomment}
\begin{eulerformula}
\[
\kappa(t)=\frac{x'(t)y''(t)-x''(t)y'(t)}{\left(x'(t)^2+y'(t)^2\right)^{3/2}}=\frac{r^2}{r^3}=\frac 1 r.
\]
\end{eulerformula}
\begin{eulercomment}
Hasil cocok dengan definisi kurvatur suatu kelengkungan.
\end{eulercomment}
\begin{eulercomment}
Kurva

\end{eulercomment}
\begin{eulerformula}
\[
y=f(x)
\]
\end{eulerformula}
\begin{eulercomment}
dapat dinyatakan ke dalam persamaan parametrik

\end{eulercomment}
\begin{eulerformula}
\[
x=t,\ y=f(t),\ \text{ dengan } x'(t)=1,\ x''(t)=0,
\]
\end{eulerformula}
\begin{eulercomment}
sehingga kurvaturnya adalah

\end{eulercomment}
\begin{eulerformula}
\[
\kappa(t) = \frac{y''(t)}{\left(1+y'(t)^2\right)^{3/2}}.
\]
\end{eulerformula}
\begin{eulercomment}
Contoh:\\
Akan ditentukan kurvatur parabola

\end{eulercomment}
\begin{eulerformula}
\[
y=ax^2+bx+c.
\]
\end{eulerformula}
\begin{eulerprompt}
>function f(x) &= a*x^2+b*x+c; $y=f(x)
\end{eulerprompt}
\begin{eulerformula}
\[
y=a\,x^2+b\,x+c
\]
\end{eulerformula}
\begin{eulerprompt}
>function k(x) &= (diff(f(x),x,2))/(1+diff(f(x),x)^2)^(3/2); $'k(x)=k(x) // kelengkungan parabola 
\end{eulerprompt}
\begin{eulerformula}
\[
k\left(x\right)=\frac{2\,a}{\left(\left(2\,a\,x+b\right)^2+1\right)  ^{\frac{3}{2}}}
\]
\end{eulerformula}
\begin{eulerprompt}
>function f(x) &= x^2+x+1; $y=f(x) // akan kita plot kelengkungan parabola untuk a=b=c=1
\end{eulerprompt}
\begin{eulerformula}
\[
y=x^2+x+1
\]
\end{eulerformula}
\begin{eulerprompt}
>function k(x) &= (diff(f(x),x,2))/(1+diff(f(x),x)^2)^(3/2); $'k(x)=k(x) // kelengkungan parabola 
\end{eulerprompt}
\begin{eulerformula}
\[
k\left(x\right)=\frac{2}{\left(\left(2\,x+1\right)^2+1\right)^{  \frac{3}{2}}}
\]
\end{eulerformula}
\begin{eulercomment}
Berikut kita gambar parabola tersebut beserta kurva kelengkungan, kurva jari-jari kelengkungan dan salah satu lingkaran oskulasi
di titik puncak parabola. Perhatikan, puncak parabola dan jari-jari lingkaran oskulasi di puncak parabola adalah

\end{eulercomment}
\begin{eulerformula}
\[
(-1/2,3/4),\ 1/k(2)=1/2,
\]
\end{eulerformula}
\begin{eulercomment}
sehingga pusat lingkaran oskulasi adalah (-1/2, 5/4).
\end{eulercomment}
\begin{eulerprompt}
>plot2d(["f(x)", "k(x)"],-2,1, color=[blue,red]); plot2d("1/k(x)",-1.5,1,color=green,>add); ...
>plot2d("-1/2+1/k(-1/2)*cos(x)","5/4+1/k(-1/2)*sin(x)",xmin=0,xmax=2pi,>add,color=blue):
\end{eulerprompt}
\eulerimg{27}{Image Use/TugasAplikom6Adi25030630001-211.png}
\begin{eulercomment}
Untuk kurva yang dinyatakan dengan fungsi implisit

\end{eulercomment}
\begin{eulerformula}
\[
F(x,y)=0
\]
\end{eulerformula}
\begin{eulercomment}
dengan turunan-turunan parsial

\end{eulercomment}
\begin{eulerformula}
\[
F_x=\frac{\partial F}{\partial x},\ F_y=\frac{\partial F}{\partial y},\ F_{xy}=\frac{\partial}{\partial y}\left(\frac{\partial F}{\partial x}\right),\ F_{xx}=\frac{\partial}{\partial x}\left(\frac{\partial F}{\partial x}\right),\ F_{yy}=\frac{\partial}{\partial y}\left(\frac{\partial F}{\partial y}\right),
\]
\end{eulerformula}
\begin{eulercomment}
berlaku

\end{eulercomment}
\begin{eulerformula}
\[
F_x dx+ F_y dy = 0\text{ atau } \frac{dy}{dx}=-\frac{F_x}{F_y},
\]
\end{eulerformula}
\begin{eulercomment}
sehingga kurvaturnya adalah

\end{eulercomment}
\begin{eulerformula}
\[
\kappa =\frac {F_y^2F_{xx}-2F_xF_yF_{xy}+F_x^2F_{yy}}{\left(F_x^2+F_y^2\right)^{3/2}}.
\]
\end{eulerformula}
\begin{eulercomment}
(Silakan Anda turunkan sendiri!)

Contoh 1:\\
Parabola 

\end{eulercomment}
\begin{eulerformula}
\[
y=ax^2+bx+c
\]
\end{eulerformula}
\begin{eulercomment}
dapat dinyatakan ke dalam persamaan implisit

\end{eulercomment}
\begin{eulerformula}
\[
ax^2+bx+c-y=0.
\]
\end{eulerformula}
\begin{eulerprompt}
>function F(x,y) &=a*x^2+b*x+c-y; $F(x,y)
\end{eulerprompt}
\begin{eulerformula}
\[
-y+a\,x^2+b\,x+c
\]
\end{eulerformula}
\begin{eulerprompt}
>Fx &= diff(F(x,y),x), Fxx &=diff(F(x,y),x,2), Fy &=diff(F(x,y),y), Fxy &=diff(diff(F(x,y),x),y), Fyy &=diff(F(x,y),y,2)  
\end{eulerprompt}
\begin{euleroutput}
  
                                2 a x + b
  
  
                                   2 a
  
  
                                   - 1
  
  
                                    0
  
  
                                    0
  
\end{euleroutput}
\begin{eulerprompt}
>function k(x) &= (Fy^2*Fxx-2*Fx*Fy*Fxy+Fx^2*Fyy)/(Fx^2+Fy^2)^(3/2); $'k(x)=k(x) // kurvatur parabola tersebut
\end{eulerprompt}
\begin{eulerformula}
\[
k\left(x\right)=\frac{2\,a}{\left(\left(2\,a\,x+b\right)^2+1\right)  ^{\frac{3}{2}}}
\]
\end{eulerformula}
\begin{eulercomment}
Hasilnya sama dengan sebelumnya yang menggunakan persamaan parabola biasa.
\end{eulercomment}
\eulerheading{Latihan}
\begin{eulercomment}
- Bukalah buku Kalkulus.\\
- Cari dan pilih beberapa (paling sedikit 5 fungsi berbeda tipe/bentuk/jenis) fungsi dari buku tersebut, kemudian definisikan di
EMT pada baris-baris perintah berikut (jika perlu tambahkan lagi).\\
- Untuk setiap fungsi, tentukan anti turunannya (jika ada), hitunglah integral tentu dengan batas-batas yang menarik (Anda
tentukan sendiri), seperti contoh-contoh tersebut.\\
- Lakukan hal yang sama untuk fungsi-fungsi yang tidak dapat diintegralkan (cari sedikitnya 3 fungsi).\\
- Gambar grafik fungsi dan daerah integrasinya pada sumbu koordinat yang sama.\\
- Gunakan integral tentu untuk mencari luas daerah yang dibatasi oleh dua kurva yang berpotongan di dua titik. (Cari dan gambar
kedua kurva dan arsir (warnai) daerah yang dibatasi oleh keduanya.)\\
- Gunakan integral tentu untuk menghitung volume benda putar kurva y= f(x) yang diputar mengelilingi sumbu x dari x=a sampai x=b,
yakni

\end{eulercomment}
\begin{eulerformula}
\[
V = \int_a^b \pi (f(x)^2\ dx.
\]
\end{eulerformula}
\begin{eulercomment}
(Pilih fungsinya dan gambar kurva dan benda putar yang dihasilkan. Anda dapat mencari contoh-contoh bagaimana cara menggambar
benda hasil perputaran suatu kurva.)\\
- Gunakan integral tentu untuk menghitung panjang kurva y=f(x) dari x=a sampai x=b dengan menggunakan rumus:

\end{eulercomment}
\begin{eulerformula}
\[
S = \int_a^b \sqrt{1+(f'(x))^2} \ dx.
\]
\end{eulerformula}
\begin{eulercomment}
(Pilih fungsi dan gambar kurvanya.)\\
- Apabila fungsi dinyatakan dalam koordinat kutub x=f(r,t), y=g(r,t), r=h(t), x=a bersesuaian dengan t=t0 dan x=b bersesuian
dengan t=t1, maka rumus di atas akan menjadi:

\end{eulercomment}
\begin{eulerformula}
\[
S=\int_{t_0}^{t_1} \sqrt{x'(t)^2+y'(t)^2}\ dt.
\]
\end{eulerformula}
\begin{eulercomment}
- Pilih beberapa kurva menarik (selain lingkaran dan parabola) dari buku  kalkulus. Nyatakan setiap kurva tersebut dalam bentuk:\\
\end{eulercomment}
\begin{eulerttcomment}
  a. koordinat Kartesius (persamaan y=f(x))
  b. koordinat kutub ( r=r(theta))
  c. persamaan parametrik x=x(t), y=y(t)
  d. persamaan implit F(x,y)=0
\end{eulerttcomment}
\begin{eulercomment}
- Tentukan kurvatur masing-masing kurva dengan menggunakan keempat representasi tersebut (hasilnya harus sama).\\
- Gambarlah kurva asli, kurva kurvatur, kurva jari-jari lingkaran oskulasi, dan salah satu lingkaran oskulasinya.
\end{eulercomment}
\begin{eulercomment}
==============================================================================
LEMBAR KERJA EULER MATH TOOLBOX (EMT): LATIHAN KALKULUS INTEGRASI
==============================================================================
\end{eulercomment}
\begin{eulerprompt}
>reset();
\end{eulerprompt}
\eulersubheading{}
\begin{eulercomment}
NO 1. DEFINISI \& INTEGRASI 5 FUNGSI BERBEDA (BISA DIINTEGRALKAN)
\end{eulercomment}
\eulersubheading{}
\begin{eulercomment}
Fungsi 1: Polinomial f1(x) = x\textasciicircum{}3 - 3*x
\end{eulercomment}
\begin{eulerprompt}
>f1 &= x^3 - 3*x;
>$'anti_f1 = integrate(f1, x)
\end{eulerprompt}
\begin{eulerformula}
\[
{\it anti\_\_f_1}=\frac{x^4}{4}-\frac{3\,x^2}{2}
\]
\end{eulerformula}
\begin{eulerprompt}
>hasil1 &= integrate(f1, x, 0, 2);
>$hasil1
\end{eulerprompt}
\begin{eulerformula}
\[
-2
\]
\end{eulerformula}
\begin{eulercomment}
Fungsi 2: Trigonometri f2(x) = sin(x)*cos(x)
\end{eulercomment}
\begin{eulerprompt}
>f2 &= sin(x)*cos(x);
>$'anti_f2 = integrate(f2, x)
\end{eulerprompt}
\begin{eulerformula}
\[
{\it anti\_\_f_2}=-\frac{\cos ^2x}{2}
\]
\end{eulerformula}
\begin{eulerprompt}
>hasil2 &= integrate(f2, x, 0, pi/2);
>$hasil2
\end{eulerprompt}
\begin{eulerformula}
\[
\frac{1}{2}
\]
\end{eulerformula}
\begin{eulercomment}
Fungsi 3: Eksponensial f3(x) = x*exp(x)
\end{eulercomment}
\begin{eulerprompt}
>f3 &= x*exp(x);
>$'anti_f3 = integrate(f3, x)
\end{eulerprompt}
\begin{eulerformula}
\[
{\it anti\_\_f_3}=\left(x-1\right)\,e^{x}
\]
\end{eulerformula}
\begin{eulerprompt}
>hasil3 &= integrate(f3, x, 0, 1);
>$hasil3
\end{eulerprompt}
\begin{eulerformula}
\[
1
\]
\end{eulerformula}
\begin{eulercomment}
Fungsi 4: Rasional f4(x) = 1/(x\textasciicircum{}2 + 1)
\end{eulercomment}
\begin{eulerprompt}
>f4 &= 1/(x^2 + 1);
>$'anti_f4 = integrate(f4, x)
\end{eulerprompt}
\begin{eulerformula}
\[
{\it anti\_\_f_4}=\arctan x
\]
\end{eulerformula}
\begin{eulerprompt}
>hasil4 &= integrate(f4, x, 0, 1);
>$hasil4
\end{eulerprompt}
\begin{eulerformula}
\[
\frac{\pi}{4}
\]
\end{eulerformula}
\begin{eulercomment}
Fungsi 5: Logaritma f5(x) = log(x)
\end{eulercomment}
\begin{eulerprompt}
>f5 &= log(x);
>$'anti_f5 = integrate(f5, x)
\end{eulerprompt}
\begin{eulerformula}
\[
{\it anti\_\_f_5}=x\,\log x-x
\]
\end{eulerformula}
\begin{eulerprompt}
>hasil5 &= integrate(f5, x, 1, e);
\end{eulerprompt}
\begin{euleroutput}
  Answering "Is e-1 positive, negative or zero?" with "positive"
\end{euleroutput}
\begin{eulerprompt}
>$hasil5
\end{eulerprompt}
\begin{eulerformula}
\[
e\,\log e-e+1
\]
\end{eulerformula}
\eulersubheading{}
\begin{eulercomment}
NO 2. FUNGSI YANG TIDAK MEMILIKI ANTI-TURUNAN ELEMENTER (3 FUNGSI)
\end{eulercomment}
\eulersubheading{}
\begin{eulercomment}
Menggunakan metode integrasi numerik bawaan EMT (romberg / integrate numerik)
Fungsi A: fA(x) = exp(-x\textasciicircum{}2) [Fungsi Gaussi]
\end{eulercomment}
\begin{eulerprompt}
>fA &= exp(-x^2);
>romberg("exp(-x^2)", 0, 1)
\end{eulerprompt}
\begin{euleroutput}
  0.746824132814
\end{euleroutput}
\begin{eulercomment}
Fungsi B: fB(x) = sin(x)/x [Fungsi Sinc]
\end{eulercomment}
\begin{eulerprompt}
>fB &= sin(x)/x;
>romberg("sin(x)/x", 0.000001, pi)
\end{eulerprompt}
\begin{euleroutput}
  1.85193605198
\end{euleroutput}
\begin{eulercomment}
Fungsi C: fC(x) = sqrt(1 + x\textasciicircum{}4)
\end{eulercomment}
\begin{eulerprompt}
>fC &= sqrt(1 + x^4);
>romberg("sqrt(1 + x^4)", 0, 1)
\end{eulerprompt}
\begin{euleroutput}
  1.08942941322
\end{euleroutput}
\eulersubheading{}
\begin{eulercomment}
NO 3. GRAFIK FUNGSI DAN DAERAH INTEGRASINYA
\end{eulercomment}
\eulersubheading{}
\begin{eulercomment}
Menggambar f1(x) = x\textasciicircum{}3 - 3*x pada rentang [0, 2] beserta daerah arsirnya
\end{eulercomment}
\begin{eulerprompt}
>reset();
>plot2d("x^3 - 3*x", 0, 2); ...
>plot2d("x^3 - 3*x", 0, 2, >filled, color=lightgray, >add);
\end{eulerprompt}
\eulersubheading{}
\begin{eulercomment}
NO 4. LUAS DAERAH DIBATASI DUA KURVA BERSILANGAN
\end{eulercomment}
\eulersubheading{}
\begin{eulercomment}
Kurva y1 = x\textasciicircum{}2 dan y2 = x + 2 (Berpotongan di x = -1 dan x = 2)
\end{eulercomment}
\begin{eulerprompt}
>reset();
>y1 &= x^2; y2 &= x + 2;
>Luas &= integrate(y2 - y1, x, -1, 2)
\end{eulerprompt}
\begin{euleroutput}
  
                                    9
                                    -
                                    2
  
\end{euleroutput}
\begin{eulercomment}
Plot dua kurva dan daerah di antaranya
\end{eulercomment}
\begin{eulerprompt}
>plot2d(["x^2", "x + 2"], -2, 3, color=[blue, red]); ...
>plot2d("x + 2", -1, 2, >filled, fillcolor=lightgray, >add); ...
>plot2d("x^2", -1, 2, >filled, fillcolor=white, >add); ...
>plot2d(["x^2", "x + 2"], -2, 3, color=[blue, red], >add);
\end{eulerprompt}
\eulersubheading{}
\begin{eulercomment}
NO 5. VOLUME BENDA PUTAR MEMENGELILINGI SUMBU X
\end{eulercomment}
\eulersubheading{}
\begin{eulercomment}
Kurva f(x) = sqrt(x) diputar mengelilingi sumbu x dari x = 0 sampai x = 4
Rumus: V = pi * integrate(f(x)\textasciicircum{}2, x, a, b)
\end{eulercomment}
\begin{eulerprompt}
>reset();
>f &= sqrt(x);
>Volume &= pi * integrate(f^2, x, 0, 4)
\end{eulerprompt}
\begin{euleroutput}
  
                                   8 pi
  
\end{euleroutput}
\begin{eulercomment}
Gambar kurva 2D pembentuknya
\end{eulercomment}
\begin{eulerprompt}
>plot2d("sqrt(x)", 0, 4);
\end{eulerprompt}
\eulersubheading{}
\begin{eulercomment}
NO 6. PANJANG KURVA KOORDINAT KARTESIUS y = f(x)
\end{eulercomment}
\eulersubheading{}
\begin{eulercomment}
Kurva y = (2/3)*x\textasciicircum{}(3/2) dari x = 0 sampai x = 3
Rumus: S = integrate(sqrt(1 + (f'(x))\textasciicircum{}2), x, a, b)
\end{eulercomment}
\begin{eulerprompt}
>reset();
>fy &= (2/3)*x^(3/2);
>dfy &= diff(fy, x);
>dS &= radcan(sqrt(1 + dfy^2));
>PanjangKartesius &= integrate(dS, x, 0, 3)
\end{eulerprompt}
\begin{euleroutput}
  
                                    14
                                    --
                                    3
  
\end{euleroutput}
\eulersubheading{}
\begin{eulercomment}
NO 7. PANJANG KURVA KOORDINAT KUTUB (POLAR)
\end{eulercomment}
\eulersubheading{}
\begin{eulercomment}
Kurva Kardioida: r(t) = 1 + cos(t) dari t = 0 sampai t = 2*pi
Rumus: S = integrate(sqrt(x'(t)\textasciicircum{}2 + y'(t)\textasciicircum{}2), t, t0, t1)
\end{eulercomment}
\begin{eulerprompt}
>reset();
>rt &= 1 + cos(t);
>xt &= rt * cos(t);
>yt &= rt * sin(t);
>dSpolar &= trigreduce(radcan(sqrt(diff(xt,t)^2 + diff(yt,t)^2)));
>PanjangPolar &= integrate(dSpolar, t, 0, 2*pi)
\end{eulerprompt}
\begin{euleroutput}
  
                                    8
  
\end{euleroutput}
\eulersubheading{}
\begin{eulercomment}
NO 8. REPRESENSI KURVA MENARIK (ASTROID) DALAM 4 BENTUK
\end{eulercomment}
\eulersubheading{}
\begin{eulerprompt}
>reset();
\end{eulerprompt}
\begin{eulercomment}
a. Koordinat Kartesius: y = (1 - x\textasciicircum{}(2/3))\textasciicircum{}(3/2) 
\end{eulercomment}
\begin{eulerprompt}
>yatas &= (1 - abs(x)^(2/3))^(3/2);
>$'y_atas = y_atas
\end{eulerprompt}
\begin{eulerformula}
\[
{\it y\_\_atas}={\it y\_\_atas}
\]
\end{eulerformula}
\begin{eulerprompt}
>ybawah &= -(1 - abs(x)^(2/3))^(3/2);
>$'y_bawah = y_bawah
\end{eulerprompt}
\begin{eulerformula}
\[
{\it y\_\_bawah}={\it y\_\_bawah}
\]
\end{eulerformula}
\begin{eulerprompt}
>plot2d("(1 - abs(x)^(2/3))^(3/2)", -1, 1, square=1, color=blue);
>plot2d("-(1 - abs(x)^(2/3))^(3/2)", -1, 1, >add, color=blue);
\end{eulerprompt}
\begin{eulercomment}
b. Koordinat Kutub: r(theta) = cos(theta)*sin(theta) (Bentuk turunan) 
\end{eulercomment}
\begin{eulerprompt}
>reset();
>t = linspace(0, 2*pi, 1000);
>r = 1 ./ (abs(cos(t)).^(2/3) + abs(sin(t)).^(2/3)).^(3/2);
>x = r .* cos(t);
>y = r .* sin(t);
>plot2d(x, y, square=1, color=red);  
\end{eulerprompt}
\begin{eulercomment}
c. Persamaan Parametrik: x(t) = cos(t)\textasciicircum{}3, y(t) = sin(t)\textasciicircum{}3
\end{eulercomment}
\begin{eulerprompt}
>xt &= cos(t)^3; 
\end{eulerprompt}
\begin{euleroutput}
  Error : ypol(t) does not produce a real or column vector
  
  Error generated by error() command
  
   %ploteval:
      error(f$|" does not produce a real or column vector"); 
  adaptiveevalone:
      s=%ploteval(g$,t;args());
  Try "trace errors" to inspect local variables after errors.
  plot2d:
      dw/n,dw/n^2,dw/n,auto;args());
\end{euleroutput}
\begin{eulerprompt}
>yt &= sin(t)^3;
\end{eulerprompt}
\begin{eulercomment}
d. Persamaan Implisit: x\textasciicircum{}(2/3) + y\textasciicircum{}(2/3) - 1 = 0
Gambar Kurva Astroid Parametrik
\end{eulercomment}
\begin{eulerprompt}
>Fimpl &= abs(x)^(2/3) + abs(y)^(2/3) - 1;
>$'F(x,y) = Fimpl
\end{eulerprompt}
\begin{eulerformula}
\[
F\left(x , y\right)=-1+x^{\frac{2}{3}}+y^{\frac{2}{3}}
\]
\end{eulerformula}
\begin{eulerprompt}
>plot2d("abs(x)^(2/3) + abs(y)^(2/3) - 1", r=1.5, level=0, square=1, color=blue);
>  
\end{eulerprompt}
\eulerheading{Barisan dan Deret}
\begin{eulercomment}
(Catatan: bagian ini belum lengkap. Anda dapat membaca contoh-contoh pengguanaan EMT dan
Maxima untuk menghitung limit barisan, rumus jumlah parsial suatu deret, jumlah tak hingga
suatu deret konvergen, dan sebagainya. Anda dapat mengeksplor contoh-contoh di EMT atau
perbagai panduan penggunaan Maxima di software Maxima atau dari Internet.)

Barisan dapat didefinisikan dengan beberapa cara di dalam EMT, di antaranya:

- dengan cara yang sama seperti mendefinisikan vektor dengan elemen-elemen beraturan
(menggunakan titik dua ":");\\
- menggunakan perintah "sequence" dan rumus barisan (suku ke -n);\\
- menggunakan perintah "iterate" atau "niterate";\\
- menggunakan fungsi Maxima "create\_list" atau "makelist" untuk menghasilkan barisan
simbolik;\\
- menggunakan fungsi biasa yang inputnya vektor atau barisan;\\
- menggunakan fungsi rekursif.

EMT menyediakan beberapa perintah (fungsi) terkait barisan, yakni:

- sum: menghitung jumlah semua elemen suatu barisan\\
- cumsum: jumlah kumulatif suatu barisan\\
- differences: selisih antar elemen-elemen berturutan

EMT juga dapat digunakan untuk menghitung jumlah deret berhingga maupun deret tak hingga,
dengan menggunakan perintah (fungsi) "sum". Perhitungan dapat dilakukan secara numerik
maupun simbolik dan eksak.

Berikut adalah beberapa contoh perhitungan barisan dan deret menggunakan EMT.
\end{eulercomment}
\begin{eulerprompt}
>1:10 // barisan sederhana
\end{eulerprompt}
\begin{euleroutput}
  [1,  2,  3,  4,  5,  6,  7,  8,  9,  10]
\end{euleroutput}
\begin{eulerprompt}
>1:2:30
\end{eulerprompt}
\begin{euleroutput}
  [1,  3,  5,  7,  9,  11,  13,  15,  17,  19,  21,  23,  25,  27,  29]
\end{euleroutput}
\eulerheading{Iterasi dan Barisan}
\begin{eulercomment}
EMT menyediakan fungsi iterate("g(x)", x0, n) untuk melakukan iterasi

\end{eulercomment}
\begin{eulerformula}
\[
x_{k+1}=g(x_k), \ x_0=x_0, k= 1, 2, 3, ..., n.
\]
\end{eulerformula}
\begin{eulercomment}
Berikut ini disajikan contoh-contoh penggunaan iterasi dan rekursi dengan EMT. Contoh
pertama menunjukkan pertumbuhan dari nilai awal 1000 dengan laju pertambahan 5\%, selama 10
periode.
\end{eulercomment}
\begin{eulerprompt}
>q=1.05; iterate("x*q",1000,n=10)'
\end{eulerprompt}
\begin{euleroutput}
           1000 
           1050 
         1102.5 
        1157.63 
        1215.51 
        1276.28 
         1340.1 
         1407.1 
        1477.46 
        1551.33 
        1628.89 
\end{euleroutput}
\begin{eulercomment}
Contoh berikutnya memperlihatkan bahaya menabung di bank pada masa sekarang! Dengan bunga
tabungan sebesar 6\% per tahun atau 0.5\% per bulan dipotong pajak 20\%, dan biaya administrasi
10000 per bulan, tabungan sebesar 1 juta tanpa diambil selama sekitar 10 tahunan akan habis
diambil oleh bank!
\end{eulercomment}
\begin{eulerprompt}
>r=0.005; plot2d(iterate("(1+0.8*r)*x-10000",1000000,n=130)):
\end{eulerprompt}
\eulerimg{27}{Image Use/TugasAplikom6Adi25030630001-227.png}
\begin{eulercomment}
Silakan Anda coba-coba, dengan tabungan minimal berapa agar tidak akan
habis diambil oleh bank dengan ketentuan bunga dan biaya administrasi
seperti di atas.
\end{eulercomment}
\begin{eulercomment}
1. Bersihkan memori total
\end{eulercomment}
\begin{eulerprompt}
>reset();
\end{eulerprompt}
\begin{eulercomment}
2. Definisikan parameter bunga (r) dan biaya administrasi (a)
\end{eulercomment}
\begin{eulerprompt}
>r = 0.005;
>a = 10000;
\end{eulerprompt}
\begin{eulercomment}
3. Hitung tabungan minimal secara presisi
\end{eulercomment}
\begin{eulerprompt}
>  Mmin = a / (0.8 * r)
\end{eulerprompt}
\begin{euleroutput}
  2500000
\end{euleroutput}
\begin{eulercomment}
4. Buktikan dengan iterasi \& plot selama 130 bulan
\end{eulercomment}
\begin{eulerprompt}
>plot2d(iterate("(1+0.8*r)*x-10000", Mmin, n=130)):
\end{eulerprompt}
\eulerimg{27}{Image Use/TugasAplikom6Adi25030630001-228.png}
\begin{eulercomment}
Berikut adalah perhitungan minimal tabungan agar aman di bank dengan
bunga sebesar r dan biaya administrasi a, pajak bunga 20\%.
\end{eulercomment}
\begin{eulerprompt}
>$solve(0.8*r*A-a,A), $% with [r=0.005, a=10] 
\end{eulerprompt}
\begin{eulerformula}
\[
\left[ A=2500.0 \right] 
\]
\end{eulerformula}
\eulerimg{0}{Image Use/TugasAplikom6Adi25030630001-230-large.png}
\begin{eulercomment}
Berikut didefinisikan fungsi untuk menghitung saldo tabungan, kemudian dilakukan iterasi.
\end{eulercomment}
\begin{eulerprompt}
>function saldo(x,r,a) := round((1+0.8*r)*x-a,2);
>iterate(\{\{"saldo",0.005,10\}\},1000,n=6)
\end{eulerprompt}
\begin{euleroutput}
  [1000,  994,  987.98,  981.93,  975.86,  969.76,  963.64]
\end{euleroutput}
\begin{eulerprompt}
>iterate(\{\{"saldo",0.005,10\}\},2000,n=6)
\end{eulerprompt}
\begin{euleroutput}
  [2000,  1998,  1995.99,  1993.97,  1991.95,  1989.92,  1987.88]
\end{euleroutput}
\begin{eulerprompt}
>iterate(\{\{"saldo",0.005,10\}\},2500,n=6)
\end{eulerprompt}
\begin{euleroutput}
  [2500,  2500,  2500,  2500,  2500,  2500,  2500]
\end{euleroutput}
\begin{eulercomment}
Tabungan senilai 2,5 juta akan aman dan tidak akan berubah nilai (jika tidak ada penarikan),
sedangkan jika tabungan awal kurang dari 2,5 juta, lama kelamaan akan berkurang meskipun
tidak pernah dilakukan penarikan uang tabungan.
\end{eulercomment}
\begin{eulerprompt}
>iterate(\{\{"saldo",0.005,10\}\},3000,n=6)
\end{eulerprompt}
\begin{euleroutput}
  [3000,  3002,  3004.01,  3006.03,  3008.05,  3010.08,  3012.12]
\end{euleroutput}
\begin{eulercomment}
Tabungan yang lebih dari 2,5 juta baru akan bertambah jika tidak ada penarikan.

Untuk barisan yang lebih kompleks dapat digunakan fungsi "sequence()". Fungsi ini menghitung
nilai-nilai x[n] dari semua nilai sebelumnya, x[1],...,x[n-1] yang diketahui.\\
Berikut adalah contoh barisan Fibonacci.

\end{eulercomment}
\begin{eulerformula}
\[
x_n = x_{n-1}+x_{n-2}, \quad x_1=1, \quad x_2 =1
\]
\end{eulerformula}
\begin{eulerprompt}
>sequence("x[n-1]+x[n-2]",[1,1],15)
\end{eulerprompt}
\begin{euleroutput}
  [1,  1,  2,  3,  5,  8,  13,  21,  34,  55,  89,  144,  233,  377,  610]
\end{euleroutput}
\begin{eulercomment}
Barisan Fibonacci memiliki banyak sifat menarik, salah satunya adalah akar pangkat ke-n suku
ke-n akan konvergen ke pecahan emas:
\end{eulercomment}
\begin{eulerprompt}
>$'(1+sqrt(5))/2=float((1+sqrt(5))/2)
\end{eulerprompt}
\begin{eulerformula}
\[
\frac{\sqrt{5}+1}{2}=1.618033988749895
\]
\end{eulerformula}
\begin{eulerprompt}
>plot2d(sequence("x[n-1]+x[n-2]",[1,1],250)^(1/(1:250))):
\end{eulerprompt}
\eulerimg{27}{Image Use/TugasAplikom6Adi25030630001-232.png}
\begin{eulercomment}
Barisan yang sama juga dapat dihasilkan dengan menggunakan loop.
\end{eulercomment}
\begin{eulerprompt}
>x=ones(500); for k=3 to 500; x[k]=x[k-1]+x[k-2]; end;
\end{eulerprompt}
\begin{eulercomment}
Rekursi dapat dilakukan dengan menggunakan rumus yang tergantung pada semua elemen
sebelumnya. Pada contoh berikut, elemen ke-n merupakan jumlah (n-1) elemen sebelumnya,
dimulai dengan 1 (elemen ke-1). Jelas, nilai elemen ke-n adalah 2\textasciicircum{}(n-2), untuk n=2, 4, 5,
....
\end{eulercomment}
\begin{eulerprompt}
>sequence("sum(x)",1,10)
\end{eulerprompt}
\begin{euleroutput}
  [1,  1,  2,  4,  8,  16,  32,  64,  128,  256]
\end{euleroutput}
\begin{eulercomment}
Selain menggunakan ekspresi dalam x dan n, kita juga dapat menggunakan fungsi.

Pada contoh berikut, digunakan iterasi

\end{eulercomment}
\begin{eulerformula}
\[
x_n =A \cdot x_{n-1},
\]
\end{eulerformula}
\begin{eulercomment}
dengan A suatu matriks 2x2, dan setiap x[n] merupakan matriks/vektor 2x1.
\end{eulercomment}
\begin{eulerprompt}
>A=[1,1;1,2]; function suku(x,n) := A.x[,n-1]
>sequence("suku",[1;1],6)
\end{eulerprompt}
\begin{euleroutput}
  Real 2 x 6 matrix
  
              1             2             5            13     ...
              1             3             8            21     ...
\end{euleroutput}
\begin{eulercomment}
Hasil yang sama juga dapat diperoleh dengan menggunakan fungsi perpangkatan matriks
"matrixpower()". Cara ini lebih cepat, karena hanya menggunakan perkalian matriks sebanyak
log\_2(n).

\end{eulercomment}
\begin{eulerformula}
\[
x_n=A.x_{n-1}=A^2.x_{n-2}=A^3.x_{n-3}= ... = A^{n-1}.x_1.
\]
\end{eulerformula}
\begin{eulerprompt}
>sequence("matrixpower(A,n).[1;1]",1,6)
\end{eulerprompt}
\begin{euleroutput}
  Real 2 x 6 matrix
  
              1             5            13            34     ...
              1             8            21            55     ...
\end{euleroutput}
\eulerheading{Spiral Theodorus}
\begin{eulercomment}
image: Spiral\_of\_Theodorus.png\\
Spiral Theodorus (spiral segitiga siku-siku) dapat digambar secara rekursif. Rumus
rekursifnya adalah:

\end{eulercomment}
\begin{eulerformula}
\[
x_n = \left( 1 + \frac{i}{\sqrt{n-1}} \right) \, x_{n-1}, \quad x_1=1,
\]
\end{eulerformula}
\begin{eulercomment}
yang menghasilkan barisan bilangan kompleks.
\end{eulercomment}
\begin{eulerprompt}
>function g(n) := 1+I/sqrt(n)
\end{eulerprompt}
\begin{eulercomment}
Rekursinya dapat dijalankan sebanyak 17 untuk menghasilkan barisan 17 bilangan kompleks,
kemudian digambar bilangan-bilangan kompleksnya.
\end{eulercomment}
\begin{eulerprompt}
>x=sequence("g(n-1)*x[n-1]",1,17); plot2d(x,r=3.5); textbox(latex("Spiral\(\backslash\) Theodorus"),0.4):
\end{eulerprompt}
\eulerimg{27}{Image Use/TugasAplikom6Adi25030630001-233.png}
\begin{eulercomment}
Selanjutnya dihubungan titik 0 dengan titik-titik kompleks tersebut menggunakan loop.
\end{eulercomment}
\begin{eulerprompt}
>for i=1:cols(x); plot2d([0,x[i]],>add); end:
\end{eulerprompt}
\eulerimg{27}{Image Use/TugasAplikom6Adi25030630001-234.png}
\begin{eulerprompt}
> 
\end{eulerprompt}
\begin{eulercomment}
Spiral tersebut juga dapat didefinisikan menggunakan fungsi rekursif, yang tidak memmerlukan
indeks dan bilangan kompleks. Dalam hal ini diigunakan vektor kolom pada bidang.
\end{eulercomment}
\begin{eulerprompt}
>function gstep (v) ...
\end{eulerprompt}
\begin{eulerudf}
  w=[-v[2];v[1]];
  return v+w/norm(w);
  endfunction
\end{eulerudf}
\begin{eulercomment}
Jika dilakukan iterasi 16 kali dimulai dari [1;0] akan didapatkan matriks yang memuat
vektor-vektor dari setiap iterasi.
\end{eulercomment}
\begin{eulerprompt}
>x=iterate("gstep",[1;0],16); plot2d(x[1],x[2],r=3.5,>points):
\end{eulerprompt}
\eulerimg{27}{Image Use/TugasAplikom6Adi25030630001-235.png}
\begin{eulercomment}
\begin{eulercomment}
\eulerheading{Kekonvergenan}
\begin{eulercomment}
Terkadang kita ingin melakukan iterasi sampai konvergen. Apabila iterasinya tidak konvergen
setelah ditunggu lama, Anda dapat menghentikannya dengan menekan tombol [ESC].
\end{eulercomment}
\begin{eulerprompt}
>iterate("cos(x)",1) // iterasi x(n+1)=cos(x(n)), dengan x(0)=1.
\end{eulerprompt}
\begin{euleroutput}
  0.739085133216
\end{euleroutput}
\begin{eulercomment}
Iterasi tersebut konvergen ke penyelesaian persamaan

\end{eulercomment}
\begin{eulerformula}
\[
x = \cos(x).
\]
\end{eulerformula}
\begin{eulercomment}
Iterasi ini juga dapat dilakukan pada interval, hasilnya adalah
barisan interval yang memuat akar tersebut.
\end{eulercomment}
\begin{eulerprompt}
>hasil := iterate("cos(x)",~1,2~) //iterasi x(n+1)=cos(x(n)), dengan interval awal (1, 2)
\end{eulerprompt}
\begin{euleroutput}
  ~0.739085133211,0.7390851332133~
\end{euleroutput}
\begin{eulercomment}
Jika interval hasil tersebut sedikit diperlebar, akan terlihat bahwa interval tersebut
memuat akar persamaan x=cos(x).
\end{eulercomment}
\begin{eulerprompt}
>h=expand(hasil,100), cos(h) << h
\end{eulerprompt}
\begin{euleroutput}
  ~0.73908513309,0.73908513333~
  1
\end{euleroutput}
\begin{eulercomment}
Iterasi juga dapat digunakan pada fungsi yang didefinisikan.
\end{eulercomment}
\begin{eulerprompt}
>function f(x) := (x+2/x)/2
\end{eulerprompt}
\begin{eulercomment}
Iterasi x(n+1)=f(x(n)) akan konvergen ke akar kuadrat 2.
\end{eulercomment}
\begin{eulerprompt}
>iterate("f",2), sqrt(2)
\end{eulerprompt}
\begin{euleroutput}
  1.41421356237
  1.41421356237
\end{euleroutput}
\begin{eulercomment}
Jika pada perintah iterate diberikan tambahan parameter n, maka hasil iterasinya akan
ditampilkan mulai dari iterasi pertama sampai ke-n.
\end{eulercomment}
\begin{eulerprompt}
>iterate("f",2,5)
\end{eulerprompt}
\begin{euleroutput}
  [2,  1.5,  1.41667,  1.41422,  1.41421,  1.41421]
\end{euleroutput}
\begin{eulercomment}
Untuk iterasi ini tidak dapat dilakukan terhadap interval.
\end{eulercomment}
\begin{eulerprompt}
>niterate("f",~1,2~,5)
\end{eulerprompt}
\begin{euleroutput}
  [ ~1,2~,  ~1,2~,  ~1,2~,  ~1,2~,  ~1,2~,  ~1,2~ ]
\end{euleroutput}
\begin{eulercomment}
Perhatikan, hasil iterasinya sama dengan interval awal. Alasannya adalah perhitungan dengan
interval bersifat terlalu longgar. Untuk meingkatkan perhitungan pada ekspresi dapat
digunakan pembagian intervalnya, menggunakan fungsi ieval().
\end{eulercomment}
\begin{eulerprompt}
>function s(x) := ieval("(x+2/x)/2",x,10)
\end{eulerprompt}
\begin{eulercomment}
Selanjutnya dapat dilakukan iterasi hingga diperoleh hasil optimal, dan intervalnya tidak
semakin mengecil. Hasilnya berupa interval yang memuat akar persamaan:

\end{eulercomment}
\begin{eulerformula}
\[
x = \frac{1}{2} \left( x + \frac{2}{x} \right).
\]
\end{eulerformula}
\begin{eulercomment}
Satu-satunya solusi adalah\\
\end{eulercomment}
\begin{eulerformula}
\[
x = \sqrt2.
\]
\end{eulerformula}
\begin{eulerprompt}
>iterate("s",~1,2~)
\end{eulerprompt}
\begin{euleroutput}
  ~1.41421356236,1.41421356239~
\end{euleroutput}
\begin{eulercomment}
Fungsi "iterate()" juga dapat bekerja pada vektor. Berikut adalah contoh fungsi vektor, yang
menghasilkan rata-rata aritmetika dan rata-rata geometri.

\end{eulercomment}
\begin{eulerformula}
\[
(a_{n+1},b_{n+1}) = \left( \frac{a_n+b_n}{2}, \sqrt{a_nb_n} \right)
\]
\end{eulerformula}
\begin{eulercomment}
Iterasi ke-n disimpan pada vektor kolom x[n].
\end{eulercomment}
\begin{eulerprompt}
>function g(x) := [(x[1]+x[2])/2;sqrt(x[1]*x[2])]
\end{eulerprompt}
\begin{eulercomment}
Iterasi dengan menggunakan fungsi tersebut akan konvergen ke rata-rata aritmetika dan
geometri dari nilai-nilai awal. 
\end{eulercomment}
\begin{eulerprompt}
>iterate("g",[1;5])
\end{eulerprompt}
\begin{euleroutput}
        2.60401 
        2.60401 
\end{euleroutput}
\begin{eulercomment}
Hasil tersebut konvergen agak cepat, seperti kita cek sebagai berikut.
\end{eulercomment}
\begin{eulerprompt}
>iterate("g",[1;5],4)
\end{eulerprompt}
\begin{euleroutput}
              1             3       2.61803       2.60403       2.60401 
              5       2.23607       2.59002       2.60399       2.60401 
\end{euleroutput}
\begin{eulercomment}
Iterasi pada interval dapat dilakukan dan stabil, namun tidak menunjukkan bahwa limitnya
pada batas-batas yang dihitung.
\end{eulercomment}
\begin{eulerprompt}
>iterate("g",[~1~;~5~],4)
\end{eulerprompt}
\begin{euleroutput}
  Interval 2 x 5 matrix
  
  ~0.999999999999999778,1.00000000000000022~     ...
  ~4.99999999999999911,5.00000000000000089~     ...
\end{euleroutput}
\begin{eulercomment}
Iterasi berikut konvergen sangat lambat.

\end{eulercomment}
\begin{eulerformula}
\[
x_{n+1} = \sqrt{x_n}.
\]
\end{eulerformula}
\begin{eulerprompt}
>iterate("sqrt(x)",2,10)
\end{eulerprompt}
\begin{euleroutput}
  [2,  1.41421,  1.18921,  1.09051,  1.04427,  1.0219,  1.01089,
  1.00543,  1.00271,  1.00135,  1.00068]
\end{euleroutput}
\begin{eulercomment}
Kekonvergenan iterasi tersebut dapat dipercepatdengan percepatan Steffenson:
\end{eulercomment}
\begin{eulerprompt}
>steffenson("sqrt(x)",2,10)
\end{eulerprompt}
\begin{euleroutput}
  [1.04888,  1.00028,  1,  1]
\end{euleroutput}
\eulerheading{Iterasi menggunakan Loop yang ditulis Langsung}
\begin{eulercomment}
Berikut adalah beberapa contoh penggunaan loop untuk melakukan iterasi yang ditulis langsung
pada baris perintah.
\end{eulercomment}
\begin{eulerprompt}
>x=2; repeat x=(x+2/x)/2; until x^2~=2; end; x,
\end{eulerprompt}
\begin{euleroutput}
  1.41421356237
\end{euleroutput}
\begin{eulercomment}
Penggabungan matriks menggunakan tanda "\textbar{}" dapat digunakan untuk menyimpan semua hasil
iterasi.
\end{eulercomment}
\begin{eulerprompt}
>v=[1]; for i=2 to 8; v=v|(v[i-1]*i); end; v,
\end{eulerprompt}
\begin{euleroutput}
  [1,  2,  6,  24,  120,  720,  5040,  40320]
\end{euleroutput}
\begin{eulercomment}
hasil iterasi juga dapat disimpan pada vektor yang sudah ada.
\end{eulercomment}
\begin{eulerprompt}
>v=ones(1,100); for i=2 to cols(v); v[i]=v[i-1]*i; end; ...
>plot2d(v,logplot=1); textbox(latex(&log(n)),x=0.5):
\end{eulerprompt}
\eulerimg{27}{Image Use/TugasAplikom6Adi25030630001-237.png}
\begin{eulerprompt}
>A =[0.5,0.2;0.7,0.1]; b=[2;2]; ...
>x=[1;1]; repeat xnew=A.x-b; until all(xnew~=x); x=xnew; end; ...
>x,
\end{eulerprompt}
\begin{euleroutput}
       -7.09677 
       -7.74194 
\end{euleroutput}
\eulerheading{Iterasi di dalam Fungsi}
\begin{eulercomment}
Fungsi atau program juga dapat menggunakan iterasi dan dapat digunakan untuk melakukan iterasi. Berikut adalah beberapa contoh
iterasi di dalam fungsi.

Contoh berikut adalah suatu fungsi untuk menghitung berapa lama suatu iterasi konvergen. Nilai fungsi tersebut adalah hasil akhir
iterasi dan banyak iterasi sampai konvergen.
\end{eulercomment}
\begin{eulerprompt}
>function map hiter(f$,x0) ...
\end{eulerprompt}
\begin{eulerudf}
  x=x0;
  maxiter=0;
  repeat
    xnew=f$(x);
    maxiter=maxiter+1;
    until xnew~=x;
    x=xnew;
  end;
  return maxiter;
  endfunction
\end{eulerudf}
\begin{eulercomment}
Misalnya, berikut adalah iterasi untuk mendapatkan hampiran akar kuadrat 2, cukup cepat,
konvergen pada iterasi ke-5, jika dimulai dari hampiran awal 2.
\end{eulercomment}
\begin{eulerprompt}
>hiter("(x+2/x)/2",2)
\end{eulerprompt}
\begin{euleroutput}
  5
\end{euleroutput}
\begin{eulercomment}
Karena fungsinya didefinisikan menggunakan "map". maka nilai awalnya dapat berupa vektor.
\end{eulercomment}
\begin{eulerprompt}
>x=1.5:0.1:10; hasil=hiter("(x+2/x)/2",x); ...
>  plot2d(x,hasil):
\end{eulerprompt}
\eulerimg{27}{Image Use/TugasAplikom6Adi25030630001-238.png}
\begin{eulercomment}
Dari gambar di atas terlihat bahwa kekonvergenan iterasinya semakin lambat, untuk nilai awal
semakin besar, namun penambahnnya tidak kontinu. Kita dapat menemukan kapan maksimum
iterasinya bertambah.
\end{eulercomment}
\begin{eulerprompt}
>hasil[1:10]
\end{eulerprompt}
\begin{euleroutput}
  [4,  5,  5,  5,  5,  5,  6,  6,  6,  6]
\end{euleroutput}
\begin{eulerprompt}
>x[nonzeros(differences(hasil))]
\end{eulerprompt}
\begin{euleroutput}
  [1.5,  2,  3.4,  6.6]
\end{euleroutput}
\begin{eulercomment}
maksimum iterasi sampai konvergen meningkat pada saat nilai awalnya 1.5, 2, 3.4, dan 6.6.

Contoh berikutnya adalah metode Newton pada polinomial kompleks berderajat 3.
\end{eulercomment}
\begin{eulerprompt}
>p &= x^3-1; newton &= x-p/diff(p,x); $newton
\end{eulerprompt}
\begin{eulerformula}
\[
x-\frac{x^3-1}{3\,x^2}
\]
\end{eulerformula}
\begin{eulercomment}
Selanjutnya didefinisikan fungsi untuk melakukan iterasi (aslinya 10 kali).
\end{eulercomment}
\begin{eulerprompt}
>function iterasi(f$,x,n=10) ...
\end{eulerprompt}
\begin{eulerudf}
  loop 1 to n; x=f$(x); end;
  return x;
  endfunction
\end{eulerudf}
\begin{eulercomment}
Kita mulai dengan menentukan titik-titik grid pada bidang kompleksnya.
\end{eulercomment}
\begin{eulerprompt}
>r=1.5; x=linspace(-r,r,501); Z=x+I*x'; W=iterasi(newton,Z);
\end{eulerprompt}
\begin{eulercomment}
Berikut adalah akar-akar polinomial di atas.
\end{eulercomment}
\begin{eulerprompt}
>z=&solve(p)()
\end{eulerprompt}
\begin{euleroutput}
  [ -0.5+0.866025i,  -0.5-0.866025i,  1+0i  ]
\end{euleroutput}
\begin{eulercomment}
Untuk menggambar hasil iterasinya, dihitung jarak dari hasil iterasi ke-10 ke masing-masing
akar, kemudian digunakan untuk menghitung warna yang akan digambar, yang menunjukkan limit
untuk masing-masing nilai awal. 

Fungsi plotrgb() menggunakan jendela gambar terkini untuk menggambar warna RGB sebagai
matriks.
\end{eulercomment}
\begin{eulerprompt}
>C=rgb(max(abs(W-z[1]),1),max(abs(W-z[2]),1),max(abs(W-z[3]),1)); ...
>  plot2d(none,-r,r,-r,r); plotrgb(C):
\end{eulerprompt}
\eulerimg{27}{Image Use/TugasAplikom6Adi25030630001-240.png}
\eulerheading{Iterasi Simbolik}
\begin{eulercomment}
Seperti sudah dibahas sebelumnya, untuk menghasilkan barisan ekspresi simbolik dengan Maxima
dapat digunakan fungsi makelist().
\end{eulercomment}
\begin{eulerprompt}
>&powerdisp:true // untuk menampilkan deret pangkat mulai dari suku berpangkat terkecil
\end{eulerprompt}
\begin{euleroutput}
  
                                   true
  
\end{euleroutput}
\begin{eulerprompt}
>deret &= makelist(taylor(exp(x),x,0,k),k,1,3); $deret // barisan deret Taylor untuk e^x
\end{eulerprompt}
\begin{eulerformula}
\[
\left[ 1+x , 1+x+\frac{x^2}{2} , 1+x+\frac{x^2}{2}+\frac{x^3}{6}   \right] 
\]
\end{eulerformula}
\begin{eulercomment}
Untuk mengubah barisan deret tersebut menjadi vektor string di EMT digunakan fungsi
mxm2str(). Selanjutnya, vektor string/ekspresi hasilnya dapat digambar seperti menggambar
vektor eskpresi pada EMT.
\end{eulercomment}
\begin{eulerprompt}
>plot2d("exp(x)",0,3); // plot fungsi aslinya, e^x
>plot2d(mxm2str("deret"),>add,color=4:6): // plot ketiga deret taylor hampiran fungsi tersebut
\end{eulerprompt}
\eulerimg{27}{Image Use/TugasAplikom6Adi25030630001-242.png}
\begin{eulercomment}
Selain cara di atas dapat juga dengan cara menggunakan indeks pada vektor/list yang
dihasilkan.
\end{eulercomment}
\begin{eulerprompt}
>$deret[3]
\end{eulerprompt}
\begin{eulerformula}
\[
1+x+\frac{x^2}{2}+\frac{x^3}{6}
\]
\end{eulerformula}
\begin{eulerprompt}
>plot2d(["exp(x)",&deret[1],&deret[2],&deret[3]],0,3,color=1:4):
\end{eulerprompt}
\eulerimg{27}{Image Use/TugasAplikom6Adi25030630001-244.png}
\begin{eulerprompt}
>$sum(sin(k*x)/k,k,1,5)
\end{eulerprompt}
\begin{eulerformula}
\[
\sin x+\frac{\sin \left(2\,x\right)}{2}+\frac{\sin \left(3\,x  \right)}{3}+\frac{\sin \left(4\,x\right)}{4}+\frac{\sin \left(5\,x  \right)}{5}
\]
\end{eulerformula}
\begin{eulercomment}
Berikut adalah cara menggambar kurva

\end{eulercomment}
\begin{eulerformula}
\[
y=\sin(x) + \dfrac{\sin 3x}{3} + \dfrac{\sin 5x}{5} + \ldots.
\]
\end{eulerformula}
\begin{eulerprompt}
>plot2d(&sum(sin((2*k+1)*x)/(2*k+1),k,0,20),0,2pi):
\end{eulerprompt}
\eulerimg{27}{Image Use/TugasAplikom6Adi25030630001-246.png}
\begin{eulercomment}
Hal serupa juga dapat dilakukan dengan menggunakan matriks, misalkan kita akan menggambar
kurva

\end{eulercomment}
\begin{eulerformula}
\[
y = \sum_{k=1}^{100} \dfrac{\sin(kx)}{k},\quad 0\le x\le 2\pi.
\]
\end{eulerformula}
\begin{eulercomment}
\end{eulercomment}
\begin{eulerprompt}
>x=linspace(0,2pi,1000); k=1:100; y=sum(sin(k*x')/k)'; plot2d(x,y):
\end{eulerprompt}
\eulerimg{27}{Image Use/TugasAplikom6Adi25030630001-247.png}
\eulerheading{Tabel Fungsi}
\begin{eulercomment}
Terdapat cara menarik untuk menghasilkan barisan dengan ekspresi Maxima. Perintah
mxmtable() berguna untuk menampilkan dan menggambar barisan dan menghasilkan barisan sebagai
vektor kolom. 

Sebagai contoh berikut adalah barisan turunan ke-n x\textasciicircum{}x di x=1.
\end{eulercomment}
\begin{eulerprompt}
>mxmtable("diffat(x^x,x=1,n)","n",1,8,frac=1);
\end{eulerprompt}
\begin{euleroutput}
          1 
          2 
          3 
          8 
         10 
         54 
        -42 
        944 
\end{euleroutput}
\eulerimg{27}{Image Use/TugasAplikom6Adi25030630001-248.png}
\begin{eulerprompt}
>$'sum(k, k, 1, n) = factor(ev(sum(k, k, 1, n),simpsum=true)) // simpsum:menghitung deret secara simbolik
\end{eulerprompt}
\begin{eulerformula}
\[
\sum_{k=1}^{n}{k}=\frac{n\,\left(1+n\right)}{2}
\]
\end{eulerformula}
\begin{eulerprompt}
>$'sum(1/(3^k+k), k, 0, inf) = factor(ev(sum(1/(3^k+k), k, 0, inf),simpsum=true))
\end{eulerprompt}
\begin{eulerformula}
\[
\sum_{k=0}^{\infty }{\frac{1}{k+3^{k}}}=\sum_{k=0}^{\infty }{\frac{  1}{k+3^{k}}}
\]
\end{eulerformula}
\begin{eulercomment}
Di sini masih gagal, hasilnya tidak dihitung.
\end{eulercomment}
\begin{eulerprompt}
>$'sum(1/x^2, x, 1, inf)= ev(sum(1/x^2, x, 1, inf),simpsum=true) // ev: menghitung nilai ekspresi
\end{eulerprompt}
\begin{eulerformula}
\[
\sum_{x=1}^{\infty }{\frac{1}{x^2}}=\frac{\pi^2}{6}
\]
\end{eulerformula}
\begin{eulerprompt}
>$'sum((-1)^(k-1)/k, k, 1, inf) = factor(ev(sum((-1)^(x-1)/x, x, 1, inf),simpsum=true))
\end{eulerprompt}
\begin{eulerformula}
\[
\sum_{k=1}^{\infty }{\frac{\left(-1\right)^{-1+k}}{k}}=-\sum_{x=1  }^{\infty }{\frac{\left(-1\right)^{x}}{x}}
\]
\end{eulerformula}
\begin{eulercomment}
Di sini masih gagal, hasilnya tidak dihitung.
\end{eulercomment}
\begin{eulerprompt}
>$'sum((-1)^k/(2*k-1), k, 1, inf) = factor(ev(sum((-1)^k/(2*k-1), k, 1, inf),simpsum=true))
\end{eulerprompt}
\begin{eulerformula}
\[
\sum_{k=1}^{\infty }{\frac{\left(-1\right)^{k}}{-1+2\,k}}=\sum_{k=1  }^{\infty }{\frac{\left(-1\right)^{k}}{-1+2\,k}}
\]
\end{eulerformula}
\begin{eulerprompt}
>$ev(sum(1/n!, n, 0, inf),simpsum=true)
\end{eulerprompt}
\begin{eulerformula}
\[
\sum_{n=0}^{\infty }{\frac{1}{n!}}
\]
\end{eulerformula}
\begin{eulercomment}
Di sini masih gagal, hasilnya tidak dihitung, harusnya hasilnya e.
\end{eulercomment}
\begin{eulerprompt}
>&assume(abs(x)<1); $'sum(a*x^k, k, 0, inf)=ev(sum(a*x^k, k, 0, inf),simpsum=true), &forget(abs(x)<1);
\end{eulerprompt}
\begin{eulerformula}
\[
a\,\sum_{k=0}^{\infty }{x^{k}}=\frac{a}{1-x}
\]
\end{eulerformula}
\begin{eulercomment}
Deret geometri tak hingga, dengan asumsi rasional antara -1 dan 1.
\end{eulercomment}
\begin{eulerprompt}
>$'sum(x^k/k!,k,0,inf)=ev(sum(x^k/k!,k,0,inf),simpsum=true)
\end{eulerprompt}
\begin{eulerformula}
\[
\sum_{k=0}^{\infty }{\frac{x^{k}}{k!}}=\sum_{k=0}^{\infty }{\frac{x  ^{k}}{k!}}
\]
\end{eulerformula}
\begin{eulerprompt}
>$limit(sum(x^k/k!,k,0,n),n,inf)
\end{eulerprompt}
\begin{eulerformula}
\[
\lim_{n\rightarrow \infty }{\sum_{k=0}^{n}{\frac{x^{k}}{k!}}}
\]
\end{eulerformula}
\begin{eulerprompt}
>function d(n) &= sum(1/(k^2-k),k,2,n); $'d(n)=d(n)
\end{eulerprompt}
\begin{eulerformula}
\[
d\left(n\right)=\sum_{k=2}^{n}{\frac{1}{-k+k^2}}
\]
\end{eulerformula}
\begin{eulerprompt}
>$d(10)=ev(d(10),simpsum=true)
\end{eulerprompt}
\begin{eulerformula}
\[
\sum_{k=2}^{10}{\frac{1}{-k+k^2}}=\frac{9}{10}
\]
\end{eulerformula}
\begin{eulerprompt}
>$d(100)=ev(d(100),simpsum=true)
\end{eulerprompt}
\begin{eulerformula}
\[
\sum_{k=2}^{100}{\frac{1}{-k+k^2}}=\frac{99}{100}
\]
\end{eulerformula}
\begin{eulerprompt}
> 
\end{eulerprompt}
\eulerheading{Deret Taylor}
\begin{eulercomment}
Deret Taylor suatu fungsi f yang diferensiabel sampai tak hingga di sekitar x=a adalah:

\end{eulercomment}
\begin{eulerformula}
\[
f(x) = \sum_{k=0}^\infty \frac{(x-a)^k f^{(k)}(a)}{k!}.
\]
\end{eulerformula}
\begin{eulerprompt}
>$'e^x =taylor(exp(x),x,0,10) // deret Taylor e^x di sekitar x=0, sampai suku ke-11
\end{eulerprompt}
\begin{eulerformula}
\[
e^{x}=1+x+\frac{x^2}{2}+\frac{x^3}{6}+\frac{x^4}{24}+\frac{x^5}{120  }+\frac{x^6}{720}+\frac{x^7}{5040}+\frac{x^8}{40320}+\frac{x^9}{  362880}+\frac{x^{10}}{3628800}
\]
\end{eulerformula}
\begin{eulerprompt}
>$'log(x)=taylor(log(x),x,1,10)// deret log(x) di sekitar x=1
\end{eulerprompt}
\begin{eulerformula}
\[
\log x=-1-\frac{\left(-1+x\right)^2}{2}+\frac{\left(-1+x\right)^3}{  3}-\frac{\left(-1+x\right)^4}{4}+\frac{\left(-1+x\right)^5}{5}-  \frac{\left(-1+x\right)^6}{6}+\frac{\left(-1+x\right)^7}{7}-\frac{  \left(-1+x\right)^8}{8}+\frac{\left(-1+x\right)^9}{9}-\frac{\left(-1  +x\right)^{10}}{10}+x
\]
\end{eulerformula}
\eulerheading{Contoh Penyelesaian Masalah Kalkulus}
\begin{eulercomment}
\end{eulercomment}
\eulersubheading{1. Limit Fungsi Trigonometri Bentuk Tak Tentu}
\begin{eulercomment}
Soal 1: Hitunglah nilai limit fungsi trigonometri berikut serta
tampilkan grafiknya di sekitar titik x = 0:

\end{eulercomment}
\begin{eulerformula}
\[
\lim_{x \to 0} \frac{1 - \cos(2x)}{x^2}
\]
\end{eulerformula}
\begin{eulercomment}
\end{eulercomment}
\begin{eulerprompt}
>$showev('limit((1 - cos(2*x))/x^2, x, 0))
\end{eulerprompt}
\begin{eulerformula}
\[
\lim_{x\rightarrow 0}{\frac{1-\cos \left(2\,x\right)}{x^2}}=2
\]
\end{eulerformula}
\begin{eulerprompt}
>plot2d("(1 - cos(2*x))/x^2", -2, 2); plot2d(0, 2, >points, style="ow", >add):
\end{eulerprompt}
\eulerimg{17}{Image Use/TugasAplikom6Adi25030630001-265.png}
\begin{eulercomment}
Penjelasan hasil output: Perhitungan simbolik menggunakan Maxima
menghasilkan nilai limit sebesar 2. Secara analitis, dengan identitas
1 - cos(2x) = 2*sin\textasciicircum{}2(x), diperoleh 2 * (sin(x)/x)\textasciicircum{}2 yang menuju
2*(1)\textasciicircum{}2 = 2 saat x -\textgreater{} 0. Pada grafik, kurva mendekati ketinggian y = 2
dari kiri maupun kanan meskipun fungsi tidak terdefinisi tepat di x =
0 (ditandai lingkaran kosong).

\end{eulercomment}
\eulersubheading{2. Turunan Fungsi Hasil Kali dan Eksponensial (Aturan Leibniz /}
\begin{eulercomment}
Product Rule)

Soal 2: Tentukan turunan pertama terhadap x dari fungsi berikut dan
hitung nilainya pada x = 1:

\end{eulercomment}
\begin{eulerformula}
\[
f(x) = x^2 e^{-x}
\]
\end{eulerformula}
\begin{eulercomment}
\end{eulercomment}
\begin{eulerprompt}
>$showev('diff(x^2 * exp(-x), x))
\end{eulerprompt}
\begin{eulerformula}
\[
\frac{d}{d\,x}\,\left(x^2\,e^ {- x }\right)=2\,x\,e^ {- x }-x^2\,e
 ^ {- x }
\]
\end{eulerformula}
\begin{eulerprompt}
>$% with x=1
\end{eulerprompt}
\begin{eulerformula}
\[
{\it \%at}\left(\frac{d}{d\,x}\,\left(x^2\,e^ {- x }\right) , x=1
 \right)=e^ {- 1 }
\]
\end{eulerformula}
\begin{eulerprompt}
>$float(%)
\end{eulerprompt}
\begin{eulerformula}
\[
{\it \%at}\left(\frac{d^{1.0}}{d\,x^{1.0}}\,\left(\frac{x^2}{
 2.718281828459045^{x}}\right) , x=1.0\right)=0.3678794411714423
\]
\end{eulerformula}
\begin{eulercomment}
Penjelasan hasil output: Dengan aturan hasil kali dan aturan rantai,
turunan pertama fungsi adalah f'(x) = (2x - x\textasciicircum{}2)*e\textasciicircum{}(-x) (atau
-e\textasciicircum{}(-x)*(x\textasciicircum{}2 - 2x)). Saat disubstitusikan x = 1, nilai eksaknya adalah
1/e dan nilai desimalnya sekitar 0.3678794411714423.

\end{eulercomment}
\eulersubheading{3. Integral Tak Tentu dengan Substitusi / Bagian}
\begin{eulercomment}
Soal 3: Tentukan antiturunan (integral tak tentu) dari fungsi berikut:

\end{eulercomment}
\begin{eulerformula}
\[
\int x \cos(x)\,dx
\]
\end{eulerformula}
\begin{eulercomment}
\end{eulercomment}
\begin{eulerprompt}
>$showev('integrate(x * cos(x), x))
\end{eulerprompt}
\begin{eulerformula}
\[
\int {x\,\cos x}{\;dx}=x\,\sin x+\cos x
\]
\end{eulerformula}
\begin{eulercomment}
Penjelasan hasil output: Menggunakan integral parsial dengan misal u =
x dan dv = cos(x) dx, diperoleh hasil antiturunan x*sin(x) + cos(x)
(ditambah konstanta integrasi C pada penulisan analitis). Jika
diturunkan kembali terhadap x, hasilnya kembali ke x*cos(x).

\end{eulercomment}
\eulersubheading{4. Integral Tentu dan Luas Daerah di Antara Dua Kurva}
\begin{eulercomment}
Soal 4: Hitunglah luas daerah yang dibatasi oleh kurva y = x\textasciicircum{}2 dan
garis y = x pada selang perpotongannya, lalu gambarkan daerah yang
diarsir:

\end{eulercomment}
\begin{eulerformula}
\[
L = \int_0^1 \left(x - x^2\right) dx
\]
\end{eulerformula}
\begin{eulercomment}
\end{eulercomment}
\begin{eulerprompt}
>plot2d(["x", "x^2"], -0.2, 1.2, color=[red, blue]); ...
>x=linspace(0,1,200); xi=flipx(x); ...
>plot2d(x|xi, x|-xi^2, >filled, style="|", fillcolor=lightgray, >add):
\end{eulerprompt}
\eulerimg{17}{Image Use/TugasAplikom6Adi25030630001-273.png}
\begin{eulerprompt}
>$showev('integrate(x - x^2, x, 0, 1))
\end{eulerprompt}
\begin{eulerformula}
\[
\int_{0}^{1}{x-x^2\;dx}=\frac{1}{6}
\]
\end{eulerformula}
\begin{eulercomment}
Penjelasan hasil output: Kurva y = x dan y = x\textasciicircum{}2 berpotongan di x = 0
dan x = 1, dengan y = x berada di atas y = x\textasciicircum{}2 pada interval [0, 1].
Integrasi x - x\textasciicircum{}2 dari 0 sampai 1 menghasilkan luas sebesar 1/6
(sekitar 0.1667), yang tampak sebagai daerah terarsir di antara kedua
kurva pada grafik.

\end{eulercomment}
\eulersubheading{5. Deret Taylor / Hampiran Fungsi di Sekitar Titik Pusat}
\begin{eulercomment}
Soal 5: Tentukan polinomial Taylor orde 5 dari fungsi eksponensial di
sekitar x = 0:

\end{eulercomment}
\begin{eulerformula}
\[
e^x \approx \sum_{k=0}^{5} \frac{x^k}{k!}
\]
\end{eulerformula}
\begin{eulercomment}
\end{eulercomment}
\begin{eulerprompt}
>&powerdisp:true;
>$'exp(x) = taylor(exp(x), x, 0, 5)
\end{eulerprompt}
\begin{eulerformula}
\[
e^{x}=1+x+\frac{x^2}{2}+\frac{x^3}{6}+\frac{x^4}{24}+\frac{x^5}{120
 }
\]
\end{eulerformula}
\begin{eulerprompt}
>plot2d(["exp(x)", &taylor(exp(x), x, 0, 5)], -1, 2, color=[blue, red]):
\end{eulerprompt}
\eulerimg{17}{Image Use/TugasAplikom6Adi25030630001-277.png}
\begin{eulercomment}
Penjelasan hasil output:
Ekspansi deret Taylor orde 5 di sekitar x = 0 menghasilkan 1 + x + x\textasciicircum{}2/2 +
x\textasciicircum{}3/6 + x\textasciicircum{}4/24 + x\textasciicircum{}5/120. Pada grafik, kurva hampiran (merah) berimpit sangat
dekat dengan kurva asli e\textasciicircum{}x (biru) di sekitar x = 0.
\end{eulercomment}
\end{eulernotebook}
\end{document}


% ====================================================================



% ====================================================================
\end{document}